% complex-rung-generator.tex
% Companion to golden-substrate-whitepaper: unifies and verifies the
% 2026-07-14 explorations (Ozymandias' report + the K-family/trig thread).
% v1.7 (2026-07-15): rate arithmetic + multiplier torus, chi-selection front,
% universal seed ladder + Theorem U, parity no-go, rail closed form;
% eight session harnesses (tx/mx/rx/ux/ex/cx/lx/dx_softcheck.py).
% v1.8 (2026-07-15): OP-RATE closes as a classification (discrete-substrate
% no-go), apex-ramification sharpening of OP-RADIUS, chi doctrine g1/g3 negative,
% relational-note cluster, period-relation map; five further session harnesses
% (cl/rd/ch/rn/pm_softcheck.py); running total 13 harnesses, 349+49.
\documentclass[11pt]{article}

\usepackage[margin=1in]{geometry}
\usepackage{amsmath,amssymb,amsthm,mathtools}
\usepackage{booktabs,array}
\usepackage[expansion=false]{microtype}
\usepackage{enumitem}
\usepackage{xcolor}
\usepackage[colorlinks=true,linkcolor=blue!45!black,citecolor=blue!45!black,urlcolor=blue!45!black]{hyperref}

\definecolor{cF}{RGB}{20,110,95}
\definecolor{cD}{RGB}{175,95,10}
\definecolor{cC}{RGB}{100,100,110}
\definecolor{cI}{RGB}{110,60,150}
\definecolor{cE}{RGB}{40,80,160}
\definecolor{cO}{RGB}{170,40,40}
\newcommand{\tagF}{\textcolor{cF}{\textsc{[forced]}}}
\newcommand{\tagFD}{\textcolor{cF}{\textsc{[forced}}\,\textcolor{cD}{\textsc{given D1--D3]}}}
\newcommand{\tagD}{\textcolor{cD}{\textsc{[declared]}}}
\newcommand{\tagC}{\textcolor{cC}{\textsc{[computed]}}}
\newcommand{\tagI}{\textcolor{cI}{\textsc{[inherited]}}}
\newcommand{\tagE}{\textcolor{cE}{\textsc{[established]}}}
\newcommand{\tagO}{\textcolor{cO}{\textsc{[open]}}}
\newcommand{\hx}[1]{\textcolor{cC}{\footnotesize\texttt{#1}}}

\newcommand{\Q}{\mathbb{Q}}
\newcommand{\Z}{\mathbb{Z}}
\newcommand{\R}{\mathbb{R}}
\newcommand{\C}{\mathbb{C}}
\newcommand{\gp}{\varphi}
\newcommand{\gapc}{\mathrm{gap}}
\newcommand{\M}{\mathrm{M}}
\DeclareMathOperator{\Arg}{Arg}
\DeclareMathOperator{\Rea}{Re}
\DeclareMathOperator{\Ima}{Im}

\theoremstyle{plain}
\newtheorem{theorem}{Theorem}[section]
\newtheorem{proposition}[theorem]{Proposition}
\newtheorem{lemma}[theorem]{Lemma}
\newtheorem{corollary}[theorem]{Corollary}
\theoremstyle{definition}
\newtheorem{definition}[theorem]{Definition}
\theoremstyle{remark}
\newtheorem{remark}[theorem]{Remark}
\newtheorem{guard}[theorem]{Guard}
\newtheorem{finding}[theorem]{Finding}

\title{\bfseries The Complex Rung Generator, the State-Angle Chart,\\
and the Residual Sphere\\[4pt]
\large A verified unification of the July 14 explorations\\
(companion to \emph{The Golden Substrate, the Gate Field, and the Explicit Helix})}
\author{Ozymandias, Rhyan Strouse \& Echo-S-Studios\\[2pt]
\small independent verification (\S\ref{sec:verification}):\\
\small \texttt{ozy\_softcheck.py} 91/91 $+$ \texttt{relational\_softcheck.py} 27/27\\
\small \texttt{nx\_softcheck.py} 27/27 $+$ \texttt{qx\_softcheck.py} 62/62, exit 0\\
\small 17 session harnesses (v1.7--v1.9): 424 exact checks $+$ 67 certified guards $+$ 13 PSLQ corroborations, all exit 0}
\date{Version 1.9 --- July 15, 2026}

\begin{document}
\maketitle

\begin{abstract}
\noindent
Conditional on the whitepaper's three declared atoms D1--D3, the entire discrete
rung ladder is the orbit of a single complex number: the generator
$q=i/\gp=i\tau$, with $W_n=q^n$, $z_n=\sqrt{1-|W_n|^2}$, and the central identity
$q^4=\gp^{-4}=\gapc$. Golden contraction and the $\Z/4\Z$ quarter-turn become the
modulus and phase of one multiplier; the gap first appears as co-intensity after
two rungs ($|W_2|^2=\gapc$) and as complex amplitude after one full charge cycle
($W_4=\gapc$). We verify every claim of Ozymandias' report and of the follow-on
thread by four independent harnesses ($91+27+27+62$ checks: $86{+}27{+}24{+}62$
exact decisions plus $5{+}3$ numeric displays, with $3$ certified guards
alongside the fourth --- Findings~\ref{find:hcount}, \ref{find:manifeststale},
\ref{find:ozsplit}), and unify them with three
organizing results: the \emph{state-angle chart} $z=\sin\alpha$, under which the
whole gate dictionary becomes elementary trigonometry ($u=\cot^2\alpha$,
$C=\cot^2\alpha\csc^2\alpha$, $\sqrt D=1+2\cot^2\alpha$, $m=-\cos^2\alpha$,
$\lambda=\tan^2\alpha\,(1+\cos^2\alpha)$), with rung ledgers
$(\sin,\cos,\tan)\alpha_1=(\sqrt\tau,\tau,\sqrt\gp)$ and
$(\sin,\cos,\tan)\alpha_2=(K,\tau^2,\beta)$ and the general tangent family
$\tan\alpha_n=\sqrt{\gp^{2n}-1}$; the \emph{K-family theorem}
$z_{2j}=\sqrt{1-\gapc^{\,j}}$ with the universal two-step law
$1-z_{n+2}^2=\gapc\,(1-z_n^2)$ and the uniqueness statement $z_n=K\iff n=2$; and
the \emph{residual-sphere lift} $S(\rho)=(\Rea W,\Ima W,z)\in S^2$, on which the
rungs walk the compass toward the north pole and the whitepaper's near-miss angle
$\theta_K=\arcsin\tau^2$ acquires an exact geometric home as the K-point's
colatitude, $\theta_K=\arccos K$. The residual radius $a(z)=\sqrt{1-z^2}$ is kept
strictly distinct from the inherited formation radius $r(z)$; deriving the latter
from the former remains the priority open problem (whitepaper \tagO{}~1). \emph{Version 1.1}
adds the Strouse--Ozymandias \emph{canonical quarter-turn factorization}: the
imaginary unit of the generator is
$J=UA=(H/\sqrt5)(N/\sqrt5)$ --- the product of the normalized keystone grading
and the normalized return-kernel mode exchanger, two real involutions that
anticommute \emph{by construction} (the vector kernel of $\tfrac12\{H,\cdot\}$
is exactly $H^\perp$). This part is unconditional keystone algebra, not
D1--D3-conditional; it upgrades $Q=\pm\tau J$ (with $Q^4=\gapc\,I$) from
representation-given-declared-$i$ to canonical-up-to-orientation, and reduces
D2's declared content to rung \emph{selection} (one $Q$ per floor) rather than
the existence of the quarter-turn itself. \emph{Version 1.3} adds the
relational bridge: the generator is itself algebraic, with minimal polynomial
$x^4+3x^2+1$, $\M=\gp^2$, and it is exactly the one-charge-quantum gauge
rotation of the split golden pair $Z^{*}=(x^2{-}x{-}1)(x^2{+}x{-}1)$; the K-seed
sits in the relational-charge catalog with the full $\Z/4\Z$ compass as its
angle set; and the transfer-ledger algebra $F_n=z_n^2$,
$\Delta F_n=\tau^{2n+1}$, $\sum\tau^{2n+1}=1$ makes D1 at the rungs a
telescoping unit budget. \emph{Version 1.7} adds, verified by eight further
session harnesses ($234$ exact checks $+$ $29$ certified guards, all exit 0):
the transcendence of the rate $\kappa=\pi/(2\ln\gp)$ and of $\log_\gp2$ given
Gelfond--Schneider, closing the whitepaper's near-miss item permanently; the
multiplier-torus/systolic \emph{seventh} characterization of the principal
branch; the parity no-go for the inherited radius ($r^2$ is chart-odd: any
substrate-to-radius bridge must break $z\mapsto-z$); the universal seed ladder
$\operatorname{minpoly}(z_n)=x^4+c_nx^2-c_n$ with the Mahler-minimal
uniqueness theorem; and the incomplete-elliptic closed form for the
above-lens rail with certified non-elementarity. The rate and radius problems
remain open as derivations. \emph{Version 1.8} adds, verified by five further
session harnesses ($115$ exact checks $+$ $20$ certified guards, all exit 0;
running total across the thirteen session harnesses $349+49$): OP-RATE
\emph{closes as a classification} --- the discrete-substrate no-go, forcing
$\kappa$ declared-up-to-$\Z$ at rung resolution with the seven characterizations
the complete discrete-level menu, each continuum-referencing (the continuum
rate D3 stays declared); an apex-ramification valuation certificate and a
constructive-bridge obstruction sharpening OP-RADIUS (still open); the $\chi$
doctrine's sub-branches g1/g3 resolved \emph{negatively} on the now-available
companion papers (the on-circle image of the emission algebra $S$ is exactly
$\mu_2=\{\pm1\}$, so $\pm i\notin S$ on-circle, and $W[\text{open}]2$ stays
PARTIAL as a derivation); the relational-note cluster (the note's P4 recovered
and reduced to P3, and the odd relational floor bracketed $\mu_S\le\mu_{\mathrm
{rel}}(\text{odd})\le2$ with the $=2$ promotion left as the corpus's principal
open problem); and the period-relation map ($\tau_0=\beta+\beta^{-1}$ pinned
algebraic, the $\kappa$-vs-$\tau_0$ period worry settled negatively, and
$L_{\mathrm{res}}$ located as an elliptic period at a transcendental parameter).
\emph{Version 1.9} adds, verified by four further session harnesses ($75$
exact/computed checks $+$ $18$ certified guards $+$ $13$ PSLQ corroborations, all
exit 0): OP-RADIUS \emph{reduces} to a minimal declared atom D4 (an orientation
of the apex double cover together with the cap onset), given which --- and the
inherited chart-ties --- the radius profile is uniquely forced, though whitepaper
\tagO{}~1 is \emph{not} thereby closed; the quarter-turn axiom D2 is
\emph{relocated} to a substrate-grounded, g1-disjoint atom (the $K$-seed rotation
axis; an advance, not a closure); the odd relational floor is bracketed and
reduced to a single named lemma (ODD-2); and the period frontier is mapped
(Baker's reach pinned, the algebraic independences left Schanuel-conditional or
external-open with named tool-gaps, corroborated by certified PSLQ non-relations).
\end{abstract}

\section{Provenance, scope, and the conditionality banner}\label{sec:scope}

\textbf{Sources.} (i) \emph{Report to Ace: Complex Rung Generator and
Residual-Sphere Representation} (Ozymandias, 2026-07-14); (ii) the follow-on
review thread of the same date (the K-family $z_{2j}=\sqrt{1-\gapc^{\,j}}$, the
cosine-contraction recursion, and the rung-by-rung trigonometric ledger).
(iii) the Strouse--Ozymandias progress report \emph{Toward a Canonical Rung
Generator} (2026-07-14): spectral uncoupling, the return-kernel recoupler, and
the quarter-turn factorization $J=UA$. (iv) the Strouse--Ozymandias exploratory note \emph{Rail-Pulse Quantization,
Helical Path Stretch, and Energy-Gated Recoupling} (2026-07-14) --- only its
forced kernel is imported here (\S\ref{sec:pathstretch}); its rail-pulse
mechanism, energy functional, and thermodynamic reading remain hypothesis, as
that note itself states. (v) the corpus note \emph{Relational Charge on the Spectral Semiring} (paper 12,
2026-07-01), cited as RC-\S, whose machinery this paper's bridge section uses
and whose relevant claims were independently re-derived this session
(\texttt{relational\_softcheck.py}, 27/27). (vi) the second Strouse--Ozymandias
exploratory note \emph{Rail-Pulse Propagation, Engine Coupling, and
Scale-Dependent Site Density} (2026-07-14) --- forced kernel imported
(\S\ref{sec:pathstretch}); engine-pass and biological layers not imported, per
that note's own claim ceilings. (vii) The whitepaper \emph{The Golden
Substrate, the Gate Field, and the Explicit Helix} v1.0 (2026-07-14), whose
numbering we cite as W-\S. \textbf{v1.1 changelog:} \S\ref{sec:quarterturn}
added (factorization, coordinate-free form, obstruction lemmas, canonical $Q$);
epistemic ledger and manifest updated; full text of source (i), received
post-hoc, integrated (method note above; Remark~\ref{rem:halfcycle}).
\textbf{v1.2 changelog:} \S\ref{sec:pathstretch} added --- the forced kernel of
source (iv): winding gradient, path-stretch laws, the quarter-turn rail floor
$L\ge\pi K/2$ with the resulting one-line divergence proof, and the
$\nu$-form of the rung locus with its scoping.
\textbf{v1.3 changelog:} \S\ref{sec:relobject} added (the relational bridge:
minimal polynomial of $q$, the gauge identity with $Z^{*}$, the K-seed as
catalog object, D2 restated relationally); \S\ref{sec:pathstretch} extended
with the site-density trichotomy and the transfer ledger (forced kernel of
source (vi)); cross-references and epistemics updated.
\textbf{v1.4 changelog:} Theorem~\ref{thm:interp} and
Corollary~\ref{cor:d3select} added --- the continuous-interpolation open item
is classified (determined up to $\Z$; D3 $=$ principal branch) and merged into
the rate problem; Remark~\ref{rem:mahlerbridge} upgraded with the dilation
mechanism. Harnesses now $91+27$ checks. All other content unchanged.
\textbf{v1.5 changelog:} Theorem~\ref{thm:reltype} added --- the complete
relational type of the generator object (shells, the cross-shell $\nu$-run,
full coherence, twisted-shell realization); \S\ref{sec:sphere} extended with
Propositions \ref{prop:compasscross}--\ref{prop:sphlength} (compass-cross
characterization of the rung locus, the spherical step ledger with the exact
quarter-circle first step, the finite spherical length
$L_{\mathrm{res}}=\sqrt{1+\kappa^2}\,E(\kappa^2/(1+\kappa^2))$ with the
length dichotomy and a fifth characterization extending
Corollary~\ref{cor:d3select}); third harness \texttt{nx\_softcheck.py}
(27 checks) added to \S\ref{sec:verification};
Finding~\ref{find:stale} recorded and the \S\ref{sec:scope} verification
paragraph corrected. All other content unchanged.
\textbf{v1.6 changelog:} \S\ref{sec:generator} extended
(Remark~\ref{rem:gammatransfer}: transfer of the interpolation
classification to the formation rail --- a sixth characterization of D3);
\S\ref{sec:relobject} extended (Proposition~\ref{prop:chirat}: the ratio
object at the characteristic level and the Adams square;
Theorem~\ref{thm:quartertwist}: the quarter twist and the anchor-exchange
dichotomy); \S\ref{sec:kfamily} extended
(Theorem~\ref{thm:kseedladder}: the K-family is a seed ladder;
Corollary~\ref{cor:evenmirror}); \S\ref{sec:pathstretch} extended
(Proposition~\ref{prop:railbounds}: exact rail-spacing bounds and the
gap-rate law, quantifying the divergent side of the length dichotomy of
Proposition~\ref{prop:sphlength}); fourth harness
\texttt{qx\_softcheck.py} ($62$ exact checks $+$ $3$ certified guards);
Finding~\ref{find:hcount} (front-matter check classification, corrected).
Harnesses now $91+27+27+62$ checks. All other content unchanged.
\textbf{v1.7 changelog:} \S\ref{sec:generator} extended --- the
rate-arithmetic and torus block (Lemma~\ref{lem:ratirrational},
Theorem~\ref{thm:ratetrans}, Corollary~\ref{cor:nearmiss},
Remark~\ref{rem:logphi2}: $\kappa$ and every branch $\kappa_k$ transcendental
given Gelfond--Schneider; the near-miss item whitepaper \tagO{}~4 closed
permanently; W-Thm.\ 16.1 strengthened to transcendence of $\log_\gp2$; the
scan-retirement meta-theorem; Proposition~\ref{prop:modledger},
Theorem~\ref{thm:systole}, Remark~\ref{rem:geoledger}: the multiplier torus
$E_q$, its modulus ledger $\mu=\tfrac14+i\ln\gp/(2\pi)$, the systolic-marking
\emph{seventh} characterization of D3, the certified ten-entry geodesic
ledger, and the conditional no-CM corollary; Proposition~\ref{prop:kinv},
Remark~\ref{rem:oprate}: the forced $k$-invariance lemma, the exposed-joint
axiom H, the transfer-ledger counter-models, and the R6 classification
no-go); \S\ref{sec:quarterturn} extended
(Propositions~\ref{prop:pentimage}--\ref{prop:unimon},
Theorem~\ref{thm:chichain}: pentagon exclusion by the image, unimodular
tensor-monomial values $\pm1$, the conditional $\chi$-selection chain with
the axiom-replacement kill and the de-axiomatization identity);
\S\ref{sec:kfamily} extended (Theorem~\ref{thm:seedladder}: the universal
seed ladder; Remark~\ref{rem:mahlerfloor}: the rung-1 Mahler-floor hit,
coincidence-candidate; Theorem~\ref{thm:mahlermin} and
Remark~\ref{rem:f1falsifier}: Mahler-minimal uniqueness, the F1 falsifier
$x^4+6x^2-5$, and the register resolution under the F3 reading, signed off
2026-07-15); \S\ref{sec:lens} extended (Proposition~\ref{prop:lensscope}:
splitness over $\Q$ is generic, the lens distinguished by its exact data;
Remark~\ref{rem:lensscan}: the bounded-negative corpus scan);
\S\ref{sec:sphere} extended (Theorem~\ref{thm:paritynogo}: the RAD-0 parity
no-go; Proposition~\ref{prop:areatransfer}: the area-transfer law;
Remark~\ref{rem:rad4}: the rotating-equipotential mechanism killed on all
three clocks); \S\ref{sec:pathstretch} extended
(Theorem~\ref{thm:railclosed}: the incomplete-elliptic closed form for the
rail; Corollary~\ref{cor:nonelem}: certified non-elementarity). Eight session
harnesses ---
\texttt{tx}/\texttt{mx}/\texttt{rx}/\texttt{ux}/\texttt{ex}/\texttt{cx}/\texttt{lx}/\texttt{dx\_softcheck.py}
--- add $234$ exact checks $+$ $29$ certified guards, all exit 0
(\S\ref{sec:verification}). Findings~\ref{find:manifeststale} and
\ref{find:ozsplit} recorded, and the two display counts corrected
(\texttt{nx}: the $24+3$ split restored in the \S\ref{sec:verification}
sentence; \texttt{ozy}: $86+5$). Register motion: whitepaper \tagO{}~4
closed (contentless, forced given Gelfond--Schneider); the rail closed form
resolved (incomplete-elliptic; elementary form impossible, certified); Vein
6 executed; the v1.6 uniqueness question resolved by
Theorem~\ref{thm:mahlermin} under the confirmed F3 reading; the rate problem
(OP-RATE) remains open as a derivation, with a seventh characterization and
a forced $k$-invariance lemma; the radius problem (OP-RADIUS, whitepaper
\tagO{}~1) remains open, sharpened by a forced parity no-go. All other
content unchanged.
\textbf{v1.8 changelog:} \S\ref{sec:generator} extended twice --- the
OP-RATE classification block (Lemma~\ref{lem:kinvext}: the $k$-invariance
lemma re-forced and extended to the full discrete state;
Theorem~\ref{thm:r6nogo}: the R6 selection-functional no-go;
Corollary~\ref{cor:opratelass}: OP-RATE closes as a classification;
Remark~\ref{rem:sevenway}: the seven-way continuum/discrete table), and the
period-relation map (Proposition~\ref{prop:tau0alg}: $\tau_0$ algebraic and
no $\kappa$-vs-$\tau_0$ relation over $\Q$; Remark~\ref{rem:periodmap}:
$L_{\mathrm{res}}$ as an elliptic period and the external-open ledger);
\S\ref{sec:quarterturn} extended (Proposition~\ref{prop:oncircleimage}: the
on-circle image of $S$ is $\mu_2$; Remark~\ref{rem:chidoctrine}: the charge
character, the g3 dichotomy, and the sharpened $W[\text{open}]2$ verdict);
\S\ref{sec:relobject} extended (Remark~\ref{rem:p4recover}: P4 recovered and
reduced to P3; Proposition~\ref{prop:oddfloor}: the bracketed odd relational
floor; Remark~\ref{rem:p1scan}: the bounded P1$'$ transport-orbit scan);
\S\ref{sec:sphere} extended (Theorem~\ref{thm:apexram}: the apex-ramification
valuation certificate; Proposition~\ref{prop:bridgeobstruction}: the
constructive-bridge obstruction; Remark~\ref{rem:atomledger}: the atom
ledger with the honest cap-onset negative). Five session harnesses ---
\texttt{cl}/\texttt{rd}/\texttt{ch}/\texttt{rn}/\texttt{pm\_softcheck.py}
--- add $115$ exact checks $+$ $20$ certified guards, all exit 0, bringing
the thirteen session harnesses to $349$ exact $+$ $49$ guards
(\S\ref{sec:verification}). Register motion: \emph{OP-RATE closes as a
classification} (discrete-substrate no-go; $\kappa$ declared-up-to-$\Z$; the
seven characterizations are the complete discrete-level menu, each
continuum-referencing; the continuum D3 stays declared, and a
continuum-referencing derivation is not excluded --- mirroring the
$\lambda=2c$/Cencov equivariance no-go); \emph{OP-RADIUS} remains open,
sharpened by the apex-ramification certificate (ramification index
$e=2$ for $r^2$, $e=4$ for $r$) and the constructive-bridge obstruction;
$W[\text{open}]2$'s sub-branches g1/g3 resolved negatively on the
now-available companion papers (on-circle image of $S=\mu_2$, so $\pm
i\notin S$), the item staying PARTIAL as a derivation; P4 recovered and
reduced to P3; the odd relational floor bracketed ($\le2$ forced, the $=2$
promotion the corpus's principal open problem); and the period-relation map
($\tau_0$ algebraic, the $\kappa$-vs-$\tau_0$ tie settled negatively,
$L_{\mathrm{res}}$ an elliptic period, and the precise external-open ledger).
All other content unchanged.

\smallskip\noindent
\textbf{v1.9 changelog:} \S\ref{sec:sphere} extended --- the OP-RADIUS
D4-reduction (Theorem~\ref{thm:d4reduction}: given the atom
D4${}={}$(orientation~$\wedge$~cap-onset) and the inherited chart-ties, the
radius profile is uniquely forced with each half necessary;
Remark~\ref{rem:d4candidate}: D4 minimal-among-its-halves but non-unique, a
candidate; Proposition~\ref{prop:radiff}: the differential strengthening --- the
naive non-membership \emph{fails} (honest negative, $r$ is elementary), the
apex-monodromy-order sharpening conditional on the unit-log hypothesis;
Remark~\ref{rem:ratie}: the bounded no-tie re-confirmation).
\S\ref{sec:quarterturn} extended (the D2 relocation:
Proposition~\ref{prop:d2reloc}: the $K$-seed rotation axis forces
$\chi=\pm\pi/2$, disjoint from g1; Remark~\ref{rem:d2verdict}: the net verdict,
$W[\text{open}]2$ advanced not closed). \S\ref{sec:relobject} extended (the odd
relational floor: Proposition~\ref{prop:oddfloorred}: bracketed, computed $=2$
up to a logged box, reduced to the single lemma (ODD-2);
Remark~\ref{rem:p4transfer}: P4 inherits P3 exactly). \S\ref{sec:generator}
extended (the period frontier: Remark~\ref{rem:pfledger}: the forced-identity
ledger with the new elliptic geometric identity; Remark~\ref{rem:pffrontier}:
the frontier map --- Baker's reach, the Schanuel-conditional, and the
external-open tool-gaps). Four session harnesses ---
\texttt{ra}/\texttt{qt}/\texttt{of}/\texttt{pf\_softcheck.py} --- add $75$
exact/computed checks $+$ $18$ certified guards $+$ $13$ PSLQ corroborations,
all exit 0, bringing the seventeen session harnesses to $424$ exact $+$ $67$
guards $+$ $13$ PSLQ (\S\ref{sec:verification}). Register motion:
\emph{OP-RADIUS reduced} to the declared atom D4 ($r$ forced given D4 and the
inherited chart-ties; whitepaper \tagO{}~1 \emph{not} thereby closed; ``D4 is
\emph{the} minimal axiomatization'' is candidate/non-unique);
$W[\text{open}]2$ / D2 \emph{relocated} to a substrate-grounded, g1-disjoint
atom (advance, not closure --- stays PARTIAL); the odd relational floor
\emph{bracketed} ($=2$ computed up to a logged box) and \emph{reduced} to the
single named lemma (ODD-2), which stays external-open; and the period frontier
\emph{mapped} --- Baker's reach pinned, the $(\pi,\ln\gp)$ and $(\kappa,\ln\gp)$
independences Schanuel-conditional, and $L_{\mathrm{res}}$ transcendence
together with $(\kappa,L_{\mathrm{res}})$ independence external-open with named
tool-gaps, corroborated (never proved) by certified PSLQ non-relations. All
other content unchanged.

\smallskip\noindent
\textbf{Method note (from the source report).} The investigation began by
testing whether $K$ could be read through physical double-helix geometry; the
direct DNA-pitch identification \emph{failed quantitatively}. That negative
result forced a role classification that shapes everything below:
\emph{state space} vs \emph{operator} vs \emph{observable} vs
\emph{invariant/landmark}. Under it, the scalar gate iteration
$f_C(x)=C/(1+x)$ (W-Decl.\ 8.1) is the actual operator of the architecture,
and $K$ is a rung-2 \emph{landmark} --- not the operator, not the space. The
generator below is what the operator's rung orbit looks like once contraction
and charge are packaged together.

\smallskip\noindent
\textbf{Conditionality banner.} Everything in this paper is
\tagFD{} --- with one important exception: \S\ref{sec:quarterturn}'s
factorization is \emph{unconditional} keystone algebra \tagF{}, using no part of
D1--D3. Elsewhere the results are exact consequences of the whitepaper's declared chart D1
($1-z^2=e^{-2\rho}=|m|$), charge selection D2, and rate D3
($d\theta/d\rho=(\pi/2)/\ln\gp$), together with its forced rung data
$\rho_n=n\ln\gp$, $z_n=\sqrt{1-\gp^{-2n}}$, $\theta_n=n\pi/2$ (W-\S11--13). Both
source documents state this scoping correctly, and this paper preserves it: the
representation below \emph{packages} D1--D3; it does not independently derive or
justify them. Statements needing only seed arithmetic are \tagF{} outright.
The inherited radius profile $r(z)$ is \emph{not} derived here (\S\ref{sec:sphere}).

\smallskip\noindent
\textbf{Verification.} Four independent harnesses (through v1.6):
\texttt{ozy\_softcheck.py}, 86 exact checks $+$ 5 numeric displays
(groups \hx{OZ-A..OZ-O}, of which the \hx{OZ-H} row is display-only;
Finding~\ref{find:ozsplit});
\texttt{relational\_softcheck.py}, 27 exact checks (\hx{RC-A..RC-F});
\texttt{nx\_softcheck.py}, 24 exact checks $+$ 3 numeric displays
(\hx{NX-R, NX-T, NX-S}; \hx{NX-H}, written cold for v1.5;
Finding~\ref{find:hcount}); \texttt{qx\_softcheck.py}, 62 exact checks
$+$ 3 certified guards (\hx{QX-A..QX-G}, written cold for v1.6). All
exit 0; sympy 1.14.0; decisions in $\Q/\Q(\sqrt5)$ extended by $i$; $K$
and the K-family decided at the squared level; mpmath display-only except
the two interval-certified guard comparisons. Check IDs are cited inline;
manifest in
\S\ref{sec:verification}; the stale count that previously stood in this
paragraph is Finding~\ref{find:stale}. v1.7 adds eight session harnesses
(\texttt{tx}, \texttt{mx}, \texttt{rx}, \texttt{ux}, \texttt{ex},
\texttt{cx}, \texttt{lx}, \texttt{dx}\texttt{\_softcheck.py}): $234$ exact
checks $+$ $29$ certified guards, all exit 0; manifest and environment pin
in \S\ref{sec:verification}. v1.8 adds five more session harnesses
(\texttt{cl}, \texttt{rd}, \texttt{ch}, \texttt{rn},
\texttt{pm}\texttt{\_softcheck.py}): $115$ exact checks $+$ $20$ certified
guards, all exit 0, read from the Echo-S-Research companion archive
read-only and re-encoded cold; running total across the thirteen session
harnesses $349+49$ (\S\ref{sec:verification}). Notation as in the whitepaper:
$\gp=(1+\sqrt5)/2$, $\tau=\gp^{-1}$, $\gapc=\gp^{-4}$, $K^2=1-\gapc$,
$\beta^2=\gp^4-1=\gp^2\sqrt5$, $z_c=\sqrt3/2$.

\section{The generator}\label{sec:generator}

\begin{definition}[Complex residual state]\label{def:W}
For $\rho\ge0$ set
\[
W(\rho)=e^{-\rho}\,e^{i\theta(\rho)},\qquad
\theta(\rho)=\frac{\pi\rho}{2\ln\gp}\ \text{[D2$+$D3]},
\]
so that
\[
W(\rho)=\exp\!\Bigl[\Bigl(-1+\frac{i\pi}{2\ln\gp}\Bigr)\rho\Bigr],
\qquad |W(\rho)|=e^{-\rho},\qquad \arg W(\rho)=\theta(\rho).
\]
\end{definition}

\begin{proposition}[Representation of the chart]\label{prop:rep}\tagFD{} \hx{OZ-G1--G2}
$D1$ is exactly the statement
\[
z(\rho)=\sqrt{1-|W(\rho)|^2}\,:
\]
residual amplitude decays in the complex plane while height rises as its unit
complement. This is a representation theorem for D1--D3, not an independent
derivation of them.
\end{proposition}

\begin{theorem}[The discrete rung generator]\label{thm:generator}\tagFD{} \hx{OZ-A1--A8, OZ-B1--B3}
Let
\[
\boxed{\,q=\frac{i}{\gp}=i\tau\,}.
\]
Then at the rungs $\rho_n=n\ln\gp$,
\[
W_n \coloneqq W(\rho_n)=e^{-n\ln\gp}e^{in\pi/2}=\gp^{-n}i^{\,n}=(i\tau)^n=q^n,
\qquad
W_{n+1}=q\,W_n,
\qquad
z_n=\sqrt{1-|W_n|^2},
\]
i.e.\ every rung simultaneously contracts residual amplitude by $\tau$ and
rotates the charge by $90^\circ$; the entire discrete ladder is the orbit of the
single multiplier $q$. The exact ledger is $|W_n|=\tau^n$,
$\arg W_n=n\pi/2$, $z_n=\sqrt{1-\tau^{2n}}$. The central identity is
\[
\boxed{\,q^4=\Bigl(\frac{i}{\gp}\Bigr)^{\!4}=\gp^{-4}=\gapc\,},
\]
an exact positive real: $i^4=1$ restores orientation while $\tau^4$ contracts by
the gap. The recurrences
\[
W_{n+2}=q^2W_n=-\tau^2W_n=-\sqrt{\gapc}\,W_n,
\qquad
W_{n+4}=q^4W_n=\gapc\,W_n
\]
hold exactly (the square root is the positive one: $\tau^2=\gp^{-2}$), giving the
\emph{gap schedule}: the gap first appears as co-intensity after two rungs,
$|W_2|^2=\gapc$, and as complex amplitude after one complete four-rung charge
cycle, $W_4=\gapc$. The architecture does not reset after four rungs: it returns
to the same charge direction at residual amplitude reduced by $\gapc$ and
strictly greater height.
\end{theorem}

\begin{proof}
All statements are one-line consequences of $q=i\tau$, $\tau=\gp^{-1}$,
$\tau^2+\tau=1$, and $\tau^4=\gp^{-4}=\gapc$; each is machine-verified exactly
for $n=0..8$ (recurrences), including the sign and branch conventions.
The two-rung law is the exact origin of the odd/even rung decomposition
(Fibonacci/Lucas branches, W-Thm.\ 8.1), and the four-rung law is the discrete
face of the whitepaper's full-turn cost $\Delta\rho=4\ln\gp=\ln\M$ where
$\M(A_{\gapc})=\gp^4$ (W-\S13): per revolution, co-intensity contracts by
$\gp^{-8}=|W_4|^2$. The imaginary unit written here as $i$ is factored
\emph{internally} from the operator layer in \S\ref{sec:quarterturn}; the
generator is moreover itself an algebraic integer, with minimal polynomial
$x^4+3x^2+1$ and a distinguished place in the relational-charge catalog
(\S\ref{sec:relobject}).
\end{proof}



\begin{theorem}[Interpolation classification]\label{thm:interp}\tagF{} \hx{OZ-O1--O4}
The continuous one-parameter matrix groups through the generator --- the curves
$G(t)=e^{tL}$ with $G(1)=Q$ --- form exactly a $\Z$-family. Solving
$e^{aI+bJ}=Q$ entrywise gives $e^a\cos b=0$ and $e^a\sin b=\tau$, forcing
$\cos b=0$ and $\sin b=+1$ (the branch $\sin b=-1$ is impossible: $e^a=-\tau$
has no real solution), hence $a=\ln\tau$ uniquely and
\[
\log Q=\bigl\{(\ln\tau)I+\bigl(\tfrac\pi2+2\pi k\bigr)J\ :\ k\in\Z\bigr\},
\qquad
W_k(\rho)=e^{-\rho}\,e^{\,i(\pi/2+2\pi k)\rho/\ln\gp}
\quad(t=\rho/\ln\gp).
\]
All branches agree at \emph{every} rung, $W_k(n\ln\gp)=(i\tau)^n$ --- the
discrete orbit cannot distinguish them --- and they are pairwise distinct as
curves (initial velocities $-1+i\kappa_k$ with
$\kappa_k=(\pi/2+2\pi k)/\ln\gp$ all distinct; adjacent branches already
differ at the half-floor, while branches with $k-k'$ even coincide there but
never as functions).
\end{theorem}

\begin{corollary}[What D3 selects]\label{cor:d3select}\tagF{}/\tagD{} \hx{OZ-O5}
The continuous residual state of Definition~\ref{def:W} is the $k=0$ branch,
and within the family the following are equivalent: (i) D3's rate; (ii) the
\emph{principal} logarithm of $Q$; (iii) minimal winding,
$|\pi/2+2\pi k|$ minimized; (iv) the shortest residual spiral per floor,
$L_k=(1-\tau)\sqrt{1+\kappa_k^2}$ minimized (speed
$|W_k'|=e^{-\rho}\sqrt{1+\kappa_k^2}$ integrated over one floor). The
interpolation problem is therefore not undetermined but \emph{determined up to
$\Z$}, with D3 exactly the principal-branch selection; deriving that selection
from the substrate is the same open problem as deriving the rate.
\end{corollary}

\begin{remark}[Transfer to the formation rail]\label{rem:gammatransfer}\tagFD{}/\tagI{} \hx{QX-F1--F6}
The classification of Theorem~\ref{thm:interp} transfers verbatim to the
formation rail: among curves
$\Gamma_{\kappa'}(\rho)=\bigl(r(z(\rho))\cos\kappa'\rho,\
r(z(\rho))\sin\kappa'\rho,\ z(\rho)\bigr)$ whose winding is a one-parameter
rotation group and which pass through every rung site, the admissible rates
are exactly $\kappa'_k=(\pi/2+2\pi k)/\ln\gp$, $k\in\Z$ --- the same
$\Z$-family, because radius and height carry no interpolation freedom ($r$
is inherited, $z$ is D1): the rung-passage condition at $n=1$ already
forces $\kappa'\ln\gp\equiv\pi/2\ (\mathrm{mod}\ 2\pi)$, and all further
rungs follow. All branches agree at every rung site and are pairwise
distinct as curves \hx{QX-F1--F2}; and the $k=0$ branch is the unique one
of shortest \emph{formation-rail} length per floor on both branches of the
profile: the speed $\smash{\sqrt{r'^2+r^2\kappa'^2+z'^2}}$ is strictly
increasing in $\kappa'^2$ wherever $r>0$ \hx{QX-F3--F5}.
Corollary~\ref{cor:d3select}'s list thus gains a sixth member --- \emph{D3
$=$ the shortest formation rail per floor} --- the rail-side companion of
the fifth (shortest total spherical length,
Proposition~\ref{prop:sphlength}). The residual-sphere curves
$S_k=(\Rea W_k,\Ima W_k,z)$ inherit the same classification, all on the
unit sphere \hx{QX-F6}.
\end{remark}

\smallskip\noindent
\textbf{Rate arithmetic and the multiplier torus (v1.7).} The remainder of
this section settles the \emph{arithmetic} of the rate $\kappa=\pi/(2\ln\gp)$
and packages the branch family of Theorem~\ref{thm:interp} into one geometric
object. Scoping first: everything below is \emph{branch-blind} --- every
statement survives the substitution $\kappa\to\kappa_k$ --- so none of it
derives D3. The rate-derivation problem (the session register's OP-RATE)
remains open throughout; what changes is its arithmetic perimeter and its
classification.

\begin{lemma}[Irrationality of the log-ratio]\label{lem:ratirrational}\tagF{} \hx{TX-A1--A7, TX-G1}
$i\pi/\ln\gp\notin\Q$. If $i\pi/\ln\gp=p/q$ with $q\ge1$, exponentiating
gives $(-1)^q=\gp^{\,p}$; but $2\gp^{\,p}=L_p+F_p\sqrt5$ exactly (exact on
the grid $p\in[-12,12]$, general by Binet) with $\sqrt5$-part $F_p/2\neq0$
for $p\neq0$, so $\gp^{\,p}\neq\pm1$ for $p\neq0$ (exact $\Q(\sqrt5)$
decisions), while $p=0$ forces $q\,i\pi=0$, impossible ($q\ge1$; $\pi>3$
interval-certified). Second, premise-independent route: $i\pi/\ln\gp$ is
purely imaginary and nonzero ($\gp>1$ exact), hence not rational.
\end{lemma}

\begin{theorem}[Transcendence of the rate]\label{thm:ratetrans}\tagF{}/\tagE{} \hx{TX-B1--B5}
Forced given the established Gelfond--Schneider theorem (Gelfond 1934;
Schneider 1934, independently; quotient form: for nonzero algebraic
$\alpha,\beta$ with fixed logarithm branches and
$\log\alpha,\log\beta\neq0$, the ratio $\log\alpha/\log\beta$ is rational or
transcendental) \tagE{}:
\[
\boxed{\ \kappa=\frac{\pi}{2\ln\gp}\ \text{is transcendental, and so is
every branch}\ \ \kappa_k=(4k+1)\,\kappa,\ k\in\Z\ }.
\]
Chain: Lemma~\ref{lem:ratirrational} kills the rational horn of G--S at
$(\alpha,\beta)=(-1,\gp)$, so $i\pi/\ln\gp$ is transcendental; the bridge
$i\pi/\ln\gp=2i\,\kappa$ with $2i$ algebraic (root of $x^2+4$) transfers
transcendence to $\kappa$ ($\overline\Q$ is a field); and
$\pi/2+2\pi k=(4k+1)\,\pi/2$ with $4k+1$ a nonzero rational for every
$k\in\Z$. \emph{Restatement-test outcome, stated explicitly:} this is
branch-blind arithmetic --- transcendence holds for every branch at once,
so nothing here selects $k=0$; OP-RATE stays open.
\end{theorem}

\begin{corollary}[Near-miss closure]\label{cor:nearmiss}\tagF{}/\tagE{} \hx{TX-C1--C7, TX-G2--G4}
$4\ln\gp/\pi=2/\kappa$ exactly, and $2/\kappa$ is transcendental
(Theorem~\ref{thm:ratetrans}), while $\tau$ is algebraic (root of the
irreducible $x^2+x-1$): the whitepaper's open item 4 --- the near miss
$4\ln\gp/\pi$ vs $\tau$, margin $5.34\times10^{-3}$ \tagC{}, there
``expected contentless'' --- closes \emph{permanently and negatively}. Both
clauses of the whitepaper's Guard 15.3 --- the whitepaper carries a
\emph{single} Guard 15.3 with both near-miss clauses; there is no Guard
15.4, which also corrects the citation in Remark~\ref{rem:thetaK} ---
promote from [computed-nonequal at 60 dps] to forced-never-equal given
G--S: equality would force $\kappa=2/\tau=2\gp$ (root of $x^2-2x-4$) in the
first clause, and $\kappa=1/\tau^2=\gp^2$ (root of $x^2-3x+1$) in the pitch
clause (via $\tan\theta_K=\tau^2/K$ and injectivity of $\tan$ on
$(0,\pi/2)$) --- both impossible. General statement: \emph{every} near miss
of the form (algebraic) vs (nonzero rational multiple of $\pi/\ln\gp$ or
$\ln\gp/\pi$) is permanently empty. Margins interval-certified at 60 dps:
$0.00533<\tau-4\ln\gp/\pi<0.00534$, $\tau^2-2\ln\gp/\pi>0.07$,
$1/40<\kappa-2\gp<3/100$; all printed decimals of Guard 15.3 reproduced
\hx{TX-G6}.
\end{corollary}

\begin{remark}[$\log_\gp2$ strengthened; scan retirement]\label{rem:logphi2}\tagF{}/\tagE{} \hx{TX-D1--D5, TX-E1--E6, TX-G5}
(a) W-Thm.\ 16.1 strengthens: $\log_\gp2=\ln2/\ln\gp$ is
\emph{transcendental} given G--S at $(\alpha,\beta)=(2,\gp)$. Its
irrationality is now carried by two premise-disjoint mechanisms, meeting
the corpus convergence standard: the whitepaper's Binet $\sqrt5$-part
argument (replicated, grid $p\in[-10,10]$, $q\in[1,8]$, both orientations)
and the new norm obstruction $|N(\gp^{\,p})|=|(\gp\psi)^p|=1$ vs
$N(2^q)=4^q\ge4$, independent of Binet positivity. (b) Scan-retirement
meta-theorem: $\kappa\notin\overline\Q$ (Theorem~\ref{thm:ratetrans}), so
\emph{no algebraic-valued ledger tie for $\kappa$ exists}: for any
nonconstant rational function $R$ over $\overline\Q$ and any algebraic
target $c$, $R(\kappa)=c$ would make $\kappa$ algebraic. The register's
empirical no-tie upgrades to a theorem-level exclusion ($17$ ledger
constants certified algebraic with irreducible integer minimal polynomials;
$\kappa$ and $2/\kappa$ retire together, product $2$; finite corroboration
$|\kappa-c|>1/40$ over the interval ledger, nearest $2\gp$ at
$\approx0.0282$). Honest scope: period relations among the corpus
transcendentals ($\kappa$, $L_{\mathrm{res}}$, and the register's $\tau_0$)
and the algebraic independence of $\pi$ and $\ln\gp$ are beyond
Gelfond--Schneider and remain external-open.
\end{remark}

\begin{proposition}[Modulus ledger of the multiplier torus]\label{prop:modledger}\tagF{} \hx{MX-A1--A4, MX-M1a--M1i, MX-M2a--M2d}
Let $E_q=\C^{\times}/q^{\Z}$ be the multiplier torus of the generator
($0<|q|=\tau<1$ exact, so $E_q$ is nondegenerate). The principal logarithm
and modulus are
\[
\log q=\ln\tau+\frac{i\pi}2=-\ln\gp+\frac{i\pi}2,
\qquad
\mu=\frac{\log q}{2\pi i}=\frac14+\frac{i\ln\gp}{2\pi}\ (\mathrm{mod}\ \Z),
\]
forced by exp-match and principal-strip membership ($\Ima\in(-\pi,\pi]$;
branches differ by $2\pi ik$, at least the strip diameter). Ledger:
$\Rea\mu=\tfrac14$ is the charge quantum --- the $c=\tfrac14$ gauge
rotation of Theorem~\ref{thm:gauge}, one charge quantum per rung;
$\Ima\mu=\ln\gp/(2\pi)>0$;
\[
\frac{\Rea\mu}{\Ima\mu}=\kappa
\qquad\text{and}\qquad
|\log q|^2=\ln^2\gp+\frac{\pi^2}4=\ln^2\gp\,\bigl(1+\kappa^2\bigr),
\]
the per-floor arc constant of Corollary~\ref{cor:d3select}(iv) read in the
log chart: $(-1+i\kappa)\ln\gp=\log q$ exactly.
\end{proposition}

\begin{theorem}[Systolic marking: the seventh characterization]\label{thm:systole}\tagF{} \hx{MX-M3a--M3f, MX-M4a--M4c}
For the period lattice $\Lambda=\Z\log q+\Z\,2\pi i$ of $E_q$,
\[
|n\log q+2\pi i\,m|^2=\ln^2\gp\,\bigl(n^2+\kappa^2(n+4m)^2\bigr)
\]
exactly, so every length comparison is linear in $\kappa^2$ with integer
coefficients, decided by exact rational arithmetic against the single
certified sandwich $10<\kappa^2<11$ (one interval guard, \hx{MX-G1}).
$\Lambda$ is \emph{branch-independent}: $\{\log_kq,\ 2\pi i\}$ is a
unimodular change of basis of $\Lambda$ for every $k$, and the length
spectrum is $k$-invariant; the \emph{marking} (choice of generator) is not
($\log_kq-\log q=2\pi ik$). The systole of $E_q$ is $|\log q|$, attained
exactly at $\pm\log q$ --- the principal-branch generators --- hence
\begin{center}
\emph{D3 $\iff$ the marking whose non-meridian generator is the systole
representative.}
\end{center}
The second-shortest class is the four-rung cycle: $4\log q-2\pi i=\ln\gapc$,
length $4\ln\gp$ --- the charge cycle $W_4=\gapc$ of
Theorem~\ref{thm:generator} \emph{is} the second-shortest closed geodesic.
\emph{Stated prominently: this is characterization number seven for
Corollary~\ref{cor:d3select}'s menu, not a derivation.} Per the restatement
test the selection content is exactly a minimality principle (the systole);
every $k$-invariant part of the construction cannot separate branches, and
OP-RATE's derivation problem is untouched.
\end{theorem}

\begin{remark}[Geodesic ledger; no-CM]\label{rem:geoledger}\tagF{}/\tagE{}
(a) \hx{MX-M5a--M5c, MX-G2--G3}\ The closed-geodesic length ledger of
$E_q$ is certified: lengths
$\ln\gp\,\sqrt{a+b\,\kappa^2}$ with $(a,b)=\bigl(n^2,(n+4m)^2\bigr)$; the
first ten distinct classes (up to sign, iterates included), ordered by
exact sandwich arithmetic and interval-corroborated at 60 dps, are
\[
(1,1),\ (16,0),\ (9,1),\ (25,1),\ (4,4)^{\times2},\ (49,1),\ (64,0),\
(36,4)^{\times2},\ (81,1),\ (1,9),
\]
with $(1,1)$ the systole, $(16,0)=4\ln\gp$ the $W_4=\gapc$ cycle, and
$(4,4)$ twice the systole (an iterate/primitive degeneracy); window
sufficiency ($|n|\le12$, $|m|\le4$) closed by exact tail identities. This
is the certified ledger on which $\pi/2$ and $\ln\gp$ combine
\emph{additively} --- the natural habitat for any future extremality
selection, which must re-pass the restatement test. (b)
\hx{MX-M6a--M6c}\ Given
Theorem~\ref{thm:ratetrans}, $\mu$ is not algebraic --- the exact bridge is
$(\mu-\tfrac14)\,\kappa=i/4$ with $\mu-\tfrac14=i\ln\gp/(2\pi)\neq0$ --- in
particular not imaginary quadratic: \emph{$E_q$ has no complex
multiplication} (standard CM criterion; conditional exactly on the
transcendence input).
\end{remark}

\begin{proposition}[$k$-invariance of the discrete state]\label{prop:kinv}\tagF{} \hx{DX-E1--E3}
The full discrete state of the ladder --- the rung values
$e^{i\kappa_kn\ln\gp}=i^{\,n}$ (exact on the grid $k\in[-3,3]$,
$n\in[0,8]$, symbolically via $e^{2\pi ikn}=1$), the heights
$z_n^2=1-\tau^{2n}$, the transfer ledger, the gate data
$(C_n,u_n,m_n,\lambda_n)$, the floor lattice $\rho_n=n\ln\gp$, and the
torus lattice $\Lambda$ --- is \emph{identical for every branch} $k\in\Z$.
Only the marking $\log_kq$ is $k$-dependent (Theorem~\ref{thm:systole}).
\end{proposition}

\begin{remark}[OP-RATE at rung resolution: sign ladder, exposed joint, ledger kill, classification no-go]\label{rem:oprate}\tagF{}/\tagO{} \hx{DX-C1--C7, DX-D1--D5, DX-E4--E5}
Four bookkeeping results at the rate problem's core; items (i) and (iii)
are forced, (ii) and (iv) are candidate-level; none closes the problem.
(i) \emph{Sign ladder, D1-free} \hx{DX-C1--C2}: symbolically
$m(w)=-C/(1+u)^2=-1/w<0$ with $w=\gp^{2n}$, so the gate multiplier is
negative at \emph{every} gate and does not alternate between gates; the
two-rung carrier is $q^2=-\tau^2$, argument $\pi$. (ii) \emph{Exposed-joint
axiom H} --- a formalized open joint, presented as such, \emph{not} a
result \hx{DX-C3--C6}: given
\begin{center}
H $\coloneqq$ ``argument of the two-rung ladder action $=$ transverse
holonomy accumulated over two floors ($=\pi$)'',
\end{center}
exact mod-$2\pi$ bookkeeping forces the per-floor step
$\delta=\pi/2\ (\mathrm{mod}\ \pi)$ --- the quantum $\pm\pi/2$ at rung
resolution --- but the residual $\pi\Z$ is not pinned: every branch
satisfies H, and the branch remains untouched. (iii) \emph{Ledger
counter-models} \hx{DX-D1--D5}: the transfer ledger
$\Delta F_n=\tau^{2n+1}$ (Proposition~\ref{prop:ledger}) does \emph{not}
force a constant angular step. Two explicit counter-models pass every
ledger identity: $\theta_n=n\pi$ (all moduli correct; fails rung passage
at $n=1$, $e^{i\pi}\tau=-\tau\neq i\tau$) and $\theta_n=n\pi/2+2\pi s_n$
with a non-constant integer lift $s_n$ (reproduces \emph{all} rung values
$i^{\,n}$ and all ledger identities). Constancy of the angular step is
therefore a lift/equivariance assumption, and the ledger route to the rate
dies by contamination. (iv) \emph{R6 no-go --- a candidate classification
statement, not asserted as a theorem} \hx{DX-E4--E5}: any rate-selection
functional that factors through the $k$-invariant discrete data of
Proposition~\ref{prop:kinv} is constant in $k$ and cannot select $k=0$;
the six existing characterizations (Corollary~\ref{cor:d3select}(i)--(iv),
Proposition~\ref{prop:sphlength}, Remark~\ref{rem:gammatransfer}) all
reference \emph{continuum} objects --- winding number, spiral speed, rail
length --- which is exactly the escape, and the separation is exhibited
($L_k$ separates branches that the discrete functionals cannot). This
supports closing OP-RATE \emph{as a classification} --- determined up to
$\Z$ at rung resolution, with selection necessarily continuum-referencing
--- at candidate level; OP-RATE \emph{as a derivation} remains open
\tagO{}.
\end{remark}

\smallskip\noindent
\textbf{OP-RATE closes as a classification (v1.8).} The candidate R6 no-go of
Remark~\ref{rem:oprate}(iv) is now promoted to a forced theorem, and its
consequence stated precisely: at rung resolution the branch $k$ --- hence the
rate $\kappa$ up to the $\Z$-family --- cannot be selected by any functional of
the discrete substrate data. This settles OP-RATE \emph{as a classification};
it is a no-go on discrete-substrate \emph{derivation}, not a derivation of
$\kappa$, and the continuum rate D3 stays declared throughout.

\begin{lemma}[Extended $k$-invariance of the discrete state]\label{lem:kinvext}\tagF{} \hx{CL-1a--CL-1h, CL-1F}
The $k$-invariance of Proposition~\ref{prop:kinv} extends, re-forced cold, to
the \emph{full} discrete state of the ladder. For every branch $k\in\Z$ the
following are identical to their $k=0$ values: the rung values
$e^{i\kappa_kn\ln\gp}=i^{\,n}$ (exact grid $k\in[-4,4]$, $n\in[0,10]$, and
symbolically via $e^{2\pi ikn}=1$); the heights $z_n^2=1-\tau^{2n}$; the gate
ladder $(C_n,u_n,m_n,\lambda_n)$ with $m_n\lambda_n=1$; the floor lattice
$\rho_n=n\ln\gp$; the transfer ledger $\Delta F_n=\tau^{2n+1}$ (sum $1$); the
torus lattice $\Lambda=\Z\log q+\Z\,2\pi i$, re-verified branch-independent by
the unimodular change of basis $\bigl(\begin{smallmatrix}1&k\\0&1\end{smallmatrix}\bigr)$
($\det=1$); the ten-class geodesic length ledger
$\ln\gp\sqrt{a+b\kappa^2}$ with $(a,b)=\bigl(n^2,(n+4m)^2\bigr)$, invariant
under the marking relabel $(n,m)\mapsto(n,m+nk)$; and the \emph{systole value}
$|\log q_0|^2=\ln^2\gp\,(1+\kappa^2)=\ln^2\gp+\pi^2/4$, the lattice minimum.
Only the \emph{marking} $\log_kq$ is $k$-dependent
($|\log_1q|^2-|\log_0q|^2=6\pi^2\ne0$, detected \hx{CL-1F}). The lemma consumes
neither D2 nor D3 (D1 enters only through the $k$-free rung modulus $\tau$).
\end{lemma}

\begin{theorem}[R6 classification no-go]\label{thm:r6nogo}\tagF{} \hx{CL-2a, CL-2F1--CL-2F2}
Define a \emph{rung-resolution selection functional} as any map
$\Phi=g\circ D$ from the $k$-invariant discrete data $D$ of
Lemma~\ref{lem:kinvext} into a totally ordered set, whose argmin/argmax
purports to select the branch $k$. Then every such functional is
\emph{constant on the $\Z$-family}, and so selects no branch: if $D(k)=D(0)$
for all $k$ (Lemma~\ref{lem:kinvext}), then $\Phi(k)=g(D(k))=g(D(0))=\Phi(0)$,
whose argmin is the whole $\Z$-family. (Witness: $F=\sum_{n<6}e^{-\rho_n}i^{\,n}$
is $k$-constant symbolically \hx{CL-2a}; the non-vacuity guard \hx{CL-2F1}
fires on any core secretly made $k$-constant.) This promotes
Remark~\ref{rem:oprate}(iv) from candidate to forced. Per the restatement
test the no-go is branch-blind \emph{by design}: it is a
characterization-of-impossibility, not a derivation of $k$.
\end{theorem}

\begin{corollary}[OP-RATE closes as a classification]\label{cor:opratelass}\tagF{} \hx{CL-3a--CL-3c}
OP-RATE closes \emph{as a classification}. What closes: the reading
``derive $\kappa$ from the discrete substrate at rung resolution'' is a no-go
(Theorem~\ref{thm:r6nogo}) --- $\kappa$ is declared-up-to-$\Z$ at the discrete
level, and selecting the principal branch \emph{provably requires} a
continuum-referencing minimality principle. The seven characterizations
(Corollary~\ref{cor:d3select}(i)--(iv), Proposition~\ref{prop:sphlength},
Remark~\ref{rem:gammatransfer}, Theorem~\ref{thm:systole}) are the
\emph{complete} discrete-level menu: the discrete level contains no selector
at all. What does \emph{not} close: the continuum rate D3 itself remains
declared --- all branch curves $W_k$ are admissible one-parameter groups
through $Q$, pairwise distinct as curves (Theorem~\ref{thm:interp}, initial
velocities $-1+i\kappa_k$) --- and a continuum-referencing derivation is not
excluded. This mirrors the $\lambda=2c$/Cencov equivariance no-go of the
companion $\lambda=2c$ paper (\texttt{2026-06-lambda-2c}, thm:cencov) exactly:
as there ``$c$ cannot be fixed by any invariance requirement \ldots\ $c$ is a
declared positive scale \ldots\ shipped decisions are $c$-invariant across the
band \ldots\ [but he] does not forbid an external structural constraint from
doing so'', here $k$ cannot be fixed by any functional of the $k$-invariant
discrete data, $\kappa$ is declared-up-to-$\Z$ with discrete decisions
$k$-invariant across the $\Z$-family, and a continuum-referencing principle may
still select. OP-RATE \emph{as a derivation} remains open \tagO{}.
\end{corollary}

\begin{remark}[The seven-way continuum/discrete table]\label{rem:sevenway}\tagF{} \hx{CL-2b(i)--(vii), CL-2c}
The content of the no-go is that each of the seven characterizations is
\emph{continuum-referencing}: its selection functional is $k$-dependent
(unique argmin at $k=0$), so it provably does not factor through the
$k$-invariant discrete data. (i) D3's declared continuum rate $\kappa$ itself
--- a \emph{restatement} (the selector $|\kappa_k-\kappa|=4|k|\kappa$ consumes
the target); (ii) the principal branch of $\log Q$, i.e.\ the continuous
argument/lift $\theta(\rho)$ (the rung data is lift-blind, $e^{i\theta}=i^n$
only); (iii) the minimal winding $|\pi/2+2\pi k|$ of the continuous curve;
(iv) the arc length $L_k=(1-\tau)\sqrt{1+\kappa_k^2}$ of the continuous
residual spiral; (v) the total spherical length $L_{\mathrm{res}}$; (vi) the
formation-rail length per floor; (vii) the systolic selection --- the flat
norm $|\log_kq|^2=\ln^2\gp(1+(4k+1)^2\kappa^2)$ of the $k$-dependent
\emph{marking} on the continuous torus $E_q$, even though the lattice and the
systole value are $k$-invariant (Lemma~\ref{lem:kinvext}). Tag discipline: the
structural conclusion ``continuum-referencing, does not factor through the
discrete data'' is \tagF{} for all seven, and rows (i)--(iv),(vii) are forced
at the exact level (integer/rational cores $(4k+1)^2$, $|4k+1|$, $4|k|$, plus
$\kappa^2>0$); the strict per-characterization ``unique argmin at $k=0$''
\emph{selection} for rows (v) and (vi) is \tagF{}\emph{ given corpus} ---
inherited, cited not re-encoded, from Proposition~\ref{prop:sphlength}
($L_{\mathrm{res}}$ strictly increasing in $|\kappa|$) and
Remark~\ref{rem:gammatransfer} (rail speed strictly increasing in
$\kappa'^2$). The no-go's core needs only that $L_{\mathrm{res}}$ and the rail
are continuous integrals, manifestly not functions of the rung samples, which
is exact. The whole block is branch-blind by design: that is the content.
\end{remark}

\smallskip\noindent
\textbf{The period-relation map (v1.8).} The scan-retirement remark
(Remark~\ref{rem:logphi2}(b)) left the period relations among the corpus
transcendentals ($\kappa$, $L_{\mathrm{res}}$, and the register's Salem-slot
constant $\tau_0$) external-open and undifferentiated. This block pins what is
algebraic, settles the $\kappa$-vs-$\tau_0$ worry, and maps the precisely
external-open residue.

\begin{proposition}[$\tau_0$ algebraic; no $\kappa$-vs-$\tau_0$ relation]\label{prop:tau0alg}\tagF{}/\tagE{} \hx{PM-1B1--PM-1B9, PM-1X1--PM-1X6, PM-A2}
The Salem-slot occupant $\tau_0=\beta+\beta^{-1}$ (the trace redirection of the
companion salem-slot note, Def.~3.2; the dominant root of the degree-$2$
trace-down) is \emph{algebraic}. For the least-degree Salem number $\beta_4$
(minimal polynomial $x^4-x^3-x^2-x+1$) it has the exact minimal polynomial
\[
\boxed{\ \operatorname{minpoly}(\tau_0)=t^2-t-3,\qquad \tau_0=\tfrac12(1+\sqrt{13})\ },
\]
a degree-$2$, irreducible, monic-integer (algebraic-integer), \emph{totally
real} (discriminant $13>0$, Sturm real-root count $2$), \emph{Perron}
($\tau_0-|\text{conjugate}|=1>0$) constant with $\tau_0>2>\gp$ (salem-slot
Thm.~3.3); the type-fact is general over all Salem redirections, and the
golden limit $\sqrt5=\gp+\gp^{-1}$ (minpoly $t^2-5$) and the Pisot
accumulation $\mu_P+\mu_P^{-1}$ (minpoly $t^3+t^2-4t-5$) are algebraic too.
Combined with $\kappa$ transcendental (inherited, Theorem~\ref{thm:ratetrans},
given Gelfond--Schneider), there is \emph{no} algebraic relation over $\Q$
between $\kappa$ and $\tau_0$: for $m(Y)=\operatorname{minpoly}(\tau_0)$ the
resultant $R(X)=\operatorname{Res}_Y(P,m)=\prod_r P(X,r)\in\Q[X]$, and a
genuine (non-vacuous) $P(\kappa,\tau_0)=0$ would give $R(\kappa)=0$ with $R\ne0$,
forcing $\kappa\in\overline\Q$ --- impossible. The register's ``$\kappa$ vs
$\tau_0$'' period worry is thus settled \emph{negatively} at the algebraic
level (a specialization of the scan-retirement meta-theorem to two variables;
$\kappa$ is consumed only as an excluded value, not derived, so this is
branch-blind exclusion, not restatement).
\end{proposition}

\begin{remark}[$L_{\mathrm{res}}$ as an elliptic period; the external-open ledger]\label{rem:periodmap}\tagF{}/\tagE{}/\tagO{} \hx{PM-1C1--PM-1C2, PM-2A, PM-2B, PM-2Bf}
The finite spherical length $L_{\mathrm{res}}=\sqrt{1+\kappa^2}\,
E\!\bigl(\kappa^2/(1+\kappa^2)\bigr)=3.7312193125\ldots$
(Proposition~\ref{prop:sphlength}) is the value of the complete elliptic
integral of the second kind at the elliptic \emph{parameter}
$m=\kappa^2/(1+\kappa^2)$ (modulus $k=\sqrt m$), via the exact reduction
$1+\kappa^2\cos^2\theta=(1+\kappa^2)(1-m\sin^2\theta)$. At an \emph{algebraic}
parameter $E$ is a Kontsevich--Zagier period (an algebraic ellipse's
quarter-perimeter); here the M\"obius inverse $\kappa^2=m/(1-m)$ makes $m$
\emph{transcendental} (since $\kappa$ is), so $L_{\mathrm{res}}$ is a period
\emph{function} evaluated at a transcendental argument --- one step beyond
Gelfond--Schneider and beyond the Schneider/Chudnovsky elliptic-period
theorems (which need an algebraic modulus or CM lattice), hence its
transcendence is \emph{not} settled here. Since $\tau_0$ is algebraic
(Proposition~\ref{prop:tau0alg}), any relation $P(\tau_0,L_{\mathrm{res}})=0$
over $\Q$ would force $L_{\mathrm{res}}\in\overline\Q$ by the same
resultant/tower argument \hx{PM-2B}, so the $(\tau_0,L_{\mathrm{res}})$ pair
\emph{reduces} to the single constant. The precise ledger: \emph{settled}
--- $\kappa$ transcendental, $\tau_0$ algebraic, no $(\kappa,\tau_0)$
relation, and $(\tau_0,L_{\mathrm{res}})$ reduced to ``is $L_{\mathrm{res}}$
transcendental''; \emph{genuinely external-open} \tagO{} --- the transcendence
of $L_{\mathrm{res}}$, and the algebraic independences of $(\kappa,
L_{\mathrm{res}})$ and of $(\pi,\ln\gp)$ (Schanuel / Kontsevich--Zagier
territory, beyond Gelfond--Schneider and Baker; the ratio $\pi/\ln\gp=2\kappa$
is settled transcendental but pair-independence is strictly stronger). The
entire map survives $\kappa\mapsto\kappa_k$ (branch-blind) and advances no
selection.
\end{remark}

\begin{remark}[The forced-identity ledger of the period frontier]\label{rem:pfledger}\tagF{}/\tagFD{} \hx{PF-3A--PF-3G, PF-3Ef, PF-G4, PF-G5}
The period frontier of Remark~\ref{rem:periodmap} sits on an exact algebraic
skeleton, every relation of which is forced. \emph{Golden trace} \tagF{}:
$\gp^2=\gp+1$, $\gp\tau=1$, $\tau^2+\tau=1$, and $\gp+\gp^{-1}=\sqrt5$ (the
salem-slot golden floor). \emph{The definitional tie} \tagFD{}: given the
declared rate $\kappa=\pi/(2\ln\gp)$ [D3],
\[
\frac2\kappa=\frac{4\ln\gp}{\pi},\qquad \frac{\pi}{\ln\gp}=2\kappa,\qquad
\pi=2\kappa\ln\gp,
\]
and the branch family $\kappa_k=(\pi/2+2\pi k)/\ln\gp=(4k+1)\kappa$ (symbolic in
$k\in\Z$, with $4k+1$ a nonzero odd rational for every $k$). \emph{The Salem
anchor} \tagF{}\emph{ given corpus}: $\tau_0=(1+\sqrt{13})/2$ satisfies
$\tau_0^2=\tau_0+3$ (minimal polynomial $t^2-t-3$, discriminant $13$), the
degree-$2$ algebraic-integer occupant of the redirection slot
(Proposition~\ref{prop:tau0alg}). \emph{The lens gate} \tagF{}: $K^2=1-\gapc$,
$\gapc=\gp^{-4}$, with minimal polynomial $x^2+5x-5$ and $K^2>0$ (exact
$\Q(\sqrt5)$ sign). \emph{The elliptic reduction} \tagF{}: symbolic in
$K,\theta$,
\[
1+K^2\cos^2\theta=(1+K^2)\bigl(1-m\sin^2\theta\bigr),\qquad
m=\frac{K^2}{1+K^2},\qquad K^2=\frac{m}{1-m},
\]
so $L_{\mathrm{res}}=\int_0^{\pi/2}\sqrt{1+\kappa^2\cos^2\theta}\;d\theta
=\sqrt{1+\kappa^2}\,E(m)$ (perturbing $m\mapsto m+10^{-3}$ breaks the identity;
falsifier \hx{PF-3Ef}). \emph{A forced geometric reading} \tagF{} (new this
session): $L_{\mathrm{res}}$ is the quarter-perimeter of the ellipse with
semi-axes $(\sqrt{1+\kappa^2},\,1)$; its linear eccentricity $c$ satisfies
$c^2=a^2-b^2=\kappa^2$ \emph{exactly} (the focal parameter is $\kappa$), and its
eccentricity-squared is $c^2/a^2=m$ \emph{exactly} (the elliptic parameter). Two
independent computations of $L_{\mathrm{res}}$ --- the closed form and direct
quadrature of the integrand --- agree to $<10^{-200}$ \hx{PF-G4}, so the frontier
is numerically well-posed; every identity here survives $\kappa\mapsto\kappa_k$
verbatim.
\end{remark}

\begin{remark}[The frontier map: Baker's reach, Schanuel's conditional, and the external-open tool-gaps]\label{rem:pffrontier}\tagE{}/\tagO{}/\tagC{} \hx{PF-2A--PF-2E, PF-1P1--PF-1P6, PF-1N1--PF-1N7, PF-G1--PF-G3}
Against this skeleton the transcendence questions sort into exactly three tiers,
each with its deciding tool or its named gap. \emph{Settled} \tagE{}: $\kappa$ is
transcendental (Gelfond--Schneider at $(-1,\gp)$, since
$\kappa=(i\pi/\ln\gp)/2i$ with $2i$ algebraic; inherited
Theorem~\ref{thm:ratetrans}); $\tau_0$ is algebraic ($t^2-t-3$); and
\emph{Baker's reach is pinned} --- $\{1,\ln2,\ln3,\ln\gp\}$ is
$\overline\Q$-linearly independent \tagE{}, because $2,3,\gp$ are
multiplicatively independent (the exact box $|a|,|b|,|c|\le6$ has
$2^a3^b\gp^c=1$ only at the origin: for $c\neq0$ the $\sqrt5$-part $F_c/2\neq0$
makes $\gp^c$ irrational, and $2^a3^b=1$ forces $a=b=0$) and Baker's
linear-forms theorem then applies. \emph{Conditional on Schanuel} \tagO{}:
$\{\ln\gp,\,i\pi\}$ is $\Q$-linearly independent ($\ln\gp$ real nonzero, $i\pi$
purely imaginary nonzero) with algebraic exponentials $\gp,-1$, so Schanuel
would give $\operatorname{trdeg}\ge2$ and hence $(\pi,\ln\gp)$ --- equivalently
$(\kappa,\ln\gp)$, since $\Q(\kappa,\ln\gp)=\Q(\pi,\ln\gp)$ --- algebraically
independent; the premise is forced but the conjecture is not, so this is
\emph{not} a theorem here. \emph{Genuinely external-open} \tagO{}, with the
tool-gap named in each case: the transcendence of $L_{\mathrm{res}}$ is untouched
because its elliptic modulus $k=\sqrt m$ is \emph{transcendental}
($m=\kappa^2/(1+\kappa^2)$, and the M\"obius inverse $K^2=m/(1-m)$ makes
$m\in\overline\Q\iff\kappa\in\overline\Q$), so the Schneider/Chudnovsky
elliptic-period theorems --- which require an algebraic modulus or CM lattice ---
are \emph{inapplicable}; and the algebraic independence of
$(\kappa,L_{\mathrm{res}})$ lies beyond even Schanuel (an elliptic-integral value
outside the exp/log framework, Kontsevich--Zagier / period-conjecture territory).
\emph{We resolve none of these.} As \tagC{} corroboration only, a
validated-detector PSLQ search --- fail-first on six known relations (including a
size-$6$ planted relation matching the flagship basis) and correctly silent on
two transcendental controls \hx{PF-1P1--PF-1P6, PF-1N1--PF-1N2} --- finds
\emph{no} integer relation of height $\le10^{12}$ at $230$ dps among the frontier
bases $\{1,\kappa,\kappa^2,L_{\mathrm{res}},L_{\mathrm{res}}^2,
\kappa L_{\mathrm{res}}\}$, $\{1,\pi,\ln\gp,\pi\ln\gp\}$,
$\{1,\ln\gp,\ln2,\ln3\}$, and the $L_{\mathrm{res}}$ algebraicity ladders to
degree $6$ \hx{PF-1N3--PF-1N7}. This \emph{corroborates, and never proves}, the
open transcendence and independence statements; PSLQ silence is evidence, not a
certificate. The entire map is branch-blind ($\kappa\mapsto\kappa_k$ verbatim)
and selects no $k$ --- a characterization, not an OP-RATE derivation.
\end{remark}

\section{The canonical quarter-turn: $J=UA$}\label{sec:quarterturn}

Everything in this section is unconditional \tagF{}: pure keystone/Clifford
algebra in the whitepaper's fixed convention
($R=\bigl(\begin{smallmatrix}0&1\\1&1\end{smallmatrix}\bigr)$, $H=2R-I=2e_1-e_2$,
$e_1=\sigma_x$, $e_2=\sigma_z$, $J=e_1e_2$; W-\S3--4), using no part of D1--D3.
It resolves the canonicality problem the July 14 report posed for itself: the
representation $W_n=(i\tau)^n$ imported $i$ from the declared winding; here the
quarter-turn is \emph{factored from two real, normalized structures already
forced in the operator layer}.

\begin{proposition}[Spectral uncoupling]\label{prop:uncouple}\tagF{} \hx{OZ-J1--J5}
The spectral projectors of the keystone,
\[
P_+=\frac{R+\tau I}{\sqrt5},\qquad P_-=\frac{\gp I-R}{\sqrt5},
\]
are complete orthogonal idempotents ($P_++P_-=I$, $P_\pm^2=P_\pm$, $P_+P_-=0$)
with $RP_+=\gp P_+$ and $RP_-=\psi P_-$, $\psi=-\tau$. The normalized
\emph{keystone grading}
\[
U \coloneqq P_+-P_-=\frac{H}{\sqrt5}=\frac{2R-I}{\sqrt5}
\]
satisfies $U^2=I$ and acts as $+I$ on the expanding golden mode and $-I$ on the
contracting one; equivalently $P_\pm=(I\pm U)/2$.
\end{proposition}

\begin{proposition}[The return-kernel recoupler]\label{prop:recoupler}\tagF{} \hx{OZ-K1--K4, OZ-L4}
Let $N=e_1+2e_2$ be the generator of the vector part of $\ker\mathcal L$,
$\mathcal L(X)=\tfrac12\{H,X\}$ (W-Thm.\ 4.2). Then $N^2=5I$ and, normalizing by
the Clifford norm, $A \coloneqq N/\sqrt{N^2}=N/\sqrt5$ satisfies
\[
A^2=I,\qquad UA=-AU,\qquad AP_+=P_-A,\qquad AP_-=P_+A :
\]
$A$ is a real involution that \emph{exchanges} the two golden modes. The
anticommutation is automatic, not a coincidence: on vectors,
$\mathcal L(v)=\langle H,v\rangle\,I$ (W-Thm.\ 4.2), so the vector kernel is
\emph{by definition} $H^\perp$ --- membership in the return kernel already
certifies orthogonality to $H$, hence $HN=-NH$.
\end{proposition}

\begin{theorem}[Canonical quarter-turn factorization; Strouse--Ozymandias]\label{thm:factorization}\tagF{} \hx{OZ-K5--K6}
With $U$ and $A$ as above,
\[
\boxed{\ J=UA=\frac{H}{\sqrt5}\cdot\frac{N}{\sqrt5}=\frac{HN}{5}\ },
\qquad J^2=UAUA=-U^2A^2=-I .
\]
Matrix witness: $H=\bigl(\begin{smallmatrix}-1&2\\2&1\end{smallmatrix}\bigr)$,
$N=\bigl(\begin{smallmatrix}2&1\\1&-2\end{smallmatrix}\bigr)$,
$HN=5\bigl(\begin{smallmatrix}0&-1\\1&0\end{smallmatrix}\bigr)=5\,e_1e_2$.
Uniqueness: $U$ is fixed by the keystone; the vector kernel is one-dimensional,
so its unit generator is $\pm A$; hence
\[
J=\pm\,UA
\]
is canonical up to the orientation sign --- exactly the freedom the whitepaper's
D2 already names as a convention. The quarter-turn no longer needs to be
imported: \emph{distinguish the modes ($U$) $+$ exchange the modes ($A$) $=$
quarter-turn ($J$)}.
\end{theorem}

\begin{corollary}[Coordinate-free form]\label{cor:coordfree}\tagF{} \hx{OZ-L3--L4}
In $\mathrm{Cl}(2,0)$ the product of any positively-oriented orthonormal pair of
vectors equals the pseudoscalar: for every rotation angle $\theta$,
$e_1'e_2'=e_1e_2=J$ exactly (a reflection flips the sign). Since
$(U,A)=(\widehat H,\widehat N)$ is an orthonormal pair \emph{by construction}
(Prop.~\ref{prop:recoupler}), the factorization $J=\pm UA$ is basis-independent
up to orientation; no choice beyond the whitepaper's fixed convention enters.
For an arbitrarily scaled kernel generator the normalization is
$A=N/\sqrt{N^2}$; with $N=e_1+2e_2$ this is $N/\sqrt5$. This closes the report's
open item (1).
\end{corollary}

\begin{remark}[Why the return layer is necessary]\label{rem:obstruction}\tagF{} \hx{OZ-L1--L2}
The commutative algebra $\R[R]=\operatorname{span}\{I,R\}$ contains no square
root of $-I$: $(aI+bR)^2=-I$ has no real solutions (exact elimination). Nor does
it contain the recoupler: everything in $\R[R]$ commutes with $H$, while
$[H,N]=2HN=10J\neq0$. The quarter-turn therefore cannot come from the keystone
\emph{alone}; it requires the keystone \emph{acting on} the matrix algebra ---
the return operator's module structure --- which is precisely where $N$ lives.
The report's initial self-criticism (``clean representation $\neq$ canonical
derivation'') identified the right gap, and the factorization fills it from
inside the existing layer.
\end{remark}

\begin{definition}[Canonical rung generator]\label{def:Q}\tagF{} \hx{OZ-K7--K9}
\[
Q \coloneqq \pm\,\tau J=\pm\,\tau UA=\pm\,\frac{\tau}{5}\,HN,
\qquad
Q^2=-\tau^2I,\qquad Q^4=\gapc\,I,\qquad Q^n=\tau^nJ^n .
\]
Under the algebra isomorphism $\R[J]\cong\C$ ($J\mapsto i$), $Q^n\cong(i\tau)^n
=W_n$; concretely, with basepoint $v_0=(1,0)^{\!\top}$,
\[
Q^nv_0=\bigl(\Rea W_n,\ \Ima W_n\bigr)^{\!\top},
\qquad\text{so}\qquad
S_n=\bigl(Q^nv_0\,;\ z_n\bigr):
\]
the residual sphere's horizontal part (\S\ref{sec:sphere}) \emph{is} the
$Q$-orbit of the east basepoint.
\end{definition}

\smallskip\noindent
\textbf{Scoping upgrade (what this does and does not close).} Before v1.1 the
generator was forced \emph{given} D1--D3, with $i$ supplied by D2. Now: the
quarter-turn operator $J$ and the rung generator $Q=\pm\tau J$ exist
canonically (up to orientation) in the substrate itself \tagF{}, and D2's
remaining declared content is \emph{selection}, not existence --- that the
winding applies exactly one $Q$ per entropy floor (D3's rate), rather than the
terrain $\pi$ or pentagon $2\pi/5$ clocks. Per the report's own ledger, still
open: rung quantization (why exactly one $Q$ per rung event --- by
Corollary~\ref{cor:d3select} this now \emph{contains} the interpolation
question, since branch selection $=$ rate selection), derivation of the
inherited radius law $r(z)$, and any physical identification (the
helicase/uncouple--recouple language was a mechanism-finding analogy only).
Coordinate-free uniqueness is closed by Corollary~\ref{cor:coordfree};
continuous interpolation is classified by Theorem~\ref{thm:interp}.

\smallskip\noindent
\textbf{The $\chi$-selection front (v1.7).} Three results sharpen the
selection scoping just stated: what the emission image $S$ (the closure of
the whitepaper's emission catalog, W-\S5) can and cannot supply for the
angular quantum $\chi$, and exactly where a $\chi$-selection argument must
stand.

\begin{proposition}[Pentagon exclusion by the image]\label{prop:pentimage}\tagF{} \hx{CX-A3, CX-B1--B5, CX-D6}
$\zeta_5$ is not in the emission image $S$: $\arg\zeta_5=2\pi/5$ lies
outside the image angle lattice $(\pi/2)\Z$ (W-Lemmas 5.4--5.5). Exactly:
$\tfrac15\notin\tfrac14\Z$ mod $1$; the pentagon classes $k/5$ ($k=1..4$)
meet $\{0,\tfrac14,\tfrac12,\tfrac34\}$ only at $0$; $\Phi_5$ (verified
$=\operatorname{minpoly}(\zeta_5)$) has no root in $\mu_4$; and the
pentagon gate multipliers are primitive fifth roots (pair product
$=\Phi_5$). This \emph{strengthens} the whitepaper's pentagon exclusion
from the principle to the image: W-Decl.\ 12.2/W-Prop.\ 14.2 exclude the
$2\pi/5$ clock \emph{by the principle} D2; here the pentagon multiplier is
not an on-circle element of $S$ at all, independently of D2.
\end{proposition}

\begin{proposition}[Unimodular tensor-monomials reach only $\pm1$]\label{prop:unimon}\tagF{} \hx{CX-A4, CX-C1--C7}
Every unimodular tensor-monomial of catalog eigenvalues has value exactly
$+1$ or $-1$: $\pm i$ is multiplicatively unreachable. Forcing: every
catalog eigenvalue has $5$-valuation $w_2\ge0$ (moduli of the form
$2^{w_0/4}3^{w_1/4}5^{w_2/4}\gp^{w_3/4}$); the only odd-quarter-argument
carriers are the K-seed's rotation pair $\pm i\beta$, each with $w_2=1>0$;
unimodularity forces the total $w_2=0$ (norming to $\Q$ with $N(\gp)=-1$
and $\gp$ not a root of unity), hence no $\pm i\beta$ factor, hence even
argument. Corroborated exhaustively on all $74613$ monomials of degree
$\le6$ ($60$ unimodular, all $\pm1$). Whether the \emph{full} image ---
closed also under addition, minimal polynomials, and $\Phi$-extraction ---
can reach $\pm i$ is \emph{not} decided by this multiplicative argument and
remains open \tagO{} \hx{CX-C6--C7}: the corpus-exhibited host is
$x^4-1=\Phi_1\Phi_2\Phi_4$ with $\Phi_4$ rooted at $\pm i$.
\end{proposition}

\begin{theorem}[The conditional $\chi$-selection chain, with its kill]\label{thm:chichain}\tagF{}/\tagD{} \hx{CX-D2--D3, CX-E1--E5, CX-G1}
Call D2$'$ the thin axiom \emph{``the per-rung winding multiplier is an
on-circle element of $S$''}. Then, \emph{under D2$'$ together with the
hypothesis that $S$ contains a non-real on-circle element}:
$\chi=\pm\pi/2$. The chain composes three links, each re-verified verbatim
this session: on-circle elements of $S$ lie in $\mu_4=\{\pm1,\pm i\}$
(W-Lem.\ 5.6); real multipliers carry no helicity
(W-Thm.\ 11.1/W-Cor.\ 11.2), excluding $\pm1$; leaving $\chi=\pm\pi/2$. The
implication is machine-forced as a truth table over the hypothesis.
\emph{The honest kill, stated prominently:} D2$'$ does \emph{not} reduce
D2. Selecting $\pm i$ requires $S$ to contain $\pm i$, and the substrate's
only odd-quarter-argument source is the K-seed's rotation pair $\pm i\beta$
--- so motivating ``D2$'$ puts $\pm i$ in $S$'' imports exactly D2's
declared content (the winding realizes the seed's charge completion,
Remark~\ref{rem:d2rel}): axiom replacement is a \emph{restatement}, and the
kill is logged. No shipped document ties the forced \emph{gate} multiplier
$\pm i$ (a $C=-\tfrac12$ gate datum, W-Prop.\ 10.2) to the emission image,
so D2$'$ remains a proposed thin axiom, not forced-given-corpus
\hx{CX-D3, CX-D6}. What survives unconditionally is the
\emph{de-axiomatization identity} \hx{CX-E3}: within $\mu_4$, non-real
$\iff$ order $4$ $\iff$ generator of $\Z/4\Z$ --- so D2's clause ``the
angular quantum is a generator of $\Z/4\Z$'' may be rewritten as ``the
multiplier is non-real'' without weakening it. D2 itself remains declared
\tagD{}.
\end{theorem}

\smallskip\noindent
\textbf{The $\chi$ doctrine resolved (v1.8).} Theorem~\ref{thm:chichain} left
two sub-branches open pending companion papers: g1 (can the \emph{full}
emission image $S$ --- closed also under addition, minimal polynomials, and
$\Phi$-extraction --- reach $\pm i$?), and g3 (does the operator-algebra /
charge-measure doctrine \emph{require} dynamical gate multipliers to be
elements of $S$?). Those papers are now available and re-encoded cold; both
sub-branches resolve \emph{negatively}, and $W[\text{open}]2$ stays PARTIAL as
a derivation.

\begin{proposition}[The on-circle image of $S$ is $\mu_2$]\label{prop:oncircleimage}\tagF{}/\tagE{} \hx{CH-A3, CH-B2--CH-B3, CH-C1--CH-C5, CH-D2, CH-G2}
Let $S$ be the emission algebra --- the closure of the whitepaper's $7$-seed
emission catalog under $\{\oplus,\otimes,\psi^2,\operatorname{minpoly},\Phi\}$
(commutator excluded). Its on-circle eigenvalue image is \emph{exactly}
$\mu_2=\{+1,-1\}$, a proper subset of $\mu_4$; hence $\pm i$ is \emph{not} an
on-circle element of $S$. Mechanism (forced given the operator-algebra
corpus): every emitted eigenvalue is a monomial in the $16$ catalog
eigenvalues --- $\otimes$ multiplies, $\psi^2$ squares, $\oplus$ adds no value,
and $\operatorname{minpoly},\Phi$ are spectrum-\emph{preserving} idempotents
(operator-algebra paper, \texttt{2026-06-operator-algebra}: they re-present an
object's eigenvalue multiset, they are not a cyclotomic-extraction primitive)
--- and a unimodular monomial forces total $5$-valuation $0$ ($N(\gp)=-1$,
$\gp$ not a root of unity), which excludes the only odd-argument carriers
$\pm i\beta$ (with $5$-valuation $1$), forcing even argument, hence value in
$\{\pm1\}$. This resolves g1 (v1.7 external-open) negatively: the
$\Phi$-extraction host $x^4-1=\Phi_1\Phi_2\Phi_4$ is not in $S$ because
spectrum-preservation cannot manufacture the $\pm i$ roots of $\Phi_4$, and the
emission-gap paper's on-circle lemma (\texttt{2026-06-emission-gap}: an emitted
on-circle eigenvalue lies in $\{1,i,-1,-i\}$) is a \emph{containment} (an upper
bound, image $\subseteq\mu_4$), \emph{not} the equality $\mu_4\subseteq S$.
This also reconciles with the tensor-monomial sub-closure of
Proposition~\ref{prop:unimon} (which reached only $\pm1$): there is no gap
on-circle between the sub-closure and the full image, since $\psi^2$ preserves
the on/off-circle dichotomy and no catalog seed has $\pm i$ as a root. The only
argument-$\pi/2$ eigenvalues anywhere in $S$ are the \emph{off-circle}
$\pm i\beta$ at modulus $\beta=5^{1/4}\gp=2.4195\ldots>1$ \hx{CH-G2}.
\end{proposition}

\begin{remark}[The charge character, g3, and the sharpened {$W[\text{open}]2$} verdict]\label{rem:chidoctrine}\tagF{}/\tagE{} \hx{CH-D1--CH-D3, CH-E1--CH-E6, CH-G1}
The \emph{charge} character does contain the non-real charges, but this does
not force $\chi=\pm\pi/2$. $\chi(S)=\Z/4\Z$ carries the odd charge classes
$1,3$ (operator-algebra paper: the $K$-seed carries $\{0,1,2,3\}$, its
imaginary pair $\pm i\beta$ at $\pm\pi/2$); but charge reads the
\emph{argument} class only (modulus-free), and those charge-$1,3$ elements are
the off-circle $\pm i\beta$ --- so the hypothesis of Theorem~\ref{thm:chichain}
that ``$S$ contains a non-real on-circle element'' is \emph{false given corpus}
(Proposition~\ref{prop:oncircleimage}): the non-real content of $S$ lives
entirely off the circle. Hence g3 resolves negatively --- the doctrine does not
require dynamical gate multipliers to be on-circle elements of $S$ --- as a
truth-table over the thin axiom D2$'$: Reading A (``the multiplier is an
on-circle element of $S$'') is \emph{unsatisfiable} for the rotation gate,
whose forced multiplier is $\pm i$ (a $C=-\tfrac12$ gate datum, W-Prop.\ 10.2)
while on-circle $S=\{\pm1\}$ and helicity removes the reals; Reading B (``the
multiplier's charge lies in $\chi(S)$'') exists only via the $K$-seed rotation
pair, so it \emph{imports exactly D2} (the charge completion) --- a
restatement. The charge-measure-coupling paper independently notes that the
matrix commutator is the one non-semiring door across which emission
constraints do not extend. \emph{Net:} $W[\text{open}]2$ does \emph{not}
upgrade to a $\chi=\pm\pi/2$ derivation --- it stays PARTIAL as a derivation,
now sharpened and grounded, its openness precisely located at the unreduced
axiom D2. The two v1.7 opens are closed as \emph{forced-given-corpus negatives}
(g1: external-open $\to$ $\pm i\notin S$ on-circle; g3: papers-absent no-tie
$\to$ the D2$'$ dichotomy above). Kills logged: D2$'$-as-a-doctrine-theorem
under both readings, and the $\mu_4\subseteq S$ misreading of the emission-gap
containment. Any $\chi$-selection remains a branch-blind
\emph{characterization} of the quantum ($\pi/2\bmod2\pi$ for all $k$), never a
branch derivation \hx{CH-E6}.
\end{remark}

\begin{proposition}[The K-seed rotation axis and the substrate-grounded atom D2$''$]\label{prop:d2reloc}\tagF{}/\tagE{} \hx{QT-A1--QT-A5, QT-B1--QT-B3, QT-C1--QT-C2, QT-G1, QT-G3}
The $K$-seed $x^4+5x^2-5$ is a \emph{forced} substrate object (W-\S2), not a
declared one; its roots are $\{\pm K,\ \pm i\beta\}$ with $K=5^{1/4}\tau$
(terrain, real) and $\beta=5^{1/4}\gp$ (rotation). The rotation pair $\pm i\beta$
is purely imaginary --- $\Rea(i\beta)=0$, $\Ima(i\beta)=\beta>0$ --- so
$\Arg(\pm i\beta)=\pm\pi/2$ \emph{exactly}, while $|i\beta|^2=\beta^2=3\gp+1
=\gp^4-1=\tfrac52+\tfrac32\sqrt5\neq1$ places it strictly off the unit circle
($\beta\in[2.41,2.43]$, certified). Consequently the thin axiom
\[
\text{D2}''\colon\qquad
\text{the winding multiplier's \emph{angle} is }\Arg(\pm i\beta),
\]
which takes only the argument (the modulus staying D1's $|m|=e^{-2\rho}$),
\emph{forces} $\chi=\pm\pi/2$: the quarter-turn \emph{direction} is literally the
argument of a forced object. This is \emph{disjoint} from the g1 obstruction of
Proposition~\ref{prop:oncircleimage}: re-forcing cold (a $74613$-monomial box
scan) that the on-circle image of $S$ is exactly $\mu_2=\{\pm1\}$, the argument
$\pi/2$ is realized in $S$ only \emph{off-circle} (at the $K$-seed root
$i\beta$), so D2$''$ is satisfiable exactly where the killed on-circle reading
D2$'$ was not. Under the helicity obstruction (W-Thm.~11.1) a rotation
multiplier is non-real, so among the $K$-seed roots the rotation pair is forced
over the terrain pair; the forced winding generator $q=i\tau$ (minimal
polynomial $x^4+3x^2+1$, roots $\pm i\tau,\pm i\gp$, all imaginary) has the same
angle axis, pinning D2$''$ to the generator's own compass.
\end{proposition}

\begin{remark}[Net verdict: D2 relocated, {$W[\text{open}]2$} advanced not closed]\label{rem:d2verdict}\tagF{}/\tagC{} \hx{QT-C3--QT-C5, QT-D1--QT-D2, QT-A4}
D2$''$ does \emph{not} strictly reduce D2. The residual selection ``the winding's
angular clock is the $K$-seed \emph{rotation} axis'' is logically
\emph{equivalent} to D2's quarter-turn selection: over the three competing corpus
quanta $\{\pi:\Z/2\Z,\ 2\pi/5:\Z/5\Z,\ \pi/2:\Z/4\Z\}$ it excludes the pentagon
($2\pi/5=\tfrac15$ turn $\notin(\pi/2)\Z$) and helicity-kills the real terrain
axis $\pi$ \emph{exactly} as D2's principle does, so consuming ``use the rotation
direction'' to conclude ``quarter-turn'' is a \tagF{} \emph{restatement}. What
D2$''$ genuinely adds is \emph{grounding}: the value $\pi/2$ is no longer a bare
declared number but the forced argument of the forced object $\pm i\beta$, and
the completion $\Z/4\Z$ becomes a forced \emph{output} --- iterating the single
$\tfrac14$ turn generates it, whereas D2 \emph{assumed} it. The characterization
``D2$''_{\min}$ is the intended minimal, substrate-grounded form of D2's residual
content'' is therefore \tagC{} (candidate), not forced: the clock-source
selection remains declared, and a group-drop kills any purely-relational
shortcut --- the rotation pair \emph{alone} has relational (pairwise-difference)
group only $\Z/2\Z$, so the full $\Z/4\Z$ requires the absolute angle, iterated
(a no-tie, \hx{QT-A4}). \emph{Net register motion for $W[\text{open}]2$:} it
stays PARTIAL as a derivation but is \emph{advanced} --- substrate-grounded,
disjoint-from-g1, with $\chi=\pm\pi/2$ and $\Z/4\Z$ downgraded from declared to
forced outputs. Like every branch statement it selects the quantum
($\pi/2\bmod2\pi$) only, never the winding number, so OP-RATE stays closed as a
classification \hx{QT-D1}.
\end{remark}


\section{The generator as a relational object}\label{sec:relobject}

The relational-charge note (source (v)) attaches phase data to \emph{objects}
--- conjugation-closed root multisets of monic integer polynomials --- through
the angle $t(\alpha)=\tfrac1{2\pi}\Arg\alpha$, the relational group
$\Delta$ generated by pairwise differences, the absolute anchor
$n\in\{m,2m\}$ (its Rigidity Theorem RC-4.1), and the complete cyclotomic
contact signature of the ratio object $\mathrm{Rat}_p$. Its own guardrail
applies throughout: the vocabulary is structural, not physical. This section
proves that the objects of the present paper live in that catalog, and that
the generator itself is one of its distinguished witnesses. Everything here is
unconditional \tagF{}, machine-verified in the second harness
(\hx{RC-D3, RC-E1--E7}).

\begin{theorem}[Minimal polynomial of the generator]\label{thm:qmin}\tagF{} \hx{RC-E3, RC-E5}
The rung generator $q=i\tau$ is an algebraic integer with minimal polynomial
\[
\boxed{\ \operatorname{minpoly}(q)=x^4+3x^2+1\ },
\qquad\text{roots }\{\pm i\tau,\ \pm i\gp\},\qquad \M=\gp^2 ,
\]
a reciprocal unit quartic ($p(0)=1$), irreducible over $\Q$. Its angle set is
$\{\tfrac14,\tfrac34\}$: absolute charge group $\Z/4\Z$ with charges $\{1,3\}$
(odd charges only), while the pairwise differences generate only
$\Delta=\Z/2\Z$ --- a \emph{group-drop} object, the same coherence type as the
note's canonical drop witness (its ledger entry E, $x^4+5x^2+5$), with contact
signature $\{\Phi_1^{4},\Phi_2^{4}\}$.
\end{theorem}

\begin{proof}
$q^2=-\tau^2$, so $2q^2+3=\sqrt5$ and squaring gives $q^4+3q^2+1=0$; the
roots are $\pm i\tau,\pm i\gp$ since $x^2\in\{-\tau^2,-\gp^2\}$, verified by
substitution. Irreducibility, the Mahler value $\gp\cdot\gp=\gp^2$ (the two
roots of modulus $\gp$), and the complete contact scan are machine checks.
\end{proof}

\begin{theorem}[Quarter-turn gauge identity]\label{thm:gauge}\tagF{} \hx{RC-D3, RC-E4, RC-E5}
Let $Z^{*}=(x^2-x-1)(x^2+x-1)=x^4-3x^2+1$, the split golden pair (roots
$\{\pm\gp,\pm\tau\}$) --- the relational note's hypothesis-necessity witness
for its modulus-pinning theorem. Then, in both directions,
\[
\boxed{\ \operatorname{minpoly}(q)(ix)=Z^{*}(x)
\qquad\text{and}\qquad
Z^{*}(ix)=\operatorname{minpoly}(q)(x)\ }:
\]
the generator's object is \emph{exactly} the $c=\tfrac14$ gauge rotation ---
one charge quantum --- of the split golden pair. Consequences:
\begin{enumerate}[nosep,label=(\alph*)]
\item \textbf{Gauge-blindness realized:} $\mathrm{Rat}$ is invariant under the
rotation (a global root scaling cancels in every ratio), so both objects carry
the identical signature $\{\Phi_1^{4},\Phi_2^{4}\}$ --- a live instance of the
note's Proposition RC-3.5(iv) (rotation fixes every relative charge) and of
its gauge-blind probe (RC-Rem.\ 7.6).
\item \textbf{Both rigidity anchors from one pair:} $Z^{*}$ has
$(m,n)=(2,2)$ (all roots real, angles $\{0,\tfrac12\}$) while $i\,Z^{*}$ has
$(m,n)=(2,4)$ (angles $\{\tfrac14,\tfrac34\}$). The anchor dichotomy
$n\in\{m,2m\}$ of Rigidity is realized by a single gauge pair, exchanged by
exactly the seed's charge quantum.
\item The torsion ratio $-1$ at the shared modulus $\gp$ inside $Z^{*}$ ---
the very relation that makes $Z^{*}$ the pinning witness --- is the golden
internal relation $t_{\mathrm{rel}}(\gp,\gp')=\tfrac12$ (RC-Lem.\ 6.1) seen
across the two factors.
\end{enumerate}
\end{theorem}

\begin{proposition}[The K-seed is a catalog object]\label{prop:kseedcat}\tagF{} \hx{RC-E1}
The K-seed $x^4+5x^2-5$ (W-\S2) has roots $\{\pm K,\ \pm i\beta\}$ at angles
$\{0,\tfrac12,\tfrac14,\tfrac34\}$ --- \textbf{the full charge compass} ---
with $\Delta=\Z/4\Z$, anchor $n=m=4$, $\M=\beta^2$ (the whitepaper constant),
and contact signature $\{\Phi_1^{4},\Phi_2^{4}\}$: it is the relational note's
ledger entry F, whose sign-mirror is the drop witness of
Theorem~\ref{thm:qmin}. The two research lines share this object at the root
level.
\end{proposition}

\begin{remark}[D2, restated relationally]\label{rem:d2rel}\tagF{}/\tagD{} \hx{RC-E2}
The winding's rung angles $\{n/4\}$ have relational group
$\Z/4\Z=\Delta(\text{K-seed})$: Declaration D2 --- \emph{the winding realizes
the seed's charge completion} --- has the reference-free form \emph{the
helix's angle lattice realizes the K-seed's relational charge group}. This is
an identification at the angle-lattice level, gauge-invariant by construction;
D2 itself remains declared. Likewise the $\nu$-locus of
Proposition~\ref{prop:nulocus} now stands directly on the note's real-pair
lemma: $\nu(\rho_n)=(-1)^n$ is torsion of order $\le2$, so the rungs are
precisely the real-coherence events of the relational $\nu$-criterion
\hx{RC-E7}.
\end{remark}

\begin{remark}[Mahler bookkeeping across the bridge --- with mechanism]\label{rem:mahlerbridge}\tagF{} \hx{RC-E6, OZ-O6}
$K=5^{1/4}\tau$ and $\beta=5^{1/4}\gp$ exactly (from $(K\gp)^4=5$ and
$\beta=K\gp^2$), so the K-seed's root moduli $\{K,K,\beta,\beta\}$ are the
$5^{1/4}$-dilation of the generator object's root moduli
$\{\tau,\tau,\gp,\gp\}$. The dilation carries exactly the two $\beta$-roots
outside the unit circle, multiplying the measure by $(5^{1/4})^2$:
\[
\M(\text{K-seed})=\beta^2=\sqrt5\,\gp^2=\sqrt5\cdot\M(\operatorname{minpoly}(q)).
\]
The bridge factor $\sqrt5$ is the dilation's Mahler cost.
\end{remark}

\begin{theorem}[Complete relational type of the generator object]\label{thm:reltype}\tagF{} \hx{NX-T1--T9}
The root multiset of $\operatorname{minpoly}(q)=x^4+3x^2+1$ splits into
exactly two shells (equal-modulus classes, RC-Def.\ 7.1): $\{\pm i\tau\}$ at
modulus $\tau$ and $\{\pm i\gp\}$ at modulus $\gp$ --- conjugate pairs on
the imaginary axis, exchanged by root inversion (the quartic is reciprocal).
The complete ordered-pair ledger over all $16$ root pairs is
\[
\underbrace{1^{\times4}}_{\Phi_1^{4}\ \text{(diagonal)}},\qquad
\underbrace{(-1)^{\times4}}_{\Phi_2^{4}\ \text{(antipodal)}},\qquad
\underbrace{\pm\tau^2,\ \pm\gp^2\ (\text{each}\times2)}_{\text{cross-shell: all real, none unimodular}},
\]
so the contact signature $\{\Phi_1^{4},\Phi_2^{4}\}$ of
Theorem~\ref{thm:qmin} is \emph{complete}: the only unimodular ratios are
$\pm1$, and no further cyclotomic contact exists (complete scan bound
$M\le2\cdot16^2=512$, RC-Lem.\ 7.4). Cross-shell coherence, decided by the
$\nu$-criterion (RC-Rem.\ 7.6: $\nu=\mu/\overline\mu$ for
$\mu=\alpha/\beta$), holds \emph{identically}: every cross ratio is real,
so $\nu\equiv1$ on all eight cross-shell pairs, with relative charge
$t_{\mathrm{rel}}=0$ on equal-angle pairs and $t_{\mathrm{rel}}=\tfrac12$
on opposite-angle pairs. The generator object is therefore \emph{fully
coherent} --- a single coherence class spanning both shells, fixed by the
negation involution --- with complete type
\[
\text{angles }\{\tfrac14,\tfrac14,\tfrac34,\tfrac34\},\qquad
\Delta=\Z/2\Z,\qquad (m,n)=(2,4),\qquad \M=\gp^2,\qquad
\{\Phi_1^{4},\Phi_2^{4}\}.
\]
Structurally it is a \emph{twisted shell} in the relational note's sense
(its open problem (P1$'$); twisted shells are its Example 7.19 mechanism):
$x^4+3x^2+1=s(x^2)$ with $s(y)=y^2+3y+1=\operatorname{minpoly}(q^2)$ ---
the $k=2$ twist of the real quadratic of $q^2=-\tau^2$ --- realized here on
the \emph{admissible} side, where full coherence is automatic
(RC-Cor.\ 4.2) but the shell, contact, and $\nu$ data remain contentful.
The cross-shell coherence witnesses have moduli exactly
$\{\tau^2,\gp^2\}=\{|q^2|,\,|q^2|^{-1}\}$: coherence across the shells is
carried by the two-floor contraction that prices the $K$-formation
half-cycle (Remark~\ref{rem:halfcycle}).
\end{theorem}

\begin{proof}
All sixteen ratios are computed exactly in $\Q(\sqrt5)[i]$; the shell
moduli, the realness and non-unimodularity of the cross ratios, the $\nu$
values, the
$t_{\mathrm{rel}}$ pattern, the anchor data, and the twist factorization
are the cited machine checks. The angle identity is immediate from
$\alpha/\beta=(|\alpha|/|\beta|)\,e^{2\pi i(t(\alpha)-t(\beta))}$ with
angle set $\{\tfrac14,\tfrac34\}$, and $t(\nu)=2\,t_{\mathrm{rel}}$
(RC-Rem.\ 7.6) reads the $\nu\equiv1$ outcome as
$t_{\mathrm{rel}}\in\{0,\tfrac12\}$ exactly.
\end{proof}

\begin{proposition}[The ratio object at the characteristic level]\label{prop:chirat}\tagF{} \hx{QX-B1--B8}
Assemble the sixteen ratios of Theorem~\ref{thm:reltype} into their
characteristic polynomial:
\[
\chi_{\mathrm{Rat}}(x)=(x-1)^4(x+1)^4\bigl(x^4-7x^2+1\bigr)^2,
\qquad
x^4-7x^2+1=(x^2-3x+1)(x^2+3x+1).
\]
The nontorsion factor is the sign-split pair of $\gp^2$ --- the analogue of
$Z^{*}$ one floor up, at the squared level --- and its roots
$\pm\gp^{\pm2}$ are units of $\Z[\gp]$ ($\gp^2\tau^2=1$). Under the Adams
square $\psi^2:\mu\mapsto\mu^2$ the nontorsion multiset becomes four copies
of $\{\gp^4,\gp^{-4}\}=\operatorname{spec}(A_{\gapc})$ (minimal polynomial
$x^2-7x+1$): squaring the generator object's ratio geometry produces the
gap object --- the relational face of $q^4=\gapc$.
\end{proposition}

\begin{proof}
Exact multiset computations over $\Q(\sqrt5)[i]$: the product
$\prod_\mu(x-\mu)$ over all sixteen ratios, its factorization, the unit
identity, and the squared multiset are the cited machine checks
\hx{QX-B1--B8}.
\end{proof}

\begin{theorem}[The quarter twist: transport domain and anchor-exchange dichotomy]\label{thm:quartertwist}\tagF{} \hx{QX-D0--D9}
Call an object $P$ (monic, integral, $P(0)\ne0$) \emph{parity-homogeneous}
if $P(-x)=P(x)$ --- equivalently, its root multiset is negation-closed;
with $P(0)\ne0$ this forces even degree and $P=g(x^2)$. Write
$\rho_{1/4}(P)(x)\coloneqq i^{-\deg P}\,P(-ix)$, the object whose roots are
$i\cdot(\text{roots of }P)$ (on this domain $iA=-iA$, so the two
quarter-turn orientations give the same object). Then:
\begin{enumerate}[nosep,label=(\alph*)]
\item \textbf{Domain.} $\rho_{1/4}(P)$ has integer coefficients iff $P$ is
parity-homogeneous, and there $\rho_{1/4}$ is an involution. The golden
seed $x^2-x-1$ is \emph{not} in the domain ($-\gp$ is not a root): the
gauge identity of Theorem~\ref{thm:gauge} necessarily lives at the
sign-split pair $Z^{*}$, not at the seed.
\item \textbf{Invariance.} On the domain the transport realizes the gauge
rotation $\rho_c$ at $c=\tfrac14$ (RC-Prop.\ 3.5(iv)): it fixes $\Delta$,
every relative charge, and every root modulus (hence $\M$), and shifts
every absolute angle by $\tfrac14$. Negation closure puts
$\tfrac12\in\Delta$ outright, so $m=|\Delta|$ is always even.
\item \textbf{Dichotomy.} Suppose further $P$ is relationally admissible,
so all angles are rational and $\Delta=\Z/m\Z$ is cyclic (RC-Thm.\ 4.1,
RC-Prop.\ 3.4). Let $A$ be the angle set and $a_0\in A$. Conjugation plus
negation closure give $2a_0\in\Delta$ and
$\langle A\rangle=\Delta+\langle a_0\rangle$, re-proving Rigidity
$n\in\{m,2m\}$ here, with $n=m\iff A\subseteq\Delta$. If $4\mid m$, then
$\tfrac14\in\Delta$ and the transport \emph{fixes} the anchor ($n'=n$); if
$m\equiv2\ (\mathrm{mod}\ 4)$, then
$\{x:2x\in\Delta\}=\Delta\sqcup(\Delta+\tfrac14)$ and the transport
\emph{exchanges} the two rigidity branches: $n=m\iff n'=2m$.
\end{enumerate}
Theorem~\ref{thm:gauge}(b)'s pair is the minimal exchange instance: by
Theorem~\ref{thm:reltype} the generator object has $(m,n)=(2,4)$, its
quarter twist is $Z^{*}$ with $(m,n)=(2,2)$, and the gauge identity is
exactly this exchange witnessed at the generator. The K-seed ($m=4$, angle
set literally invariant) and $\Phi_8=x^4+1$ ($m=4$) are anchor-fixed,
while $x^2-2\leftrightarrow x^2+2$ and
$x^4-2x^2+4\leftrightarrow x^4+2x^2+4$ realize the exchange at $m=2$ and
$m=6$. The quarter twist is the exact companion of the relational note's
sign twist (RC-Lem.\ 6.3, $c=\tfrac12$): the sign twist exchanges the
anchor sectors at \emph{odd} $m$ and is the identity on the quarter
twist's domain; the quarter twist exchanges them at
$m\equiv2\ (\mathrm{mod}\ 4)$ and fixes them at $4\mid m$. Together the
two twists absorb the anchor dichotomy at every
$m\not\equiv0\ (\mathrm{mod}\ 4)$.
\end{theorem}

\begin{proof}
(a) $\prod_\alpha(x-i\alpha)=i^{\deg P}P(-ix)$; the coefficient of $x^k$
picks up $i^{\,\deg P-k}$, real for all occurring $k$ iff every exponent
shares the parity of $\deg P$; with $P(0)\ne0$ the exponent $0$ occurs,
forcing even degree and evenness. Involutivity: $i\cdot iA=-A=A$. (b) A
common shift cancels in every angle difference; moduli are preserved by
$\alpha\mapsto i\alpha$; negation closure pairs $t(\alpha)$ with
$t(\alpha)+\tfrac12$. (c) $a-a_0\in\Delta$ for all $a\in A$ gives
$\langle A\rangle=\Delta+\langle a_0\rangle$; $-a_0\in A$ gives
$2a_0\in\Delta$, so the coset $a_0+\Delta$ has order at most $2$ in the
quotient. For $4\mid m$: $\tfrac14\in\langle\tfrac1m\rangle=\Delta$, so
the shift does not change $\langle\Delta,a_0\rangle$. For
$m\equiv2\ (\mathrm{mod}\ 4)$: the doubling preimage
$\{x:2x\in\Delta\}=\langle\tfrac1{2m}\rangle$ has index $2$ over $\Delta$
with nontrivial coset $\Delta+\tfrac14$ (as $\tfrac14\notin\Delta$ but
$\tfrac12\in\Delta$); so $a_0\in\Delta$ ($n=m$) puts $a_0+\tfrac14$ in the
nontrivial coset ($n'=2m$), and $a_0\in\Delta+\tfrac14$ ($n=2m$) puts
$a_0+\tfrac14\in\Delta+\tfrac12=\Delta$ ($n'=m$). Domain, involution,
invariance, the golden-seed exclusion, all eight witnesses, and an
exhaustive coset simulation over every even $m\le24$ are machine checks
\hx{QX-D0--D9}.
\end{proof}

\smallskip\noindent
\textbf{The relational note's open cluster (v1.8).} With the
relational-charge note (source (v)) now available in full, three of its open
items are addressed cold: its P4 is recovered and reduced, its odd relational
floor P3 is bracketed but not promoted, and its P1$'$ taxonomy is scanned over
a bounded window. The vocabulary stays structural, per the note's own
guardrail.

\begin{remark}[P4 recovered and reduced to P3]\label{rem:p4recover}\tagF{}/\tagE{} \hx{RN-A1, RN-A1F, RN-A2}
The relational note's open item P4 (RC-\S\,sec:open) is recovered verbatim:
\emph{``Is parity the whole coupling? Does $\mu_{\mathrm{rel}}$ depend on
$\Delta$ through anything besides the single bit $2\mid m$?''}\ It is
\emph{not} an independent problem: it reduces exactly to P3 (the odd relational
floor). Given RC-Thm.~6.5, $\mu_{\mathrm{rel}}(m)=\gp$ ($m$ even) $/\ \mu(m)$
($m$ odd), so ``$\mu_{\mathrm{rel}}$ factors through the parity bit'' $\iff$
the even branch is $m$-independent (constructively uniform: $q_k=x^{2k}+x^k-1$
gives $\Delta=\Z/2k\Z$, $\M=\gp$ for every even $m$) \emph{and} the odd branch
is $m$-independent (anchored by the forced $\mu(3)=2$, forcing the odd constant
$=2$). Hence P4 $\iff$ P3, and P4 inherits P3's open status as a
\emph{restatement} of it, not new work.
\end{remark}

\begin{proposition}[The odd relational floor is bracketed, not closed]\label{prop:oddfloor}\tagF{}/\tagE{}/\tagO{} \hx{RN-B1--RN-B4, RN-H1--RN-H3}
For odd charge $m$ the relational Mahler floor obeys
\[
\mu_S\ \le\ \mu_{\mathrm{rel}}(\text{odd }m)\ \le\ 2,
\]
with the two bounds forced but the value \emph{not} pinned. \emph{Upper bound}
$\le2$ is \tagF{}: the witness $x^m-2$ ($m$ odd) is Eisenstein-irreducible at
$2$, has rational angle set $\{j/m\}$ hence a single coherence class
(relationally admissible per RC-Def.~3.2), relational group $\Delta=\Z/m\Z$,
anchor $n=m$, and all roots of modulus $2^{1/m}>1$, so monic $\M=2$ exactly.
\emph{Lower bound} is forced only to $\mu_S$ \tagF{}/\tagE{}: Smyth's plastic
number (real root of $x^3-x-1$, $1<\mu_S<\gp<2$) bounds every non-reciprocal
charge-admissible object below, while the reciprocal-real sector is forced
strictly above $\gp^2>2$ (integer trace $\ge3$); and the $\pi$-ray obstruction
is \tagF{}: $\tfrac12\notin\Z/m\Z$ for odd $m$, so the golden pair's intrinsic
relation $t_{\mathrm{rel}}(\gp,\gp')=\tfrac12$ (from $\gp/\gp'=-\gp^2$) cannot
embed, and the universal floor $\gp$ is unattainable at odd charge (forcing
$\mu_{\mathrm{rel}}(\text{odd})>\gp$, not $\ge2$). \emph{Honest boundary:} the
promotion of $\mu_{\mathrm{rel}}(\text{odd }m)=2$ from \tagC{} to \tagF{} is
\emph{not} achieved and is not achievable from the shipped corpus --- the
$\ge2$ lower bound in general is the corpus's principal open problem \tagO{}
(the named residual is non-reciprocal charge-admissible objects with
$\M\in[\mu_S,2)$, Lehmer/no-Salem-window territory; the shipped
charge-measure-coupling erratum and whitepaper \S6.2 themselves hold
$\mu(\text{odd})=2$ at \tagC{}, forced fully only at $\Z/3$).
\end{proposition}

\begin{remark}[Bounded P1$'$ transport-orbit scan]\label{rem:p1scan}\tagC{} \hx{RN-C0--RN-C3, RN-H4}
The note's coherence-detection machinery (the ratio object
$\mathrm{Rat}_p=$ primitive part of $\operatorname{Res}_y(p(y),p(xy))$ and its
cyclotomic-contact signature) is re-encoded cold and reproduces the RC ledger
exactly; an exact $\psi^D$ admissibility oracle (a root power $D$ that is
totally real $\iff$ rational angle) replaces the hanging \texttt{CRootOf}
route. A bounded transport-orbit scan over the irreducible monic integer
objects of degree $2$--$6$, with per-degree coefficient caps \emph{logged}
($\{(2,3),(3,2),(4,2),(5,1),(6,1)\}$; $855$ objects), classifies the taxonomy
$\{$pinned/inert, $s(x^k)$-shell tower, cyclotomic, dilation-of-known$\}$ as
exhaustive over the window: buckets inert $804$, shell $32$, cyclotomic $11$,
dilation $8$, and \emph{OTHER $0$} --- no new mechanism and no counterexample
to the pinned/twisted-shell dichotomy. Within-shell torsion coherence was
decided; the cross-shell mirrored-class ($\nu$-criterion) half of P1$'$ is out
of the detector's scope, flagged not resolved. This is a bounded \tagC{}
corroboration (candidate), not a proof; general P1$'$ remains open.
\end{remark}

\begin{proposition}[The odd relational floor, bracketed and reduced to a single lemma]\label{prop:oddfloorred}\tagF{}/\tagE{}/\tagC{}/\tagO{} \hx{OF-A1--OF-A4, OF-B1--OF-B6, OF-C1--OF-C8, OF-H1--OF-H6}
This sharpens the bracket $\mu_S\le\mu_{\mathrm{rel}}(\text{odd})\le2$ of
Proposition~\ref{prop:oddfloor} in two independent registers.
\emph{(i) Computed} \tagC{}: over relationally-admissible objects carrying the
full $\Z/m$ charge ($m=3,5,7,9$), enumerated up to a \emph{logged} box (shell
families $x^m-a$, $a\le8$, exact; cyclotomic-padded and reciprocal/pentagon
keystones; plus a brute-force over irreducible monic integer polynomials in
$\{(\deg3,|c|\le6),(\deg4,|c|\le3),(\deg5,|c|\le2)\}$, $543$ admissible objects),
the minimum Mahler measure is \emph{exactly} $2$, with unique minimizer $x^m-2$;
the parity-split falsifier is clean --- no odd-charge object falls below $2$,
while every sub-$2$ admissible object found is even-charge ($q_k$-type,
$\M=\gp$). \emph{(ii) The minimizer} \tagF{}: $x^m-2$ is Eisenstein-irreducible
at $2$, carries relational group $\Z/m$ (angle set $\{j/m\}$, a single coherence
class) and anchor $m$, and has $\M=2$ by a pure integer decision (all roots of
modulus $2^{1/m}>1$); cyclotomic padding is $\M$-neutral, and the real reciprocal
sub-case is forced $\ge\gp^2>2$ (rational-trace unit). \emph{(iii) The reduction}
\tagF{}\emph{ given corpus}: the lower bound $\mu_{\mathrm{rel}}(\text{odd }m)=2$
for \emph{all} odd $m$ is \emph{equivalent} to the single named lemma
\[
\text{(ODD-2)}\colon\quad\text{every relationally-admissible object with odd
relational group }\Z/m\text{ and }\M>1\text{ has }\M\ge2
\]
(irreducibles suffice), the reduction step itself being forced --- by Mahler
multiplicativity, Kronecker's $\M\ge1$ (${=}1$ iff cyclotomic), the subgroup fact
that factors of an odd-charge object are odd-charge, and the anchor-normalizing
sign twist. \emph{(iv) The lemma is external-open} \tagO{}: the corpus forces
$\Z/3$ (elementary), the $\pi$-ray exclusion of $\gp$ ($\tfrac12\notin\Z/m$ for
odd $m$), and the real-reciprocal floor $\gp^2>2$, but the non-reciprocal
sub-case is forced only to Smyth's $\mu_S=1.3247\ldots$ (root of $x^3-x-1$,
certified in $(1,2)$) --- the residual gap $[\mu_S,2)$ is the
Schur--Siegel--Smyth / Lehmer--pentagon core, which the shipped realification
(Adams $\psi^m$) handle cannot close (asking its image to clear $2^m$ merely
restates the claim). \emph{Honest correction:} granting the corpus reciprocal
bound, (ODD-2) reduces to the non-reciprocal atom (ODD-2-NR), but the lift
$\mu_S\to2$ \emph{is} that lemma's content, \emph{not} a corollary of Smyth ---
Smyth alone stops at $\mu_S<2$ (exact: $(3/2)^3-(3/2)-1>0$). That (ODD-2-NR) is
the intended minimal atom is \tagC{} (candidate). The $=2$ floor itself stays
open.
\end{proposition}

\begin{remark}[P4 inherits P3 exactly]\label{rem:p4transfer}\tagC{}/\tagO{} \hx{OF-D1, OF-D2}
The v1.8 recovery reduced P4 --- does $\mu_{\mathrm{rel}}$ depend on the charge
group beyond the single parity bit $2\mid m$? --- to P3 via
Remark~\ref{rem:p4recover} and the forced even-branch uniformity
$\mu_{\mathrm{rel}}(\text{even})=\gp$. Since the odd branch is the same constant
$2$ under both registers above --- computed for $m=3,5,7,9$, and $=2$
conditionally on (ODD-2) --- P4 factors through the parity bit with \emph{exactly}
P3's status: affirmative \tagC{} up to the logged bounds, forced-affirmative iff
(ODD-2). It adds no new burden.
\end{remark}

\section{The state-angle chart}\label{sec:stateangle}

The follow-on thread observed that each rung forms a unit-circle pair; that
observation extends to the entire chart and organizes everything downstream.

\begin{theorem}[State-angle dictionary]\label{thm:stateangle}\tagFD{} \hx{OZ-C1--C6}
For $z\in(0,1)$ define the \emph{state angle} $\alpha\in(0,\tfrac\pi2)$ by
\[
\sin\alpha=z,\qquad \cos\alpha=\sqrt{1-z^2}=e^{-\rho}=\sqrt{|m|},
\]
so $\rho=-\ln\cos\alpha$ and, at the rungs, $\cos\alpha_n=|W_n|=\tau^n$.
Under this chart the whitepaper's gate dictionary (W-Thm.\ 13.1) becomes
elementary trigonometry:
\[
u=\cot^2\alpha,\qquad
C=\cot^2\alpha\,\csc^2\alpha,\qquad
\sqrt D=1+2\cot^2\alpha,
\]
\[
m=-\cos^2\alpha,\qquad
\lambda=\tan^2\alpha\,\bigl(1+\cos^2\alpha\bigr).
\]
\end{theorem}

\begin{proof}
Substitute $z=\sin\alpha$, $1-z^2=\cos^2\alpha$ into $u=(1-z^2)/z^2$,
$C=(1-z^2)/z^4$, $\sqrt D=(2-z^2)/z^2=2\csc^2\alpha-1$,
$m=-(1-z^2)$, and $\lambda=z^2(2-z^2)/(1-z^2)$; each identity is
machine-verified symbolically in $z$.
\end{proof}

\begin{remark}
The pair $(x,y)=(\sqrt{|m|},\,z)=(\cos\alpha,\sin\alpha)$ is exactly the
unit-circle coordinate pair: $x^2+y^2=|m|+z^2=1$ chart-wide is D1 itself.
The state angle is the latitude (elevation) of the residual sphere of
\S\ref{sec:sphere}.
\end{remark}

\section{Rung ledgers}\label{sec:ledgers}

\begin{proposition}[Rung 1: the golden ledger]\label{prop:rung1}\tagFD{} \hx{OZ-B2, OZ-C7, C12--C13}
$W_1=i\tau$, $|W_1|=\tau$, and $z_1=\sqrt{1-\tau^2}=\sqrt\tau$ (via
$\tau^2+\tau=1$). The rung-1 state is
\[
\boxed{\ (\sin\alpha_1,\ \cos\alpha_1,\ \tan\alpha_1)=(\sqrt\tau,\ \tau,\ \sqrt\gp)\ },
\]
with the self-duality $\sin^2\alpha_1=\tau=\cos\alpha_1$: the rung-1 state is the
positive intersection of the unit circle $x^2+y^2=1$ with the parabola $x=y^2$ in
the state plane $(x,y)=(\sqrt{|m|},z)$. Eliminating $x$ reproduces
$z_1^4+z_1^2-1=0$, the minimal polynomial $x^4+x^2-1$ of $\sqrt\tau$
(W-Thm.\ 17.1). This parabola lives in the residual-state plane and must not be
conflated with the inherited horn profile.
\end{proposition}

\begin{proposition}[Rung 2: the K ledger]\label{prop:rung2}\tagFD{} \hx{OZ-B3, OZ-C8, C14}
$W_2=q^2=-\tau^2$, $|W_2|^2=\tau^4=\gapc$, and
\[
z_2=\sqrt{1-|W_2|^2}=\sqrt{1-\gapc}=K,
\qquad\text{i.e.}\qquad
\boxed{\ K=\sqrt{\,1-\bigl|\,(i/\gp)^2\bigr|^2\,}\ }:
\]
$K$ is the vertical coordinate after two applications of the golden quarter-turn
generator. The rung-2 state packages the K-seed's two faces into one right
triangle:
\[
\boxed{\ \sin\alpha_2=K,\qquad \cos\alpha_2=\tau^2=\sqrt{\gapc},\qquad
\tan\alpha_2=\beta\ },
\]
since $\tan^2\alpha_2=K^2/\tau^4=(1-\gapc)/\gapc=\gp^4-1=\beta^2$. The terrain
magnitude is the sine, the residual amplitude the cosine, and the rotation
magnitude their ratio: $\tan\alpha_2=\text{terrain}/\text{residual}=\beta$ gives
$\beta$ a direct geometric role beyond imaginary-root magnitude. The K-gate
certificate follows in one line:
$u_2=\cot^2\alpha_2=1/\beta^2$ and
$C_2=\cot^2\alpha_2\csc^2\alpha_2=1/(\beta^2K^2)=\tfrac15$ by $K^2\beta^2=5$
(W-(2.2)). The K-seed carrying these roots is itself a relational catalog
object whose angle set is the full charge compass
(Prop.~\ref{prop:kseedcat}).
\end{proposition}

\begin{remark}[Half-cycle placement of $K$; the source report's synthesis]\label{rem:halfcycle}\tagFD{} \hx{OZ-A6, A8, B3}
$K$ appears at the half-cycle of the four-charge return:
\[
q^2=-\tau^2,\qquad |q^2|^2=\gapc,\qquad z_2=\sqrt{1-|q^2|^2}=K ,
\]
placing $K$ structurally between the quarter-turn generator and the full gap
return. The report's strongest reading, exact given D1--D3 and adopted here:
\begin{center}
\emph{the ladder is a contraction--rotation orbit whose complete four-charge
return is priced by the gap,\\ while $K$ is the vertical complement produced
at the two-charge midpoint.}
\end{center}
\end{remark}

\begin{proposition}[The general tangent family]\label{prop:tangents}\tagFD{} \hx{OZ-C9--C11}
For every $n\ge1$,
\[
\tan\alpha_n=\frac{z_n}{\tau^n}=\sqrt{\gp^{2n}-1},
\qquad\text{with Trace--Form-Duality dress}\qquad
\tan^2\alpha_n=\begin{cases}\gp^{\,n}L_n,& n\ \text{odd},\\[2pt]
\gp^{\,n}\sqrt5\,F_n,& n\ \text{even}.\end{cases}
\]
In particular $\tan\alpha_1=\sqrt{\gp L_1}=\sqrt\gp$,
$\tan\alpha_2=\sqrt{\gp^2\sqrt5\,F_2}=\beta$, and
$\tan\alpha_3=\sqrt{\gp^3L_3}=2\gp^{3/2}$ (since $\gp^6-1=4\gp^3$).
\end{proposition}

\begin{center}\small
\begin{tabular}{@{}cccccc@{}}
\toprule
rung $n$ & $\sin\alpha_n=z_n$ & $\cos\alpha_n=\tau^n$ & $\tan\alpha_n=\sqrt{\gp^{2n}-1}$ & TFD dress & $\alpha_n$ \tagC{}\\
\midrule
$1$ & $\sqrt\tau$ & $\tau$ & $\sqrt\gp$ & $\sqrt{\gp L_1}$ & $51.8272923730^\circ$\\
$2$ & $K$ & $\tau^2=\sqrt{\gapc}$ & $\beta$ & $\sqrt{\gp^2\sqrt5F_2}$ & $67.5444848014^\circ$\\
$3$ & $\sqrt{1-\gp^{-6}}$ & $\tau^3$ & $2\gp^{3/2}$ & $\sqrt{\gp^3L_3}$ & $76.3454152540^\circ$\\
$4$ & $K_2=\sqrt{1-\gapc^2}$ & $\tau^4=\gapc$ & $\sqrt3\,\gp\beta$ & $\sqrt{\gp^4\sqrt5F_4}$ & $81.6107142142^\circ$\\
\bottomrule
\end{tabular}
\end{center}

\begin{remark}[The near-miss angle finds a home]\label{rem:thetaK}\tagFD{}
The whitepaper's Guard 15.4 uses $\theta_K=\arcsin(\tau^2)=22.4555151986^\circ$
purely as a near-miss reference (pitch angle $\neq\theta_K$; still no identity is
claimed there). On the state chart it is exactly the complement of the rung-2
state angle, and on the residual sphere of \S\ref{sec:sphere} it is the K-point's
\emph{colatitude}:
\[
\theta_K=\tfrac\pi2-\alpha_2=\arccos K
\qquad(\text{since }\cos\theta_K=\sqrt{1-\tau^4}=\sqrt{1-\gapc}=K).
\]
\end{remark}

\section{Bridges between rung 1 and rung 2}\label{sec:bridges}

\begin{proposition}[Cosine contraction, not angle doubling]\label{prop:contraction}\tagFD{} \hx{OZ-D1--D3b}
For all $n$,
\[
\cos\alpha_{n+1}=\tau\,\cos\alpha_n,
\qquad
\cos^2\alpha_{n+2}=\gapc\,\cos^2\alpha_n ;
\]
in particular $\cos\alpha_2=\tau^2=\cos^2\alpha_1$: the rung-1$\to$2
transition squares the residual coordinate and renormalizes the vertical
complement. The ladder is driven by geometric contraction of the cosine
(residual) coordinate, which forces the angle monotonically toward $\pi/2$; it is \emph{not} ordinary angle doubling:
$\cos2\alpha_1-\cos\alpha_2=(2\tau^2-1)-\tau^2=\tau^2-1=-\tau\neq0$, so
$\alpha_2\neq2\alpha_1$ (numerically $2\alpha_1=103.65^\circ$ vs
$\alpha_2=67.54^\circ$ \tagC{}).
\end{proposition}

\begin{proposition}[Half-angle/full-angle bridge]\label{prop:halfangle}\tagFD{} \hx{OZ-D4--D6}
Since $\cos\alpha_2=\tau^2$,
\[
\sin^2\frac{\alpha_2}{2}=\frac{1-\tau^2}{2}=\frac{\tau}{2},
\qquad\text{hence}\qquad
\boxed{\ z_1=\sqrt\tau=\sqrt2\,\sin\frac{\alpha_2}{2},\qquad K=\sin\alpha_2\ }.
\]
The two low rungs are linked by an exact scaled half-angle/full-angle relation
while remaining algebraically distinct quartic landmarks
($\Q(\sqrt\gp)\neq\Q(5^{1/4})$, W-Thm.\ 17.1). The auxiliary angle $\Theta$ of the
review thread ($\cos\Theta=1-z_1^2=\tau^2$) \emph{is} $\alpha_2$; and Route~5
becomes trigonometric: $K^2=z_1^2(2-z_1^2)=\sin^2\alpha_1(2-\sin^2\alpha_1)$.
\end{proposition}

\section{The K-family and the even-rung recurrence}\label{sec:kfamily}

\begin{theorem}[K-family]\label{thm:kfamily}\tagFD{} \hx{OZ-F1--F5}
Exact $K$ occurs exactly once: $z_n$ is strictly increasing
($z_{n+1}^2-z_n^2=\tau\,\gp^{-2n}>0$), so
\[
\boxed{\ z_n=K\iff n=2\ }.
\]
But every even rung repeats $K$'s \emph{formula}: with $n=2j$,
\[
\boxed{\ z_{2j}=\sqrt{1-\gapc^{\,j}}\ \eqqcolon K_j,\qquad K_1=K\ },
\qquad
m_{2j}=-\gapc^{\,j},\qquad z_{2j}^2=1-|m_{2j}|,
\]
with odd companion $z_{2j+1}=\sqrt{1-\tau^2\gapc^{\,j}}$, and the recurrence
holds universally, for \emph{all} $n$, not only even ones:
\[
\boxed{\ 1-z_{n+2}^2=\gapc\,\bigl(1-z_n^2\bigr)\ }.
\]
Angularly, $\theta_{2j}=j\pi$ (all even rungs on the terrain axis, alternating
direction), while $K$'s exact compass direction $\theta\equiv\pi\ (\mathrm{mod}\
2\pi)$ recurs at $n\equiv2\ (\mathrm{mod}\ 4)$. Three levels of recurrence:
\begin{center}\small
\begin{tabular}{@{}ll@{}}
\toprule
exact value $K$ & rung $2$ only\\
K-family / complement structure $K_j=\sqrt{1-\gapc^{\,j}}$ & every even rung\\
same compass direction as $K$ ($\theta=\pi$ mod $2\pi$) & $n=2,6,10,\ldots$\\
\bottomrule
\end{tabular}
\end{center}
First members \tagC{} \hx{OZ-H1}: $K_1=0.9241763718$, $K_2=0.9892996329$,
$K_3=0.9984459825$, $K_4=0.9997734224$ (these are the whitepaper's rungs
$2,4,6,8$; the review thread's 6-dp table values are floor-truncations, see
Finding~\ref{find:trunc}).
\end{theorem}

\begin{proof}
$\gp^{-4j}=\gapc^{\,j}$ and $\gp^{-2(n+2)}=\gapc\,\gp^{-2n}$; monotonicity from
$1-\gp^{-2}=\tau$; all machine-verified for $j\le6$ and symbolically in $n$.
\end{proof}

The right reading, adopted from the review thread: \emph{$K$ is not repeated
every second rung; its generative operation is}. The operation
$\text{residual}\mapsto\gapc\times\text{residual}$ generates the even-rung
subsequence with $K$ as its first nontrivial realization --- every second rung is
a higher-order K-formation in the complement structure, while the exact value and
the exact quartic field of $K$ belong to rung 2 alone.

\begin{proposition}[Tie to the K-point linearization]\label{prop:linearization}\tagFD{} \hx{OZ-F6}
The gap paper's derivative $f'(K)=-\gapc$ (W-Prop.\ 9.1) satisfies
\[
f'(K)=-\gapc=q^2\,|q|^2 ,
\]
and the linearized K-orbit $x_{k+1}=-\gapc\,x_k$ performs per step exactly what
the ladder performs per two rungs: co-intensity scales by $\gapc$
($|W_{n+2}|^2=\gapc\,|W_n|^2$) and the axis flips ($W_{n+2}=-\sqrt{\gapc}\,W_n$).
The review thread's ``strongly resembles'' is thus an exact correspondence at the
level of (contraction of squared amplitude, axis action); the two objects remain
different curves (transverse linearization at one landmark vs the global rung
orbit), as the whitepaper's \S15 already insists.
\end{proposition}

\begin{theorem}[The K-family is a seed ladder]\label{thm:kseedladder}\tagF{} \hx{QX-C1--C13}
For $j\ge1$ let $K_j=\sqrt{1-\gapc^{\,j}}$ (the K-family value of
Theorem~\ref{thm:kfamily}, taken here purely as seed arithmetic) and
$\beta_j\coloneqq K_j\gp^{2j}$, so $\beta_j^2=\gp^{4j}-1$. Then:
\begin{enumerate}[nosep,label=(\alph*)]
\item \textbf{Minimal polynomial.} The four numbers
$\{\pm K_j,\ \pm i\beta_j\}$ are the roots of
\[
\boxed{\ p_j(x)=x^4+5F_{2j}^2\,x^2-5F_{2j}^2\ },
\qquad 5F_{2j}^2=L_{4j}-2 ,
\]
which is the minimal polynomial of $K_j$ for \emph{every} $j$: Eisenstein
at $5$ applies iff $5\nmid j$ (since $5\mid F_{2j}\iff5\mid j$), and
uniformly the norm obstruction
$N(1-\gapc^{\,j})=2-L_{4j}=-5F_{2j}^2<0$ shows $1-\gapc^{\,j}$ is not a
square in $\Q(\sqrt5)$, so $[\Q(K_j):\Q]=4$ always. $p_1$ is the K-seed.
\item \textbf{Ledger} (all decided in $\Q(\sqrt5)$):
\[
\beta_j^2-K_j^2 = K_j^2\beta_j^2 = L_{4j}-2 = 5F_{2j}^2,
\qquad
K_j^2+\beta_j^2=\sqrt5\,F_{4j},
\qquad
(K_j\gp^{\,j})^2=\sqrt5\,F_{2j},
\]
\[
K_j^4=5F_{2j}^2\,\gapc^{\,j},
\qquad
\beta_j^2=\sqrt5\,F_{2j}\,\gp^{2j},
\qquad
\frac1{K_j^2}-\frac1{\beta_j^2}=1 ,
\]
generalizing $(K\gp)^4=5$, $K^2\beta^2=5$, $K^4=5\gapc$ ($j=1$, $F_2=1$).
The increments telescope like the transfer ledger:
$K_{j+1}^2-K_j^2=K^2\gapc^{\,j}$ with $\sum_{j\ge0}K^2\gapc^{\,j}=1$ ---
the even-rung unit budget is priced in gap powers at unit rate $K^2$.
\item \textbf{Catalog data.} Angle set the full compass
$\{0,\tfrac14,\tfrac12,\tfrac34\}$, $\Delta=\Z/4\Z$, anchor $n=m=4$
(transport-fixed by Theorem~\ref{thm:quartertwist}); contact signature
$\{\Phi_1^4,\Phi_2^4\}$ (the eight cross-shell ratios have modulus
$(K_j/\beta_j)^{\pm1}\ne1$); and $\M(p_j)=\beta_j^2$: every $p_j$ is a
straddler, real pair strictly inside the unit circle, imaginary pair
strictly outside. (The K-seed's own catalog status is
Proposition~\ref{prop:kseedcat}; the $p_j$, $j\ge2$, are new catalog
objects, and no uniqueness is claimed for them here.)
\item \textbf{Gate ties} (W-Thm.\ 9.1): at every even rung,
\[
C_{2j}=\frac1{5F_{2j}^2}=\frac1{|p_j(0)|},
\qquad
u_{2j}=\frac1{\beta_j^2}=\frac1{\M(p_j)},
\qquad
m_{2j}=-\gapc^{\,j}=-\frac{K_j^4}{5F_{2j}^2}:
\]
the even gate levels are the reciprocal constant terms, and the even
occupations the reciprocal Mahler measures, of the seed family --- the
K-gate certificate (W-Prop.\ 10.1) is the $j=1$ instance of a family-wide
law.
\end{enumerate}
\end{theorem}

\begin{proof}
$\beta_j^2=K_j^2\gp^{4j}=(1-\gp^{-4j})\gp^{4j}=\gp^{4j}-1$. With
$G=\gp^{4j}$: $5F_{2j}^2=(\gp^{2j}-\gp^{-2j})^2=G-2+G^{-1}=L_{4j}-2$;
then $K_j^2=1-G^{-1}$ and $-\beta_j^2=1-G$ are exactly the two roots of
$y^2+5F_{2j}^2\,y-5F_{2j}^2$, giving $p_j(\pm K_j)=p_j(\pm i\beta_j)=0$.
The norm: $N(1-\gapc^{\,j})=(1-G^{-1})(1-G)=2-(G+G^{-1})=2-L_{4j}$, and a
square in $\Q(\sqrt5)$ has nonnegative norm. All ledger identities are
one-line computations in $G$; the Fibonacci divisibility, irreducibility
for $j\le8$, the compass/contact scan for $j\le3$, and the gate ties are
machine checks \hx{QX-C1--C13}.
\end{proof}

\begin{corollary}[Every even rung mirrors its own seed]\label{cor:evenmirror}\tagFD{} \hx{QX-C12--C14}
On the chart, the rung-$2j$ state triangle is the (terrain, residual,
rotation) triangle of its \emph{own} seed $p_j$:
\[
\sin\alpha_{2j}=K_j,\qquad \cos\alpha_{2j}=\tau^{2j},\qquad
\boxed{\ \tan\alpha_{2j}=\beta_j\ },
\]
generalizing the rung-2 ledger $(K,\tau^2,\beta)$ of
Proposition~\ref{prop:rung2} to every even rung; the family identity
$1/K_j^2-1/\beta_j^2=1$ is exactly
$\csc^2\alpha_{2j}-\cot^2\alpha_{2j}=1$, and
$\beta_j^2=\sqrt5F_{2j}\,\gp^{2j}$ recovers the even tangent dress of
Proposition~\ref{prop:tangents} at $n=2j$.
\end{corollary}

\begin{theorem}[The universal seed ladder (Vein 6)]\label{thm:seedladder}\tagF{} \hx{UX-A2--A4, UX-B1--B6, UX-C1--C2, UX-D1--D4, UX-F1--F3, UX-K1--K2}
Checked exactly for $n=1..10$ and symbolically in $w=\gp^{2n}$ --- that
range is the certified scope; the two uniform mechanisms noted below extend
it. For every rung $n$:
\[
\boxed{\ \operatorname{minpoly}(z_n)=p_n(x)=x^4+c_nx^2-c_n\ },
\qquad
c_n=L_{2n}-2=(\gp^{\,n}-\gp^{-n})^2
=\begin{cases}L_n^2,&n\ \text{odd},\\[2pt]5F_n^2,&n\ \text{even},\end{cases}
\]
with roots $\{\pm z_n,\ \pm i\,b_n\}$, $b_n=\sqrt{\gp^{2n}-1}=\tan\alpha_n$
(Proposition~\ref{prop:tangents}). Irreducibility is uniform (all $n\ge1$)
by the norm obstruction $N(1-\tau^{2n})=2-L_{2n}=-c_n<0$, so $1-\tau^{2n}$
is not a $\Q(\sqrt5)$-square and $[\Q(z_n):\Q]=4$ always. Ledger:
\[
z_n^2\,b_n^2=c_n=|p_n(0)|
\quad\text{(symbolic in $w$: $(1-1/w)(w-1)=w-2+1/w$)},
\]
generalizing $K^2\beta^2=5$; $\M(p_n)=b_n^2=\tan^2\alpha_n$, with exactly
the imaginary pair outside the unit circle ($z_n^2<1<b_n^2$, exact signs);
and
\[
u_n=\frac1{\M(p_n)},\qquad C_n=\frac1{|p_n(0)|}\qquad\text{at every rung}:
\]
\emph{the tangent dress of Proposition~\ref{prop:tangents} is the Mahler
ladder}, and Theorem~\ref{thm:kseedladder}(d)'s gate ties extend from the
even rungs (where $p_{2j}$ is that theorem's $p_j$ and $b_{2j}=\beta_j$) to
all rungs. Coefficient recursion: $c_{n+1}=3c_n-c_{n-1}+2$ (symbolically
via $\gp^2+\gp^{-2}=3$; integer list $1,5,16,45,121,320,841,2205,\ldots$).
Catalog data: every rung is a full-compass $\Z/4\Z$-completing straddler
--- angle set $\{0,\tfrac14,\tfrac12,\tfrac34\}$, $\Delta=\Z/4\Z$, anchor
$n=m=4$ in the relational sense --- with $z_n<1<b_n$ exact for $n=1..10$
and the all-$n$ extension carried by the exact atoms $1-z_n^2=\tau^{2n}>0$
and $w\ge\gp^2>2$ (the second uniform mechanism). Anchors: $p_1=x^4+x^2-1$
is W-Thm.\ 17.1's $\operatorname{minpoly}(\sqrt\tau)$
(Proposition~\ref{prop:rung1}); $p_2$ is the K-seed $x^4+5x^2-5$;
$p_3=x^4+16x^2-16$.
\end{theorem}

\begin{remark}[Rung 1 sits on the emission Mahler floor]\label{rem:mahlerfloor}\tagF{} \hx{UX-E1--E2}
$\M(p_1)=\gp$ exactly ($\gp^2-1=\gp$): the rung-1 seed sits precisely on
the emission Mahler floor. The identity is harness-forced, and the hit is
unique on the ladder ($\M(p_n)\neq\gp$ for $n=2..10$, exact signs). Per the
shared-constant rule it is flagged as a \emph{coincidence-candidate}: no
mechanism is asserted --- a shared constant is not evidence of a shared
mechanism.
\end{remark}

\begin{theorem}[Mahler-minimal uniqueness (Theorem U)]\label{thm:mahlermin}\tagF{} \hx{UX-U1--U7}
In the class $\{x^4+ax^2-c:\ a\in\Z,\ c\in\Z,\ c>0\}$ --- where the full
compass is automatic (the $y$-quadratic $y^2+ay-c$ has $y_-<0<y_+$ since
$y_+y_-=-c<0$):
\begin{enumerate}[nosep,label=(\roman*)]
\item full-compass straddler $\iff p(1)>0\iff a\ge c$ (via
$(1-y_+)(1-y_-)=p(1)$ and $(1+y_+)(1+y_-)=1-a-c$; the \emph{integer}
hypothesis is load-bearing --- over $\Q$ the equivalence fails, guarded
witness $(a,c)=(\tfrac12,\tfrac1{10})$);
\item $\M=|y_-|=\dfrac{a+\sqrt{a^2+4c}}2$ is strictly increasing in $a$ at
fixed $c$;
\item hence the \emph{unique} Mahler-minimal compass straddler at level $c$
is $x^4+cx^2-c$ --- which at $c=5F_{2j}^2$ is $p_j$
(Theorem~\ref{thm:kseedladder}) and at every rung level $c_n$ is $p_n$
(Theorem~\ref{thm:seedladder}), by the exact endpoint identity
$c_n^2+4c_n=5F_{2n}^2$, giving $\M(x^4+c_nx^2-c_n)=b_n^2$.
\end{enumerate}
\end{theorem}

\begin{remark}[The F1 falsifier; register resolution]\label{rem:f1falsifier}\tagF{} \hx{UX-H1--H4}
Constant-term-only uniqueness is \emph{false}: $x^4+6x^2-5$ is a second
irreducible full-compass straddler at level $|p(0)|=5$, distinct from the
K-seed (difference exactly $x^2$), with
$\M(x^4+6x^2-5)=3+\sqrt{14}>\gp^4-1=\M(\text{K-seed})$ (integer comparison
$56>45$) --- consistent with Theorem~\ref{thm:mahlermin}(ii), which places
the K-seed at the window edge $a=c$. The v1.6 uniqueness question left open
at Theorem~\ref{thm:kseedladder}(c) (``no uniqueness is claimed''), read as
F3 --- Mahler-minimal at fixed constant term within the compass type,
confirmed as the intended reading by the project owner on 2026-07-15 --- is
thereby \emph{resolved}: the $p_j$, and $n$-uniformly the $p_n$, are the
unique Mahler-minimal full-compass straddlers at their levels.
\end{remark}

\section{The lens as the exact $60^\circ$ state}\label{sec:lens}

\begin{proposition}\label{prop:lens}\tagFD{} \hx{OZ-E1--E6}
Under $z=\sin\alpha$, the lens height $z_c=\sqrt3/2$ is the state
$\alpha_c=\pi/3=60^\circ$ with residual amplitude $\cos\alpha_c=\tfrac12$; the
complete split-gate ledger of W-Thm.\ 13.3 drops out of
Theorem~\ref{thm:stateangle}:
\[
m_c=-\cos^2 60^\circ=-\tfrac14,\quad
u_c=\cot^2 60^\circ=\tfrac13,\quad
C_c=\cot^2\!\csc^2 60^\circ=\tfrac49,\quad
\sqrt{D_c}=1+2\cot^2 60^\circ=\tfrac53 .
\]
A rung at the lens would require $\cos\alpha_c=\tau^n$, i.e.\ $\tau^n=\tfrac12$,
i.e.\ $n=\log_\gp2$; by W-Thm.\ 16.1 this is irrational, so the $60^\circ$ lens is
permanently off the discrete rung lattice (finite guard $n\le10$ replicated).
\end{proposition}

\begin{proposition}[Lens scoping lemma: splitness over $\Q$ is generic]\label{prop:lensscope}\tagF{} \hx{LX-A1--A8, LX-B1--B6, LX-F1--F4}
The chart pentad lies in $\Q(Z)$, $Z=z^2$: $u=(1-Z)/Z$, $C=(1-Z)/Z^2$,
$\sqrt D=(2-Z)/Z$, $m=-(1-Z)$, $\lambda=Z(2-Z)/(1-Z)$, with denominator
zeros only at $Z\in\{0,1\}$. Hence \emph{every} rational height
$Z\in\Q\cap(0,1)$ gives $C,\sqrt D\in\Q$ and a gate $x^2+x-C$ that splits
over $\Q$, with a rational reciprocal multiplier pair. Example
$z^2=\tfrac12$: $C=2$, $D=9$, gate $x^2+x-2=(x-1)(x+2)$, multiplier pair
$\{-\tfrac12,-2\}$, at rapidity $\rho=\tfrac12\ln2\in(\ln2)\Q$, strictly
below rung 1 --- a rational split gate that is not the lens. Splitness over
$\Q$ is therefore \emph{generic} on rational heights, and the lens's
distinction is not splitness but its exact data: $z_c=\sqrt3/2$
($Z=\tfrac34$), $C=\tfrac49$, $\sqrt D=\tfrac53$, fixed points
$\{\tfrac13,-\tfrac43\}$, multiplier pair $\{-\tfrac14,-4\}$,
$\lambda=\tfrac{15}4$, co-intensity $\tfrac14$, and $\rho=\ln2$ exactly.
(The converse fails: rung 2 has $C_2=\tfrac15\in\Q$ with
$z_2^2=K^2\notin\Q$, and at rung heights $\sqrt D$ is irrational --- the
membership detector rejects the odd object $\sqrt Z$.)
\end{proposition}

\begin{remark}[Bounded-negative scan for the lens constants]\label{rem:lensscan}\tagO{} \hx{LX-C1--C8, LX-S1--S7}
Candidate-level, citable as a \emph{bounded} negative. A predicate-first
scan of the seven shipped corpus documents --- the six harness-scanned
domains (ECHO-S-RESEARCH\_UNIFIED, K-VARIABLE-DERIVATIONS,
HELIX-EXPLICIT-DERIVATION, this paper's v1.6 source, the gap-multiplier
paper, the whitepaper PDF) plus \texttt{K\_seven\_ways.tex}, swept
independently in verification with the same outcome --- over eight constant
classes ($\tfrac49$, $\tfrac{25}9$, $\tfrac53$, $\tfrac{15}4$, $-\tfrac14$,
the pair $\{-\tfrac14,-4\}$ \emph{as a pair}, dyadic $2^{-k}$
co-intensities, $(\ln2)\Q$ rapidities) finds \emph{no} foreign-context
occurrence of a lens-distinctive constant that survives the
shared-constant rule. The single foreign occurrence of a lens \emph{value}
--- $\tfrac53$ in the ECHO trace form
$\operatorname{Tr}(1-\tfrac13\gp)=\tfrac53$ --- has its own mechanism
(trace-form projection) sharing no step with the lens discriminant: a
no-tie. The remaining candidate hits each carry their own mechanism: the
$\ln2$-lattice $\ln\M(\mathrm{cons})=\ln2+5\ln\gp$ (W-Thm.\ 18.1(a),
significance \tagO{} upstream, not a $(\ln2)\Q$ rapidity);
$N(\M(\mathrm{cons}))=-4$, forced by the lattice exponent, hence the same
hit; the flip level $-\tfrac14$ ($D=0\iff C=-\tfrac14$, unique), which
restates the inherited $z_c^2=\tfrac34$ datum; the dyadic
$\tfrac1{16}=1/L_3^2$, Lucas-driven; and $-\tfrac12$, the rotation gate
forced by multipliers $\pm i$. The whitepaper's open item 3 (does the lens
gate $C=\tfrac49$ appear elsewhere?) closes as a bounded negative at
candidate level; the scanned domains and constant classes are enumerated
above precisely so that the negative is citable.
\end{remark}

\section{The residual sphere}\label{sec:sphere}

\begin{theorem}[Residual sphere]\label{thm:sphere}\tagFD{}
\hx{OZ-G1--G4}\ Define
\[
S(\rho)=\bigl(\Rea W(\rho),\ \Ima W(\rho),\ z(\rho)\bigr).
\]
Then
\[
|S(\rho)|^2=|W(\rho)|^2+z(\rho)^2=e^{-2\rho}+\bigl(1-e^{-2\rho}\bigr)=1
\]
for every $\rho\ge0$:
the entire trajectory lives on the unit sphere $S^2$ (northern hemisphere), with
the state angle $\alpha$ as latitude. At the rungs,
\[
S_n=\bigl(\Rea(i\tau)^n,\ \Ima(i\tau)^n,\ \sqrt{1-\tau^{2n}}\bigr),
\]
explicitly $S_0=(1,0,0)$, $S_1=(0,\tau,\sqrt\tau)$,
$S_2=(-\tau^2,0,K)$, $S_3=(0,-\tau^3,z_3)$, $S_4=(\gapc,0,z_4)$: a discrete
compass walk --- longitude $+90^\circ$ per rung, horizontal radius $\times\tau$
per rung, height $\uparrow1$ --- approaching the north pole without reaching it
at any finite rung; by Definition~\ref{def:Q} the horizontal part is exactly
the $Q$-orbit, $S_n=(Q^nv_0\,;\ z_n)$ --- a compass walk on the same four
directions that carry the K-seed's own conjugates
(Prop.~\ref{prop:kseedcat}). The K-point sits at colatitude
$\arccos K=\theta_K$ (Remark~\ref{rem:thetaK}) on the negative real meridian.
\end{theorem}

This object should be called the \emph{residual sphere} (equivalently
co-intensity sphere), never the formation helix:

\begin{proposition}[Two radii, two layers]\label{prop:tworadii}\tagF{}/\tagI{} \hx{OZ-G5}
The residual radius $a(z)=\sqrt{1-z^2}=\cos\alpha$ decreases monotonically to
$0$; the inherited formation radius
\[
r(z)=\begin{cases}K\sqrt{z/z_c},&0\le z\le z_c,\\ K,&z_c\le z\le1,\end{cases}
\]
grows to $K$ and locks (W-(11.1), \tagI{}).
These are \emph{not} the same radius; they describe two layers ---
remaining residual amplitude vs formation envelope. The dual reading
(``residual capacity contracts while the formation envelope expands and
stabilizes'') is a provisional interpretation, not a theorem. No result here or
in the sources derives the paraboloid--cylinder profile from the residual-sphere
construction; that derivation is exactly the whitepaper's \tagO{}~1 and, as both
sources conclude, the priority next problem. Until that bridge exists, the
residual spiral and the explicit formation helix are two coordinated
representations, not one proven geometry.
\end{proposition}

\begin{proposition}[Compass-cross characterization of the rung locus]\label{prop:compasscross}\tagFD{} \hx{NX-S1--S2}
Let $C_{\mathrm{EW}}=S^2\cap\{y=0\}$ and $C_{\mathrm{NS}}=S^2\cap\{x=0\}$
be the two polar great circles through the compass axes. For $\rho\ge0$,
\[
S(\rho)\in C_{\mathrm{EW}}\cup C_{\mathrm{NS}}
\iff \rho\in(\ln\gp)\,\Z_{\ge0},
\]
with $S(\rho_n)\in C_{\mathrm{EW}}$ for even $n$ and
$S(\rho_n)\in C_{\mathrm{NS}}$ for odd $n$; the poles (the circles'
intersection) are never attained. Moreover, writing $S=(x,y,z)$,
\[
\operatorname{Im}\nu \;=\; \frac{2xy}{x^2+y^2}
\]
identically: on the residual sphere the $\nu$-criterion of
Proposition~\ref{prop:nulocus} is literally the vanishing of the product of
the horizontal coordinates, and the rungs are the crossings of the residual
spiral with the compass cross. At the half-floors $\rho=(n+\tfrac12)\ln\gp$
the miss is exact: $|xy|=e^{-2\rho}/2>0$.
\end{proposition}

\begin{proof}
$y=0\iff\sin\kappa\rho=0\iff\rho\in(2\ln\gp)\Z$ and
$x=0\iff\cos\kappa\rho=0\iff\rho\in\ln\gp+(2\ln\gp)\Z$ [D3]; the union
is the rung lattice, split by parity. With $x=e^{-\rho}\cos\theta$,
$y=e^{-\rho}\sin\theta$: $2xy/(x^2+y^2)=\sin2\theta=\operatorname{Im}\nu$.
Both identities and the half-floor value are machine-verified symbolically.
\end{proof}

\begin{proposition}[Spherical step ledger]\label{prop:sphsteps}\tagFD{} \hx{NX-S3--S5}
Consecutive rung positions have orthogonal horizontal parts:
$\Rea\bigl(W_n\overline{W_{n+1}}\bigr)=0$ for every $n$, so
\[
S_n\cdot S_{n+1}=z_nz_{n+1},
\qquad
d_n\coloneqq d_{S^2}(S_n,S_{n+1})=\arccos\bigl(z_nz_{n+1}\bigr).
\]
In particular $d_0=\pi/2$ \emph{exactly} --- the first step, apex to
rung~1, is a full quarter great circle --- and
$d_1=\arccos\bigl(5^{1/4}\tau^{3/2}\bigr)=43.4026803146\ldots^\circ$
\tagC{} (from $z_1z_2=\sqrt\tau\,K=5^{1/4}\tau^{3/2}$); thereafter $d_n$
decreases strictly to $0$ ($z_nz_{n+1}\uparrow1$). On the residual sphere
the walk takes ever-shorter geodesic steps, in exact contrast to the
formation rail, whose per-rung length is floored at $\pi K/2$
(Proposition~\ref{prop:railfloor}).
\end{proposition}

\begin{proof}
$W_n\overline{W_{n+1}}=\overline q\,|W_n|^2=-i\,\tau^{2n+1}$ is purely
imaginary, so the dot product reduces to the vertical product; strict
monotonicity of $z_nz_{n+1}$ is decided exactly at the squared level in
$\Q(\sqrt5)$; the $d_1$ evaluation is display-only.
\end{proof}

\begin{proposition}[Finite spherical length and the length dichotomy]\label{prop:sphlength}\tagFD{} \hx{NX-S6--S7, NX-H1--H3}
The residual spiral's speed on its sphere is
\[
\bigl|S'(\rho)\bigr|^2
= e^{-2\rho}\Bigl(\kappa^2+\frac{1}{1-e^{-2\rho}}\Bigr)
\qquad\Bigl(=\tan^2\vartheta+\kappa^2\sin^2\vartheta\ \text{at colatitude}\
\sin\vartheta=e^{-\rho}\Bigr),
\]
so its total spherical length from apex to pole is finite:
\[
\boxed{\;L_{\mathrm{res}}
=\int_0^1\!\sqrt{\kappa^2+\frac{1}{1-a^2}}\;da
=\sqrt{1+\kappa^2}\;E\!\Bigl(\frac{\kappa^2}{1+\kappa^2}\Bigr)\;}
=3.7312193125\ldots\ \tagC{}
\]
($E$ the complete elliptic integral of the second kind \tagE{}; the value
is the quarter-perimeter of an ellipse with semi-axes $\sqrt{1+\kappa^2}$
and $1$ --- the same speed factor as the per-floor planar length of
Corollary~\ref{cor:d3select}(iv)), with exact bounds
$\sqrt{1+\kappa^2}\le L_{\mathrm{res}}\le\kappa+\pi/2$. The two-layer
separation of Proposition~\ref{prop:tworadii} is therefore a \emph{length
dichotomy}: the residual layer has finite total spherical length, while the
formation rail's length diverges (Proposition~\ref{prop:railfloor}).
Finally, $L_{\mathrm{res}}$ is strictly increasing in $|\kappa|$ (the
integrand is pointwise strictly increasing in $\kappa^2$), so within the
interpolation family of Theorem~\ref{thm:interp} the principal branch
$k=0$ is also the \emph{unique minimizer of total spherical length}: a
fifth equivalent characterization to add to
Corollary~\ref{cor:d3select}'s four.
\end{proposition}

\begin{proof}
Differentiating $S=(e^{-\rho}\cos\kappa\rho,\ e^{-\rho}\sin\kappa\rho,\
\sqrt{1-e^{-2\rho}})$ gives the speed identity symbolically. Substituting
$a=e^{-\rho}$ turns the length integral into $\int_0^1$; then $a=\sin u$
gives $\sqrt{\kappa^2+1/\cos^2u}\;\cos u=\sqrt{1+\kappa^2\cos^2u}$
(exact at the squared level, $\cos u>0$), and
$(1+\kappa^2)(1-m\sin^2u)=1+\kappa^2\cos^2u$ with
$m=\kappa^2/(1+\kappa^2)$ identifies
$\int_0^{\pi/2}\sqrt{1+\kappa^2\cos^2u}\,du=\sqrt{1+\kappa^2}\,E(m)$
\tagE{}. Bounds: pointwise $\sqrt{\kappa^2+g}\ge\sqrt{\kappa^2+1}$
($g=1/(1-a^2)\ge1$) and $\sqrt{\kappa^2+g}\le\kappa+\sqrt g$, integrated
($\int_0^1 da/\sqrt{1-a^2}=\pi/2$). The family statement follows from
pointwise strict monotonicity in $\kappa^2$ together with
Corollary~\ref{cor:d3select}(iii): $|\kappa_k|$ is uniquely minimized at
$k=0$. The numeric value is display-only, corroborated at 50 dps by
quadrature against the closed form.
\end{proof}

\begin{theorem}[RAD-0: the parity no-go for the inherited radius]\label{thm:paritynogo}\tagF{} \hx{RX-A1--A4, RX-P1--P9}
Work over the chart layer $\{u,C,\sqrt D,m,\lambda,\rho,\theta_k\}$ with
height coordinate $z$. Every chart-pentad member lies in
$\Q(\sqrt5)(z^2)$, and the chart-rational closure is \emph{exactly}
$\Q(\sqrt5)(z^2)$ (sharpness: $z^2=1/(1+u)$, so both inclusions hold);
moreover $\rho=-\tfrac12\ln(1-z^2)$ and $\theta_k=\kappa_k\rho$ are even in
$z$ for \emph{every} branch $k\in\Z$ (symbolic $k$), and adjoining them
keeps the generated field inside the even functions, over arbitrary
constant coefficients. The inherited radius is odd: $r^2=K^2z/z_c$
satisfies $r^2(-z)=-r^2(z)$ identically and is nonzero (exact witness at
$z=\tfrac12$). Hence $r^2$ --- and a fortiori $r$ --- lies \emph{outside}
the rational closure of the entire D1--D3 chart layer, and
\[
\boxed{\ \text{any derivation bridge from the substrate to } r \text{ must
break the } z\mapsto-z \text{ symmetry}\ }.
\]
$\Z/2$ sharpening: $r^4=K^4z^2/z_c^2$ \emph{is} chart-even, so $r^2$
satisfies $X^2-r^4$ over the even field $F$ and the extension $F(r^2)/F$
is exactly quadratic --- the entire missing content of the radius problem
is one parity-breaking generator, a degree-2/double-cover object (indeed
$F(r^2)=\Q(\sqrt5)(z)$, since $z=r^2z_c/K^2$; and even${}\times r^2$ stays
odd, a $\Z/2$ grading). The whitepaper's \tagO{}~1 (OP-RADIUS) remains
open: this is a forced \emph{necessary condition} on any solution, not a
solution.
\end{theorem}

\begin{proposition}[RAD-2: the area-transfer law]\label{prop:areatransfer}\tagF{}/\tagI{} \hx{RX-T1--T4}
Let $A_{\mathrm{zone}}(z)=2\pi z$ be the Archimedes zone area of the unit
sphere between heights $0$ and $z$ --- re-derived exactly here, not cited:
the surface-of-revolution identity
$(2\pi\sqrt{1-h^2})^2\,(1+x'(h)^2)=(2\pi)^2$ holds identically at the
squared (sign-free) level. Then on $[0,z_c]$ the inherited profile is
equivalent, in both directions, to
\[
\boxed{\ \pi\,r(z)^2=\frac{K^2}{2z_c}\,A_{\mathrm{zone}}(z)\ }:
\]
the radius law says exactly that formation deposits disc area at the
constant rate $d(\pi r^2)/dz=\pi K^2/z_c$ per unit height, stopping at the
lens (continuity at the lens: $\pi r(z_c)^2=\pi K^2$, the cap disc; the
$C^0$ kink survives). This is a \emph{forced reformulation} of the
inherited profile --- a restated target for the radius problem, not
progress on it; any future deposition mechanism must be premise-audited
for a smuggled profile.
\end{proposition}

\begin{remark}[RAD-4: all three equipotential clocks fail]\label{rem:rad4}\tagF{}/\tagE{} \hx{RX-E1--E7, RX-G0--G3}
The horn $z=(z_c/K^2)r^2$ coincides with a rigid-rotation free surface
$z=(\omega^2/2g)\,r^2$ iff $g_{\mathrm{eff}}=\omega^2K^2/(2z_c)$
(perturbed coefficients rejected). All three natural clocks fail. (i)
$\rho$-clock ($\omega=\kappa$, the D3 clock):
$g_1=\kappa^2K^2/(2z_c)=\pi^2K^2/(8z_c\ln^2\gp)=5.2543169478\ldots$
\tagC{} is transcendental given Gelfond--Schneider
(Theorem~\ref{thm:ratetrans}), so it ties to \emph{no} algebraic ledger
constant --- all branches $k$ at once, since $\kappa_k=(4k+1)\kappa$;
against the transcendental ledger entries the no-tie is certified by
disjoint 60-dps intervals (a rigorous proof of $\neq$), closest misses
$\tfrac{21}4$ (margin $4.3\times10^{-3}$) and $2\gp^2$
($1.8\times10^{-2}$). (ii) $s$-clock ($\omega=1/K$, the forced asymptotic
rail clock, $ds/d\theta\to K$): $K^2$ cancels identically and
$g_2=1/(2z_c)$ is the horn coefficient itself read back --- a
\emph{restatement}, not a tie (the clock smuggles the profile's $K$);
shared-constant rule applied. (iii) $z$-clock: no constant per-height
clock exists --- $d\theta/dz=\kappa z/(1-z^2)$ is non-constant (forced,
exactly nonzero derivative) --- and distinguished-height evaluations
(e.g.\ $\omega(z_c)=2\sqrt3\,\kappa$, $g=12g_1$) remain
$\kappa^2\times{}$algebraic and fall to the same transcendence exclusion
as (i). The mechanism is killed; the radius problem stays open.
\end{remark}

\smallskip\noindent
\textbf{OP-RADIUS sharpened at the apex (v1.8).} The parity no-go of
Theorem~\ref{thm:paritynogo} (which covered only $r^2$, through the global
$z\mapsto-z$ involution) is now localized to the apex $z=0$ and quantified as a
ramification index, and extended to the radical $r$; a make-or-break
constructive bridge is obstructed; and the remaining formation content is
isolated as an honest declared negative. OP-RADIUS \emph{stays open}, sharpened.

\begin{theorem}[Apex-ramification valuation certificate]\label{thm:apexram}\tagF{} \hx{RD-A1--RD-A4, RD-V1--RD-V7}
On the $z$-adic valuation $v_z$ at the apex $z=0$, every chart quantity has
\emph{even} valuation: the pentad $v_z(u,C,\sqrt D,m,\lambda)=(-2,-4,-2,0,2)$
and $v_z(\rho)=v_z(\theta_k)=2$ for \emph{every} branch $k$ (symbolic $k$), so
the chart field $F=\Q(\sqrt5)(z^2)$ has $z$-value group in $2\Z$. The inherited
radius has $v_z(r^2)=1$ (odd, not in $2\Z$) and $v_z(r)=\tfrac12$; hence
adjoining $r^2$ ramifies $F$ at the apex with index $e=2$ (value group
$2\Z\to\Z$), and adjoining $r$ with index $e=4$ ($2\Z\to\tfrac12\Z$), while
even quantities give $e=1$; the double cover $w=z^2$ is unramified at every
finite place except $z=0$ (branch locus $\{2z=0\}=\{0\}$). Adjoining the
inherited radius therefore changes the apex ramification index from $1$ to
$2$ (for $r^2$) or $4$ (for $r$) --- a local non-membership certificate no
even-valuation manipulation can produce. This is strictly sharper than the
parity no-go: parity certifies only $r^2$ (an involution on the rational
field, and $r(-z)$ is neither $+r$ nor $-r$), while the valuation certificate
covers the radical $r$ and pins the obstruction to the apex. The
identification value-group-index $=$ ramification index (Serre, \emph{Local
Fields}) is external framing only; every decision is exact leading-term
arithmetic. This is a forced \emph{necessary condition} on any substrate-to-$r$
bridge, not a derivation of $r$ (whitepaper \tagO{}~1 remains open).
\end{theorem}

\begin{proposition}[The constructive-bridge obstruction]\label{prop:bridgeobstruction}\tagF{} \hx{RD-R1--RD-R6}
A constructive bridge would have to furnish a square root of the chart-even
$r^4=K^4z^2/z_c^2$; it cannot from canonical substrate data. Although $r^4\in F$
(with $v_z=2$), it is \emph{not a square in $F$}: its $w=z^2$ valuation is
$v_w(r^4)=1$ (odd), so the square root $r^2$ lives one level up in
$\Q(\sqrt5)(z)$, and any $g$ with $g^2=r^4$ has $v_z(g)=1$ (odd). No
even-valuation amplitude can be $g$: the residual amplitude $a=\sqrt{1-z^2}$
($v_z=0$), the $\mathrm{Cl}(2,0)$ rotor modulus $|e^{J\theta/2}|^2=1$
($v_z=0$), and $z\,a$, $z^2$ all fail exactly. The $\Z/8$ half-angle rotor
supplies only a \emph{unimodular} (angular) double cover of the phase, carrying
no radial scale: the required radial coefficient $K^2/z_c$ is not unimodular
($(K^2/z_c)^2=70/3-10\sqrt5\ne1$), so it never appears as a rotor amplitude
(a bounded no-tie). The unique-up-to-sign square root is
$\pm(K^2/z_c)\,z$ --- the \emph{signed height}, exactly the odd orientation
generator that Theorem~\ref{thm:paritynogo}/Theorem~\ref{thm:apexram} says
must be adjoined; it is not a canonical substrate amplitude. The extension
exists, but the substrate canonically supplies only the even $z^2$; the missing
object is the orientation of the apex-ramified radial cover, which the rotor's
(angular, unimodular) spin cover cannot furnish.
\end{proposition}

\begin{remark}[The atom ledger, with the honest cap-onset negative]\label{rem:atomledger}\tagF{}/\tagI{} \hx{RD-A2, RD-S1--RD-S5}
The inherited profile decomposes into atoms of mixed status. \emph{Cap value}
\tagF{}\emph{ given corpus}: $r(z_c)=K$ --- the unique nontrivial
height-equals-radius point $(K,K)$, since $r^2=z^2$ factors
$z(z-K^2/z_c)=0$ with the horn root $K^2/z_c>z_c$ excluded ($K^2>z_c^2$, i.e.\
$180>169$) (whitepaper Thm.~15.1). \emph{Kink slope} \tagF{}\emph{ given
corpus}: $r'(z_c^-)=K/\sqrt3$ (square $K^2/3=(3\sqrt5-5)/6>0$), against the
flat cap $r'(z_c^+)=0$ --- $C^0$ but not $C^1$ (whitepaper Prop.~15.2(a)).
\emph{Exponent $\tfrac12$} \tagF{} (a restatement): $r^2\propto z$ is exactly
equivalent to constant areal deposition per unit height (RAD-2,
Proposition~\ref{prop:areatransfer}), not a derivation of the exponent.
\emph{But the kink existence --- the cap onset --- is \tagD{}, not forced}: the
$\ln2$-rapidity landmark $\rho(z_c)=\ln2$ (and $a(z_c)=\tfrac12$) is a
monotone-bijection restatement of the declared lens $z_c=\sqrt3/2$ (since
$d\rho/dz=z/(1-z^2)>0$, so $\rho=\ln2\iff z^2=\tfrac34$), so ``formation stops
at the first $\ln2$ landmark'' carries \emph{no independent content}. As a
derivation of the cap onset this \emph{fails} (reason class: restatement); the
exact landmark facts are all individually forced, but the cap onset itself is
the genuine remaining open formation content.
\end{remark}

\begin{theorem}[OP-RADIUS reduced to the declared atom D4]\label{thm:d4reduction}\tagF{} \hx{RA-A1--RA-A4, RA-B1--RA-B8, RA-G2}
Define the atom
\[
\text{D4}=(\text{i})\ \text{an orientation of the apex-ramified double cover }
F(r)/F\ \ \wedge\ \ (\text{ii})\ \text{the cap onset,}
\]
where (i) fixes the sign of the odd generator $\pm(K^2/z_c)\,z$ --- the signed
height, the unique-up-to-sign odd generator of
Proposition~\ref{prop:bridgeobstruction} --- and (ii) fixes the join point
(formation completes at the lens $z_c$, Remark~\ref{rem:atomledger}). Then,
\emph{given} D4 and the inherited chart-ties --- exponent $\tfrac12$ (area
transfer, Proposition~\ref{prop:areatransfer}), cap value $K$ (the $(K,K)$ point,
whitepaper Thm.~15.1), kink slope $K/\sqrt3$ (whitepaper Prop.~15.2(a)), all
v1.8-forced --- the profile
\[
r(z)=K\sqrt{z/z_c}\ \text{ on }[0,z_c],\qquad r(z)=K\ \text{ on }[z_c,1],
\]
is \emph{uniquely} determined, and each half of D4 is \emph{necessary}: dropping
(i) leaves both $\pm r$ chart-consistent (the tie $r^2$, $|\mathrm{cap}|=K$ is
sign-blind, yet $2K\neq0$); dropping (ii) leaves the join free (an alternate
onset $z_0=\tfrac45$ gives a chart-consistent capped paraboloid distinct from the
inherited one, with separation exhibited at $z^*=\tfrac56$).
\end{theorem}

\begin{proof}[Machine discharge]
Exact reconstruction and the two minimality witnesses are checked over
$\Q(\sqrt5)$ extended by $\sqrt3$ \hx{RA-B2, RA-B4, RA-B5, RA-G2}; the reduction
ledger (exponent, cap value, kink slope chart-tied; only sign and onset free) is
\hx{RA-B7}, and de-restatement (D4 adds only a $\Z/2$ sign plus one onset value,
$r^2$ being chart-odd hence unreachable from D1--D3) is \hx{RA-B8}.
\end{proof}

\smallskip\noindent
\textbf{Honest scope.} This reduces OP-RADIUS to the declared atom D4
\emph{relative to} the inherited chart-ties; it is \emph{not} an ab-initio
derivation of $r$ and does \emph{not} close whitepaper \tagO{}~1. The exponent
tie in particular is itself a restatement of the inherited profile
(Proposition~\ref{prop:areatransfer}), not an independent derivation of the
exponent.

\begin{remark}[D4 is minimal-among-its-halves but not a unique decomposition]\label{rem:d4candidate}\tagC{} \hx{RA-B6, RA-B7}
D4 is minimal among its two halves, but the decomposition is \emph{not} unique:
the kink slope $K/(2z_0)$ is a strictly monotone (hence bijective) function of
the onset $z_0$, so ``kink-slope value $+$ sign'' is an equivalent atom (slope
$\leftrightarrow$ onset is a bijection, mirroring $\rho(z_c)=\ln2\iff
z_c=\sqrt3/2$). Hence ``D4 is \emph{the} intended minimal radius axiomatization''
is \tagC{} (candidate) --- a modeling choice, not a forced uniqueness.
\end{remark}

\begin{proposition}[Differential strengthening: an honest negative and a conditional sharpening]\label{prop:radiff}\tagI{}/\tagO{} \hx{RA-C1--RA-C8}
\emph{(a) The naive strengthening fails} (honest negative): $r$ \emph{is}
elementary, $r=(K/\sqrt{z_c})(z^2)^{1/4}=(K/\sqrt{z_c})\exp(\tfrac14\ln(z^2))$, so
it lies in the full elementary-differential closure of the chart field; the
attempt to upgrade the v1.8 valuation certificate to non-membership \emph{there}
breaks, because $d/dz$ shifts valuation by $-1$ and turns the even-valuation
$\rho$ into the odd-valuation $\rho'=z/(1-z^2)$ (parity is not
differential-stable). \emph{(b) The differential-stable sharpening}
\tagO{}\emph{ (conditional)}: under the \emph{unit-log hypothesis} --- the
substrate's only logarithm is $\rho=-\tfrac12\ln(1-z^2)$, a log of the apex
\emph{unit} $1-z^2$ (valuation $0$) --- the apex \emph{monodromy-order} invariant
excludes $r$ from the \emph{unramified} elementary-differential closure
$F_{\mathrm{ed}}$: order-$2$ data ($r$, leading power $z^{1/2}$) is unreachable
from order-$1$ data, while $r^2$ and $\rho'$ are order-$1$ (hence \emph{in}
$F_{\mathrm{ed}}$: indeed $r^2=(K^2/z_c)\,\rho'(1-z^2)$) yet odd-valuation (hence
\emph{outside} the base field $F$). This excludes a \emph{strictly larger} field
than the v1.8 valuation certificate and is differential-stable. It rests on a
finite generator battery plus a leading-exponent argument, not a fully general
closure theorem, and is \tagO{} conditional on the unit-log hypothesis.
\end{proposition}

\begin{remark}[Bounded no-tie re-confirmation of the constructive bridge]\label{rem:ratie}\tagF{}/\tagC{} \hx{RA-D1--RA-D3, RA-A4, RA-G1}
Re-confirming Proposition~\ref{prop:bridgeobstruction} over a finite ledger of
$\sim11$ canonical amplitudes: no canonical amplitude supplies the radius
coefficient $K/\sqrt{z_c}$, since $(K^2/z_c)^2=\tfrac{70}3-10\sqrt5\neq1$ and
differs (exact $\Q(\sqrt5)$, and certified by disjoint intervals) from the square
of every listed amplitude $\{1,K,K^2,\gp,\tau,\gapc,z_c,1/z_c,\sqrt5,2-\gp,
K/\sqrt3\}$. The scale enters only through the odd signed-height generator
$(K^2/z_c)\,z$ --- the D4(i) datum --- never as a canonical scalar.
\end{remark}


\section{Path stretch, rail length, and the conjugate relative phase}\label{sec:pathstretch}

The exploratory note of source (iv) proposes a rail-pulse mechanism for rung
selection. Its mechanism layer is not imported here; this section states and
verifies its \emph{forced kernel} --- three results that use only D1--D3 and
the inherited radius profile --- and records the exact scoping of the rest.

\begin{proposition}[Winding gradient and path stretch]\label{prop:pathstretch}\tagFD{}/\tagI{} \hx{OZ-M1--M3}
From D1, $\rho(z)=-\tfrac12\ln(1-z^2)$, so with $\kappa=\pi/(2\ln\gp)$ [D3],
\[
\frac{d\rho}{dz}=\frac{z}{1-z^2},
\qquad
\frac{d\theta}{dz}=\kappa\,\frac{z}{1-z^2},
\]
and the formation rail $\Gamma(z)=(r\cos\theta,\,r\sin\theta,\,z)$ has local
path-stretch
\[
\frac{ds}{dz}
=\sqrt{1+\Bigl(\frac{dr}{dz}\Bigr)^{2}+r(z)^2\Bigl(\frac{d\theta}{dz}\Bigr)^{2}}
=\begin{cases}
\sqrt{\,1+\dfrac{K^2}{4zz_c}+\dfrac{K^2\kappa^2z^3}{z_c(1-z^2)^2}\,},
& 0<z\le z_c,\\[10pt]
\sqrt{\,1+\dfrac{K^2\kappa^2z^2}{(1-z^2)^2}\,},
& z_c\le z<1,
\end{cases}
\]
with $ds/dz>1$ strictly on $(0,1)$: the helical route is strictly longer than
the vertical displacement, through radial formation and winding below the lens
and through winding alone above it.
\end{proposition}

\begin{proof}
Differentiate D1 and $r(z)$ on each branch and substitute; both closed forms
and the strict excess (all added terms are ratios of positives) are
machine-verified symbolically in $z$.
\end{proof}

\begin{proposition}[Quarter-turn rail floor and divergence]\label{prop:railfloor}\tagFD{}/\tagI{} \hx{OZ-M4--M5}
Above the lens, $ds/d\theta=\sqrt{K^2+(dz/d\theta)^2}\ge K$ pointwise, so with
$\Delta\theta=\pi/2$ per rung the rail length between adjacent rungs satisfies
\[
\boxed{\;L_{n\to n+1}\ \ge\ \frac{\pi K}{2}\;}=1.4516928502\ldots\ \tagC{}
\qquad(n\ge2):
\]
every rung interval retains at least a quarter-circle of rail, even as the
vertical increments $z_{n+1}-z_n\to0$. Consequently the cumulative rail length
to the top is infinite while the height stays finite --- a one-line sharpening
of the whitepaper's logarithmic length divergence (W-Prop.\ 15.2c) into a
uniform per-rung floor. Exactly above the lens one also has the identity
$(ds/dz)^2-\bigl(K\kappa z/(1-z^2)\bigr)^2=1$, and the minorant integrates to
$-\tfrac12K\kappa\ln(1-z^2)\to\infty$.
\end{proposition}

\begin{proposition}[$\nu$-form of the rung locus]\label{prop:nulocus}\tagFD{} \hx{OZ-M6--M7}
The conjugate pair $\{W,\overline W\}$ exists automatically --- the angle
multiset of a conjugation-closed object is closed under negation (the standing
fact of the relational-charge note, its \S2.2) --- and winds oppositely:
$\theta_\pm(\rho)=\pm\kappa\rho$. Its normalized relative state
\[
\nu(\rho)\ \coloneqq\ \frac{W(\rho)}{\overline{W}(\rho)}\Big/\Bigl|\tfrac{W}{\overline W}\Bigr|
\;=\;e^{2i\theta(\rho)}
\qquad\text{(matrix form: } \mathcal R(\rho)=e^{2J\kappa\rho}\text{)}
\]
is real (scalar) exactly on the rung lattice:
\[
\operatorname{Im}\nu(\rho)=\sin\!\Bigl(\frac{\pi\rho}{\ln\gp}\Bigr)=0
\iff \rho=n\ln\gp,
\qquad
\nu(\rho_n)=(-1)^n ,
\]
and the sign $(-1)^n$ is the Trace--Form-Duality parity bit: $W_n$ is real for
even $n$ (Fibonacci/E--W rungs), imaginary for odd $n$ (Lucas/N--S rungs).
Equivalently, the rungs are exactly the heights at which $W$ is coherent with
the real axis in the sense of the relational note's $\nu$-criterion
($\alpha/\overline\alpha$ a root of unity --- here $\nu(\rho_n)=\pm1$, torsion
of order $\le2$).
\end{proposition}


\begin{proposition}[Site-density trichotomy]\label{prop:trichotomy}\tagFD{}/\tagI{}
\hx{OZ-N1--N2}\ Inverting the rung law gives the continuous floor coordinate ---
and three site densities that must be kept apart. The floor coordinate is
\[
n(z)=\frac{\rho(z)}{\ln\gp}=-\frac{\ln(1-z^2)}{2\ln\gp},
\qquad
\frac{dn}{dz}=\frac{z}{(1-z^2)\ln\gp}\ \xrightarrow[z\to1^-]{}\ \infty
\]
(at the lens, $dn/dz=2\sqrt3/\ln\gp=7.1987042602\ldots$ \tagC{} \hx{OZ-H5}).
All three densities are exact:
\begin{center}\small
\begin{tabular}{@{}lll@{}}
\toprule
measure & behaviour & source\\
\midrule
angular & fixed: $\Delta\theta=\pi/2$, four phase classes per turn & D2$+$D3\\
vertical projection & divergent: $dn/dz\to\infty$ as $z\to1$ & D1\\
rail distance (above lens) & bounded below and convergent: $\Delta s_n\downarrow \pi K/2$ & Prop.~\ref{prop:railfloor}\\
\bottomrule
\end{tabular}
\end{center}
The third entry sharpens Proposition~\ref{prop:railfloor}:
$dz/d\theta=(1-z^2)/(\kappa z)$ is strictly decreasing to $0$, so the rail
spacing $\Delta s_n=\int\sqrt{K^2+(dz/d\theta)^2}\,d\theta$ is strictly
decreasing with limit \emph{exactly} $\pi K/2$ --- the quarter-turn floor is
attained in the limit. Apparent crowding of rung sites is a projection
artifact of the D1 chart, not an increase in phase classes per revolution.
\end{proposition}

\begin{proposition}[Reversion norm and the transfer ledger]\label{prop:ledger}\tagFD{} \hx{OZ-N3--N5}
With reversion $\widetilde{aI+bJ}=aI-bJ$ (an anti-automorphism),
\[
Q\widetilde Q=\tau^2 I,
\qquad
I_n \coloneqq Q^n\widetilde{Q^n}=\tau^{2n}I=|W_n|^2 I ,
\]
and the complementary ledger
\[
\boxed{\ F_n \coloneqq 1-I_n = z_n^2\ }
\qquad\text{(D1 at the rungs, in ledger form)},
\]
with per-rung increment, via $1-\tau^2=\tau$,
\[
\Delta F_n=F_{n+1}-F_n=\tau^{2n}(1-\tau^2)=\tau^{2n+1}
=z_{n+1}^2-z_n^2 ,
\qquad
\sum_{n\ge0}\tau^{2n+1}=\frac{\tau}{1-\tau^2}=1 \ \ \text{exactly}.
\]
The odd golden powers $\tau,\tau^3,\tau^5,\ldots$ are precisely the rung
increments of $z^2$, they telescope to the unit ($1-\tau^{2N}=z_N^2$), and the
whole normalized budget is distributed over infinitely many rungs with none
consuming it at finite $n$. The per-step split is $(\tau^2,\tau)$:
transmitted co-intensity $\tau^2$, complementary increment $\tau$. All of this
is exact chart algebra; the source note's reading of $\Delta F_n$ as a
physically \emph{deposited} pulse fraction (its engine-pass model) is an
interpretation and is not imported.
\end{proposition}

\begin{remark}[Scoping: what the rails do and do not close]\label{rem:railscope}
Proposition~\ref{prop:nulocus} removes two of the exploratory note's three
assumptions: the second rail and its counter-rotation are not hypotheses but
the conjugate, forced by conjugation closure. What remains outside the forced
scope is exactly what that note's own \S10(8) concedes: (a) reading the real
crossings of $\nu$ as \emph{recoupling events that cause} rungs is an
interpretation --- with the existing rate $\kappa$ the locus merely coincides
with the ladder already forced by the gate multipliers $|m_n|=\gp^{-2n}$; and
(b) the rate itself is D3 imported, not derived --- for a general rate the
locus is $\rho_n=n\pi/(2\kappa)$, so the construction reproduces the golden
floors only because $\kappa=\pi/(2\ln\gp)$ was put in. Rung quantization (one
$Q$ per floor) therefore remains the open item named in
\S\ref{sec:quarterturn}; the energy condition, the thermodynamic reading, and
the mechanism chain of source (iv) are not imported. The second note (source
(vi)) is handled identically: its forced kernel is
Propositions~\ref{prop:trichotomy}--\ref{prop:ledger}; its engine-pass model,
pulse-budget reading, and biological program remain hypothesis --- correctly
fenced by that note's own claim ceilings and four-layer hierarchy
(geometry $\to$ kinematics $\to$ dynamics $\to$ biology), which this paper
endorses as discipline without importing as content.
\end{remark}

\begin{proposition}[Rail bounds and the gap-rate law]\label{prop:railbounds}\tagFD{}/\tagI{} \hx{QX-E1--E9}
For $n\ge2$ write $t_n=\tau^{2n}$, and on $[\theta_n,\theta_{n+1}]$ set
$x=(dz/d\theta)^2/K^2=t^2/\bigl(K^2\kappa^2(1-t)\bigr)$ with
$t=e^{-2\theta/\kappa}=1-z^2$, so that
$\Delta s_n=K\int_{\theta_n}^{\theta_{n+1}}\sqrt{1+x}\;d\theta$. With
\[
I_1(t)=\frac1{4K\kappa}\Bigl[\ln\frac{1-\tau^2t}{1-t}-\tau t\Bigr]
=\frac1{4K\kappa}\sum_{m\ge2}\frac{1-\tau^{2m}}{m}\,t^m\;>\;0,
\]
\[
I_2(t)=\frac{F(t)-F(\tau^2t)}{16K^3\kappa^3},
\quad
F(t)=\frac1{1-t}+3\ln(1-t)-3(1-t)+\frac{(1-t)^2}2,
\quad
0<I_2(t)\le\frac{1-\tau^8}{64K^7\kappa^3}\,t^4,
\]
the rail spacing of Proposition~\ref{prop:trichotomy} obeys the exact
elementary two-sided bound
\[
\boxed{\ \frac{\pi K}2+I_1(t_n)-I_2(t_n)\ \le\ \Delta s_n\ \le\
\frac{\pi K}2+I_1(t_n)\ }
\]
(lower bound via $\sqrt{1+x}\ge1+\tfrac x2-\tfrac{x^2}8$, valid since
$x\le\gapc^2/(\kappa^2K^4)<8$ on the range, interval-certified
\hx{QX-G1}). Consequences:
\[
\lim_{n\to\infty}\frac{\Delta s_n-\pi K/2}{\gapc^{\,n}}
=\frac{K}{8\kappa}=\frac{K\ln\gp}{4\pi}=0.0353900591\ldots\ \tagC{},
\qquad
\frac{\Delta s_{n+1}-\pi K/2}{\Delta s_n-\pi K/2}\longrightarrow\gapc,
\]
and since $\gapc^{\,n}=\tau^{4n}=|W_n|^4$,
\[
\Delta s_n-\frac{\pi K}2\ \sim\ \frac{K\ln\gp}{4\pi}\,|W_n|^4:
\]
the excess rail over the quarter-turn floor is priced by the square of the
co-intensity, and the divergent side of the length dichotomy
(Proposition~\ref{prop:sphlength}) is thereby quantified: the rail
diverges at exactly $\pi K/2$ per rung, approached at the gap rate.
Structural note on the exact integral: under the radical,
$(1-z^2)^2+K^2\kappa^2z^2=z^4+(K^2\kappa^2-2)z^2+1$ factors reciprocally
as $(z^2+a)(z^2+1/a)$ with $a+1/a=K^2\kappa^2-2$ (real $a>0$ since
$K\kappa>2$, interval-certified \hx{QX-G2}), placing $\Delta s_n$ in
incomplete-elliptic form --- in contrast with the residual sphere, whose
total length closes in a \emph{complete} elliptic value
(Proposition~\ref{prop:sphlength}); no elementary closed form is claimed
here, and the bounds above carry the elementary content. \hx{QX-E9}
\end{proposition}

\begin{proof}
From D1$+$D3, $dz/d\theta=(1-z^2)/(\kappa z)$ \hx{QX-E1}. The square-root
bounds are the polynomial identities $(1+x/2)^2-(1+x)=x^2/4$ and
$(1+x)-(1+x/2-x^2/8)^2=x^3(8-x)/64$ \hx{QX-E2}. Substituting
$t=e^{-2\theta/\kappa}$ ($d\theta=-\kappa\,dt/2t$) turns
$\tfrac K2\!\int x\,d\theta$ into
$\tfrac1{4K\kappa}\!\int_{\tau^2t_n}^{t_n}\!\tfrac s{1-s}\,ds$ and
$\tfrac K8\!\int x^2\,d\theta$ into
$\tfrac1{16K^3\kappa^3}\!\int_{\tau^2t_n}^{t_n}\!\tfrac{s^3}{(1-s)^2}\,ds$,
with the stated antiderivatives \hx{QX-E3--E4}; the $m=1$ series term
cancels ($1-\tau^2=\tau$), the $m=2$ term gives the constant via
$1-\tau^4=K^2$ and $1/(8\kappa)=\ln\gp/(4\pi)$ \hx{QX-E5}, and the
majorant uses $(1-s)^2\ge K^4$ on $s\le\gapc$ with
$\int s^3=t^4(1-\tau^8)/4$ \hx{QX-E6}. The ratio law is the series
quotient \hx{QX-E7--E8}. A 50-dps quadrature corroborates the bracket for
$n=2..5$ \hx{QX-G3}; display-only.
\end{proof}

\begin{theorem}[Incomplete-elliptic closed form for the rail]\label{thm:railclosed}\tagFD{}/\tagI{} \hx{EX-A1--A4, EX-B1--B8, EX-G3--G5}
Let $a>1$ be the larger root of $a+1/a=K^2\kappa^2-2$ (so
$a\cdot(1/a)=1$; $a$ real with $a>1$ strictly, interval-certified via
$K\kappa>2$, replicating \hx{QX-G2}), and set
\[
\theta(z)=\arctan\!\bigl(z\sqrt a\bigr),\qquad
m=1-\frac1{a^2},\qquad
n_c=1+\frac1a .
\]
Above the lens the rail integrand is
\[
\frac{ds}{dz}=\frac{\sqrt{(z^2+a)(z^2+1/a)}}{1-z^2}
\]
(from $\Delta s_n=\int\sqrt{K^2+(dz/d\theta)^2}\,d\theta$ and
$dz/d\theta=(1-z^2)/(\kappa z)$, with the reciprocal factorization of
Proposition~\ref{prop:railbounds}), and its antiderivative is
\[
\boxed{\ G(z)=-z\sqrt{\frac{z^2+a}{z^2+1/a}}
\;+\;\sqrt a\,E(\theta\,|\,m)
\;-\;a^{-3/2}\,F(\theta\,|\,m)
\;+\;(a+1)\,a^{-3/2}\,\Pi\!\bigl(n_c;\,\theta\,|\,m\bigr)\ },
\]
with $\Delta s_n=G(z_{n+1})-G(z_n)$ at $z_n=\sqrt{1-\tau^{2n}}$
($F,E,\Pi$ the incomplete elliptic integrals of the first, second, and
third kind \tagE{}). Certificate: $dG/dz-{}$integrand${}\equiv0$
exact-symbolically in $(a,z)$ on $a>0$, $0<z<1$ (the forced certificate),
assembled from three kind-pure lemmas through the exact partial-fraction
split of the numerator (the $(1-z^2)^{-1}$ pole, of exact weight $w+2$
with $w=K^2\kappa^2-2$, is what produces the third-kind $\Pi$ term); FTC
applicability on $[z_n,z_{n+1}]\subset(0,1)$ via
$1-n_c\sin^2\theta=(1-z^2)/(1+az^2)>0$. Numeric corroboration: the 50-dps
closed form matches quadrature of $\Delta s_n$ to $<10^{-40}$ on the
\hx{QX-G3} grid $n=2..5$ and sits inside the bracket of
Proposition~\ref{prop:railbounds}. The ``no elementary closed form is
claimed here'' scoping of that proposition is hereby completed on the
positive side: the rail side of the length dichotomy
(Proposition~\ref{prop:sphlength}) is \emph{closed in incomplete-elliptic
terms}.
\end{theorem}

\begin{corollary}[Certified non-elementarity of the rail antiderivative]\label{cor:nonelem}\tagF{}/\tagE{} \hx{EX-C1--C5, EX-G1--G2}
Forced given Liouville's theorem on integration in finite terms
(Rosenlicht 1972; algorithmic form Trager 1984) \tagE{}: the
antiderivative of the above-lens rail integrand is \emph{not} an
elementary function --- the per-rung rail length has no elementary closed
form as a function of its endpoints. All premises of the anchor are
machine-certified: the kernel curve $y^2=(z^2+a)(z^2+1/a)$ is genuinely
elliptic --- $\operatorname{disc}(z^4+wz^2+1)=16(w-2)^2(w+2)^2\neq0$ with
$w=K^2\kappa^2-2>2$, four distinct roots, genus 1 (at $a=1$ it would
collapse to $(z^2+1)^2$ and integrate elementarily); the first-kind
coefficient $c_F=-a^{-3/2}$ is nonzero \emph{unconditionally}
($c_F\,a^{3/2}=-1$ for every $a>0$); $c_E=\sqrt a$ and
$c_\Pi=(a+1)a^{-3/2}$ are nonzero; and the $\Pi$ parameters are
non-degenerate ($m\in(0,1)$ and $n_c\in(1,2)$ interval-certified;
$n_c-1$, $n_c-m$, $1-m$ all exactly nonzero, with the exact interchange
tie $n_c(2-n_c)=m$, so no special-case reduction applies), the $z=\pm1$
pole pair genuine with residues $\mp K\kappa/2\neq0$ and the exact ledger
tie $y(\pm1)^2=(1+a)(1+1/a)=K^2\kappa^2$ (the kernel over the lens poles
is the squared rail-speed constant). The closed-form dichotomy is thus
fully resolved: complete-elliptic on the residual side
(Proposition~\ref{prop:sphlength}); incomplete-elliptic, and provably
non-elementary --- certified rather than conjectured --- on the rail
side, whose elementary content is carried by the bracket and gap-rate law
of Proposition~\ref{prop:railbounds}.
\end{corollary}

\section{Epistemic ledger and findings}\label{sec:epistemics}

\textbf{Forced unconditionally, new in v1.6} \hx{QX-B, QX-C, QX-D}: the
characteristic level of the ratio object ($\chi_{\mathrm{Rat}}$, the
sign-split pair of $\gp^2$, the Adams square onto
$\operatorname{spec}(A_{\gapc})$, Proposition~\ref{prop:chirat}); the
quarter twist --- transport domain $=$ parity-homogeneous objects,
gauge-rotation invariance, anchor-exchange dichotomy (exchange iff
$m\equiv2$ mod $4$; the generator/$Z^{*}$ pair is the minimal instance and
the K-seed is anchor-fixed) (Theorem~\ref{thm:quartertwist}); the K-family
seed ladder $p_j=x^4+5F_{2j}^2x^2-5F_{2j}^2$ with uniform irreducibility
by norm obstruction, the full ledger, and the gate ties
$C_{2j}=1/|p_j(0)|$, $u_{2j}=1/\M(p_j)$
(Theorem~\ref{thm:kseedladder}).

\smallskip\noindent
\textbf{Forced given D1--D3, new in v1.6} \hx{QX-E, QX-F, QX-C12--C14}:
the exact rail-spacing bounds
$\pi K/2+I_1-I_2\le\Delta s_n\le\pi K/2+I_1$ with the gap-rate law
$\Delta s_n-\pi K/2\sim(K\ln\gp/4\pi)\,|W_n|^4$
(Proposition~\ref{prop:railbounds}); the even-rung seed mirror
$(\sin,\cos,\tan)\alpha_{2j}=(K_j,\tau^{2j},\beta_j)$
(Corollary~\ref{cor:evenmirror}); the transfer of the interpolation
classification to the formation rail, adding the sixth D3 characterization
(Remark~\ref{rem:gammatransfer}).

\smallskip\noindent
\textbf{Forced unconditionally, new in v1.5} \hx{NX-T1--T9}: the complete
relational type of the generator object (Theorem~\ref{thm:reltype}) --- two
shells $\{\pm i\tau\}$, $\{\pm i\gp\}$; contact signature
$\{\Phi_1^{4},\Phi_2^{4}\}$ complete under the full scan bound; all eight
cross-shell ratios real ($\pm\tau^2$, $\pm\gp^2$) with $\nu\equiv1$,
hence full coherence in one class; anchor $(m,n)=(2,4)$ with
$\Delta=\Z/2\Z$; the twisted-shell realization $x^4+3x^2+1=s(x^2)$,
$s=\operatorname{minpoly}(q^2)$.

\smallskip\noindent
\textbf{Forced given D1--D3, new in v1.5} \hx{NX-S1--S7, NX-H1--H3}: the
compass-cross characterization of the rung locus with
$\operatorname{Im}\nu=2xy/(x^2+y^2)$
(Proposition~\ref{prop:compasscross}); the spherical step ledger
$S_n\cdot S_{n+1}=z_nz_{n+1}$ with $d_0=\pi/2$ exactly and
$d_n\downarrow0$ (Proposition~\ref{prop:sphsteps}); the finite total
spherical length
$L_{\mathrm{res}}=\sqrt{1+\kappa^2}\,E\bigl(\kappa^2/(1+\kappa^2)\bigr)$
with the residual/formation length dichotomy and the fifth
(spherical-length) characterization of the principal branch
(Proposition~\ref{prop:sphlength}).

\smallskip\noindent
\textbf{Forced unconditionally, new in v1.4} \hx{OZ-O1--O6}: the
interpolation classification $\log Q=\{(\ln\tau)I+(\pi/2+2\pi k)J\}$ with
rung-agreement and curve-distinctness; the four equivalent characterizations
of the $k=0$ branch (D3 $=$ principal log $=$ minimal winding $=$ shortest
residual spiral per floor); the dilation mechanism $K=5^{1/4}\tau$,
$\beta=5^{1/4}\gp$ behind the bridge factor $\sqrt5$.

\smallskip\noindent
\textbf{Forced unconditionally, new in v1.3} (bridge to the relational-charge
layer, machine-verified \hx{RC-D3, RC-E1--E7}): the minimal polynomial
$x^4+3x^2+1$ of $q$ with $\M=\gp^2$ and its group-drop coherence type; the
gauge identity $\operatorname{minpoly}(q)(ix)=Z^{*}$ with both rigidity
anchors realized by the pair $(Z^{*},\,iZ^{*})$; the K-seed as catalog object
(compass angle set, $\Delta=\Z/4\Z$, $\M=\beta^2$); the relational
restatement of D2.

\smallskip\noindent
\textbf{Forced unconditionally, new in v1.1} (pure keystone algebra,
machine-verified \hx{OZ-J, OZ-K, OZ-L}): the spectral projectors and grading
$U=P_+-P_-=H/\sqrt5$; the return-kernel recoupler $A=N/\sqrt5$ with
$UA=-AU$ and mode exchange $AP_\pm=P_\mp A$; the factorization
$J=\pm UA=HN/5$ with $J^2=-I$, coordinate-free up to orientation; the
obstructions $i\notin\R[R]$ and $N\notin\R[R]$; the canonical rung generator
$Q=\pm\tau J$ with $Q^4=\gapc\,I$ and the orbit identity $Q^nv_0=(\Rea W_n,
\Ima W_n)$.

\smallskip\noindent
\textbf{Forced given D1--D3, new in the July 14 explorations} (all
machine-verified here): the packaged representation
$W(\rho)=e^{-\rho}e^{i\theta(\rho)}$ and $z=\sqrt{1-|W|^2}$; the generator
$W_n=q^n=(i\tau)^n$ with $q^4=\gapc$ and the two-/four-rung recurrences; the gap
schedule (co-intensity at two rungs, amplitude at four); the state-angle
dictionary (Thm.~\ref{thm:stateangle}); the rung ledgers
$(\sqrt\tau,\tau,\sqrt\gp)$ and $(K,\tau^2,\beta)$ with the general family
$\tan\alpha_n=\sqrt{\gp^{2n}-1}$; the bridges $\cos\alpha_2=\cos^2\alpha_1$,
$z_1=\sqrt2\sin(\alpha_2/2)$, $K=\sin\alpha_2$; the K-family
$K_j=\sqrt{1-\gapc^{\,j}}$ with uniqueness of exact $K$, the universal two-step
law, and the three-level recurrence; the lens as the $60^\circ$ state; the
path-stretch laws, the quarter-turn rail floor $L\ge\pi K/2$ with its
divergence corollary, the $\nu$-form of the rung locus with parity
$\nu(\rho_n)=(-1)^n$, the site-density trichotomy with
$\Delta s_n\downarrow\pi K/2$, and the transfer ledger $F_n=z_n^2$,
$\Delta F_n=\tau^{2n+1}$, $\sum\tau^{2n+1}=1$ (\S\ref{sec:pathstretch}); the
residual-sphere lift and the colatitude identity $\theta_K=\arccos K$; the
linearization tie $f'(K)=q^2|q|^2$. None of these is a new substrate axiom; all
are compressed representations and corollaries, exactly as the source report
states.

\smallskip\noindent
\textbf{Declared upstream} \tagD{}/\tagI{}: D1, D2, D3, the orientation
convention, and the inherited radius profile (whitepaper \S11).

\smallskip\noindent
\textbf{Not established} (carried verbatim from the source report, endorsed):
biological or DNA correspondence (the direct DNA-pitch test \emph{failed}
quantitatively and is recorded as a negative result; the helicase and
uncouple--recouple language of the v1.1 report functioned as a mechanism-finding
analogy only); physical universality; the residual sphere as fundamental
ontology; derivation of the formation radius from the complex generator;
independent validation of D1--D3 by the representation; \emph{rung
quantization} (why exactly one $Q$ per rung event --- $Q$ explains propagation,
not yet selection); \emph{continuous interpolation as selection} (Theorem~\ref{thm:interp}
classifies the interpolations: a $\Z$-family agreeing at every rung; what the
discrete powers leave undetermined is exactly the principal-branch choice,
which is D3 --- so this item merges into the rate problem); the
rail-pulse mechanism of sources (iv) and (vi) --- recoupling as \emph{cause},
the energy functional $\mathcal E$, the engine-pass/deposition reading of the
transfer ledger, and the biological comparison program --- which their own
epistemic sections correctly leave as hypothesis.
Coordinate-free uniqueness of the factorization, listed open in the source
report, is closed by Corollary~\ref{cor:coordfree}.

\begin{finding}[E-1: scoping slip in the source report's \S11]\label{find:scope}
The report's ``Forced in the whitepaper'' list includes ``quarter-turn charge
increments'' without the ``given D2'' qualifier, inconsistently with its own
``Declared'' list (which correctly names D2) and with its overall conditional
framing. Cosmetic; the report's mathematics nowhere relies on the unqualified
reading. This paper's banner (\S\ref{sec:scope}) restores the qualifier
everywhere.
\end{finding}

\begin{finding}[E-2: truncated table values in the review thread]\label{find:trunc}
The K-family table prints $K_2\approx0.989299$ and $K_3\approx0.998445$; the
exact values $0.9892996329\ldots$ and $0.9984459825\ldots$ round to
$0.989300$ and $0.998446$ at 6 dp --- the printed values are floor-truncations
\hx{OZ-H1}. Display-only; no decision weight. ($K_1$, $K_4$, and $\gapc$ are
correctly rounded.)
\end{finding}

\begin{finding}[E-3: stale verification paragraph in \S\ref{sec:scope}]\label{find:stale}
Through v1.3 and v1.4 the \S\ref{sec:scope} \textbf{Verification}
paragraph still carried the v1.0 harness description (``53 exact checks'',
ID range \hx{OZ-A1..OZ-H3}) while the title line, abstract, and manifest
(\S\ref{sec:verification}) had advanced to 91 checks across groups
\hx{OZ-A..OZ-O}. Display-only; no decision weight. Corrected in v1.5.
\end{finding}

\begin{finding}[E-4: front-matter check classification in v1.5]\label{find:hcount}
The v1.5 title line and \S\ref{sec:scope} verification paragraph counted
all $27$ checks of \texttt{nx\_softcheck.py} as exact, although its group
\hx{NX-H} comprises three 50-dps numeric display/corroboration checks
(so: $24$ exact $+$ $3$ numeric). The three numeric checks carry no
decision weight --- the harness's own header declares them
display/corroboration only, and instrumentation confirms every passing
decision was reached on the exact symbolic path. Display-only; corrected
in v1.6, which states the split and drops ``all exact'' from the title
line.
\end{finding}

\begin{finding}[E-5: stale third-harness count in the v1.6 manifest]\label{find:manifeststale}
The v1.6 \S\ref{sec:verification} third-harness sentence retained ``27
exact checks'' for \texttt{nx\_softcheck.py}, although
Finding~\ref{find:hcount} had already introduced the $24+3$ split and the
front matter (title line and the \S\ref{sec:scope} verification paragraph)
had been corrected accordingly. Display-only; no decision weight.
Corrected in v1.7.
\end{finding}

\begin{finding}[E-6: exact-count wording for the first harness]\label{find:ozsplit}
The description ``91 exact checks'' for \texttt{ozy\_softcheck.py}
(\S\ref{sec:scope} and \S\ref{sec:verification}) contradicts that
harness's own manifest table, whose \hx{OZ-H} row lists $5$ numeric
displays inside the $91$ total; under the classification standard of
Finding~\ref{find:hcount} it should read $86$ exact $+$ $5$ numeric
displays. The five numeric checks carry no decision weight. Display-only;
corrected in v1.7 (the $91/91$ and $27/27$ title-line pass ratios stand,
as pass ratios).
\end{finding}

\section{Verification manifest}\label{sec:verification}

\texttt{ozy\_softcheck.py}: 86 exact checks $+$ 5 numeric displays (the
\hx{OZ-H} row below; Finding~\ref{find:ozsplit}), exit 0; sympy 1.14.0 over
$\Q(\sqrt5)[i]$, $K$ decided at $K^2$, mpmath 50 dps display-only.

\begin{center}\footnotesize\setlength{\tabcolsep}{4pt}
\begin{tabular}{@{}lp{11.6cm}c@{}}
\toprule
group & content & checks\\
\midrule
OZ-A & generator $q=i\tau$: powers, $q^4=\gapc$, one/two/four-rung recurrences, gap schedule & 8\\
OZ-B & chart consistency $z_n=\sqrt{1-|W_n|^2}$; $z_1=\sqrt\tau$; $z_2=K=\sqrt{1-|q^2|^2}$ & 3\\
OZ-C & state-angle dictionary (5 fields, symbolic); rung tangents, TFD dress, $2\gp^{3/2}$; self-duality; quartic; K-certificate & 14\\
OZ-D & bridges: cosine contraction, not-doubling guard, half-angle, Route-5 trig, $\Theta=\alpha_2$ & 7\\
OZ-E & lens $60^\circ$ ledger; off-lattice guard & 6\\
OZ-F & K-family, odd companion, universal two-step law, monotone uniqueness, compass, linearization tie & 7\\
OZ-G & residual sphere: $|S|\equiv1$ (continuous and at rungs), exact $S_0..S_4$, two-radii separation & 5\\
OZ-J & spectral uncoupling: projectors, grading $U=H/\sqrt5$, mode action & 5\\
OZ-K & recoupler $A$, factorization $HN=5J$, $J=\pm UA$, generator $Q$, orbit identity & 9\\
OZ-L & obstructions $i,N\notin\R[R]$; coordinate-free pseudoscalar; kernel $=H^\perp$ & 4\\
OZ-M & path stretch (both branches), quarter-turn floor, divergence, rails algebra $\mathcal R=e^{2J\kappa\rho}$, $\nu$-locus \& parity & 7\\
OZ-N & floor coordinate \& density trichotomy ($\Delta s_n\downarrow\pi K/2$), reversion norm, transfer ledger, unit budget & 5\\
OZ-O & interpolation classification ($\log Q$ family, rung agreement, curve distinctness, minimality), dilation mechanism & 6\\
OZ-H & numeric displays: K-family table (truncation-aware), state angles, tangents, $\pi K/2$, lens density & 5\\
\midrule
 & \textbf{total} & \textbf{91}\\
\bottomrule
\end{tabular}
\end{center}

\smallskip\noindent
\smallskip\noindent
\textbf{Second harness} \texttt{relational\_softcheck.py}: 27 exact checks,
exit 0 --- cold re-derivation of the relational-charge layer this paper's
bridge uses: rigidity anchors (RC-A), the golden layer and sign twist (RC-B),
contact machinery with the sieved completeness bound (RC-C), pinning
corroborations including Lehmer's degree-100 ratio object and the
$\beta_4\times$Lehmer mixed scan (RC-D), the bridge derivations of
\S\ref{sec:relobject} (RC-E), and the (EG) quadratic case (RC-F). No defects
found in the checked claims of the relational note.

\smallskip\noindent
\textbf{Third harness} \texttt{nx\_softcheck.py} (this compilation,
written cold from the v1.4 statements): 24 exact checks $+$ 3 numeric
displays (Finding~\ref{find:hcount}; the stale ``27 exact'' that stood in
this sentence through v1.6 is Finding~\ref{find:manifeststale}), exit 0;
sympy
1.14.0 over $\Q(\sqrt5)[i]$ with rational-arithmetic sign decisions for
$\Q(\sqrt5)$ comparisons; mpmath 50 dps display-only. Groups: \hx{NX-R}
(replication of the load-bearing claims: minimal polynomial and roots,
Mahler $\gp^2$, contact signature, gauge identity, rung sphere identity,
$\nu$-locus parity, monotone heights --- 8); \hx{NX-T} (complete
relational type: shells, cross-shell ratios, realness and
non-unimodularity, $\nu\equiv1$, the $t_{\mathrm{rel}}$ pattern,
coherence completeness, anchor data, twisted-shell factorization --- 9);
\hx{NX-S} (compass cross, the $xy$-form of $\nu$, step orthogonality,
quarter-circle first step, monotone steps, spherical speed, elliptic length
form and bounds --- 7); \hx{NX-H} (numeric displays: $L_{\mathrm{res}}$,
$d_1$, bracket corroboration --- 3).

\smallskip\noindent
\textbf{Fourth harness} \texttt{qx\_softcheck.py} (this compilation,
written cold for v1.6): 62 exact checks $+$ 3 guards, exit 0; sympy
1.14.0 over $\Q(\sqrt5)[i]$, the K-family decided at the squared level;
the guards are two interval-certified transcendental comparisons (mpmath
\texttt{iv}) and one 50-dps quadrature corroboration, none load-bearing.
\begin{center}\footnotesize\setlength{\tabcolsep}{4pt}
\begin{tabular}{@{}lp{11.6cm}c@{}}
\toprule
group & content & checks\\
\midrule
QX-A & cold replication of the layer v1.6 builds on: generator and minimal polynomial, gauge identity, chart dictionary, even gate ladder, two-step law, transfer ledger & 15\\
QX-B & ratio object at the characteristic level: 16-ratio multiset, $\chi_{\mathrm{Rat}}$ and its factorization, torsion/nontorsion split, golden units, Adams square (Prop.~\ref{prop:chirat}) & 8\\
QX-C & K-family seed ladder: $p_j$, uniform irreducibility by norm, Eisenstein availability, full ledger, compass and contact scan, gate ties, even mirror (Thm.~\ref{thm:kseedladder}, Cor.~\ref{cor:evenmirror}) & 14\\
QX-D & quarter twist: domain, involution, golden-seed exclusion, eight witnesses, exhaustive coset simulation $m\le24$ (Thm.~\ref{thm:quartertwist}) & 10\\
QX-E & rail bounds: derivative law, root bounds, antiderivatives, series constants, majorant, ratio law, reciprocal factorization (Prop.~\ref{prop:railbounds}) & 9\\
QX-F & rail transfer: rung agreement, branch distinctness, $k=0$ minimality, monotone speed, sphere transfer (Rem.~\ref{rem:gammatransfer}) & 6\\
QX-G & guards: $x_{\max}<8$ and $K\kappa>2$ (interval-certified); 50-dps bracket of $\Delta s_n$, $n=2..5$ & 3\\
\midrule
 & \textbf{total (exact $+$ guards)} & \textbf{62$+$3}\\
\bottomrule
\end{tabular}
\end{center}

\smallskip\noindent
\textbf{v1.7 session harnesses} (eight, written cold this session, one per
dossier; every \hx{TX-}, \hx{MX-}, \hx{RX-}, \hx{UX-}, \hx{EX-},
\hx{CX-}, \hx{LX-}, and \hx{DX-} check ID cited in the v1.7 sections
above appears verbatim in the corresponding file):
\begin{center}\footnotesize\setlength{\tabcolsep}{4pt}
\begin{tabular}{@{}lp{10.7cm}c@{}}
\toprule
harness & content & exact $+$ guards\\
\midrule
\texttt{tx\_softcheck.py} & rate arithmetic: irrationality lemma, transcendence chain, near-miss closure of Guard 15.3, $\log_\gp2$, scan retirement (Lem.~\ref{lem:ratirrational}--Rem.~\ref{rem:logphi2}) & $30+6$\\
\texttt{mx\_softcheck.py} & multiplier torus: modulus ledger, systole and marking, restatement test, geodesic ledger, no-CM (Prop.~\ref{prop:modledger}--Rem.~\ref{rem:geoledger}) & $39+3$\\
\texttt{rx\_softcheck.py} & radius: parity no-go with $\Z/2$ sharpening, area-transfer law, the three equipotential clocks (Thm.~\ref{thm:paritynogo}--Rem.~\ref{rem:rad4}) & $25+4$\\
\texttt{ux\_softcheck.py} & universal seed ladder $n=1..10$ $+$ symbolic-in-$w$, tangent-dress Mahler ladder, Theorem U window and monotonicity, F1 falsifier (Thm.~\ref{thm:seedladder}--Rem.~\ref{rem:f1falsifier}) & $35+1$\\
\texttt{ex\_softcheck.py} & rail closed form: partial-fraction assembly, exact-derivative certificate, ellipticity, coefficient nonvanishing (Thm.~\ref{thm:railclosed}, Cor.~\ref{cor:nonelem}) & $17+5$\\
\texttt{cx\_softcheck.py} & $\chi$-selection: pentagon-by-image, unimodular monomials, conditional chain, de-axiomatization (Props.~\ref{prop:pentimage}--\ref{prop:unimon}, Thm.~\ref{thm:chichain}) & $34+2$\\
\texttt{lx\_softcheck.py} & lens: scoping lemma, rational-height splits, bounded-negative scan ($7$ scan lines, bookkeeping, counted separately) (Prop.~\ref{prop:lensscope}, Rem.~\ref{rem:lensscan}) & $26+2$\\
\texttt{dx\_softcheck.py} & OP-RATE core: gate-ladder replication, $\pm Q$ uniqueness, sign ladder and axiom H, ledger counter-models, $k$-invariance and R6 no-go (Prop.~\ref{prop:kinv}, Rem.~\ref{rem:oprate}) & $28+6$\\
\midrule
 & \textbf{total} & $\mathbf{234+29}$\\
\bottomrule
\end{tabular}
\end{center}
Environment pin: Python 3.12.10, sympy 1.14.0, mpmath 1.3.0 (the
\texttt{py} launcher); every exact decision in $\Q/\Q(\sqrt5)$ extended by
$i$, with $\Q(\sqrt5)$ signs by rational arithmetic and $K$ decided at the
squared level; the guards are interval-certified transcendental
comparisons (mpmath \texttt{iv}, 60 dps), 50-dps quadrature
corroborations, or falsifier rejections (\texttt{dx}), none load-bearing;
all eight exit 0; each dossier passed two-lane adversarial verification
(independent harness audit plus mathematical audit) before inclusion here.

\smallskip\noindent
\textbf{v1.8 session harnesses} (five, written cold this session, one per
dossier; every \hx{CL-}, \hx{RD-}, \hx{CH-}, \hx{RN-}, and \hx{PM-} check ID
cited in the v1.8 sections above appears in the corresponding file --- the
\hx{CL-2b(i)--(vii)} row suffixes are assembled as the seven-row loop runs):
\begin{center}\footnotesize\setlength{\tabcolsep}{4pt}
\begin{tabular}{@{}lp{10.7cm}c@{}}
\toprule
harness & content & exact $+$ guards\\
\midrule
\texttt{rd\_softcheck.py} & OP-RADIUS: apex-ramification valuation certificate ($e=2$ for $r^2$, $e=4$ for $r$) and the constructive-bridge obstruction; atom ledger with the cap-onset negative (Thm.~\ref{thm:apexram}--Rem.~\ref{rem:atomledger}) & $22+2$\\
\texttt{cl\_softcheck.py} & OP-RATE R6 classification no-go: extended $k$-invariance, the selection-functional no-go, the seven-way table, the Cencov-mirror closure (Lem.~\ref{lem:kinvext}--Rem.~\ref{rem:sevenway}) & $20+5$\\
\texttt{ch\_softcheck.py} & $\chi$ doctrine g1/g3 resolved negatively: on-circle image of $S=\mu_2$, the charge character, the D2$'$ dichotomy (Prop.~\ref{prop:oncircleimage}, Rem.~\ref{rem:chidoctrine}) & $26+2$\\
\texttt{rn\_softcheck.py} & relational note: P4 recovered/reduced to P3, the bracketed odd floor, the bounded P1$'$ transport-orbit scan (Rem.~\ref{rem:p4recover}--Rem.~\ref{rem:p1scan}) & $20+4$\\
\texttt{pm\_softcheck.py} & period-relation map: $\tau_0$ algebraic and no $\kappa$-vs-$\tau_0$ relation, $L_{\mathrm{res}}$ as an elliptic period, the external-open ledger (Prop.~\ref{prop:tau0alg}, Rem.~\ref{rem:periodmap}) & $27+7$\\
\midrule
 & \textbf{total} & $\mathbf{115+20}$\\
\bottomrule
\end{tabular}
\end{center}
Across the thirteen session harnesses (\texttt{tx}, \texttt{mx}, \texttt{rx},
\texttt{ux}, \texttt{ex}, \texttt{cx}, \texttt{lx}, \texttt{dx}, \texttt{cl},
\texttt{rd}, \texttt{ch}, \texttt{rn}, \texttt{pm}) the running total is
$\mathbf{349}$ exact checks $+$ $\mathbf{49}$ certified guards, all exit 0.
Environment pin: Python 3.12.10, sympy 1.14.0, mpmath 1.3.0 (the \texttt{py}
launcher); every exact decision in $\Q/\Q(\sqrt5)$ extended by $i$ (and by
$\sqrt{13}$ / $5^{1/4}$ where a Salem trace-down or an off-circle modulus
enters), with $\Q(\sqrt5)$ signs by rational arithmetic; the guards are
interval-certified transcendental comparisons (mpmath \texttt{iv}, 60 dps),
$50$--$140$-dps quadrature/\texttt{pslq}/decimal corroborations, or fail-first
falsifier rejections, none load-bearing; all five exit 0; each dossier passed
two-lane adversarial verification, with the companion papers read from the
Echo-S-Research archive read-only and their machinery re-encoded cold.

\smallskip\noindent
\textbf{v1.9 session harnesses} (four, written cold this session, one per
dossier; every \hx{RA-}, \hx{QT-}, \hx{OF-}, and \hx{PF-} check ID cited in the
v1.9 sections above appears in the corresponding file):
\begin{center}\footnotesize\setlength{\tabcolsep}{4pt}
\begin{tabular}{@{}lp{10.7cm}c@{}}
\toprule
harness & content & exact $+$ guards\\
\midrule
\texttt{ra\_softcheck.py} & OP-RADIUS D4-reduction ($D4\Rightarrow r$ unique $+$ minimality) and the differential strengthening (honest negative $+$ conditional apex-monodromy-order sharpening) (Thm.~\ref{thm:d4reduction}--Rem.~\ref{rem:ratie}) & $23+4$\\
\texttt{qt\_softcheck.py} & D2 relocation: the $K$-seed rotation axis forces $\chi=\pm\pi/2$ off-circle (disjoint from g1); net verdict $W[\text{open}]2$ advanced not closed (Prop.~\ref{prop:d2reloc}, Rem.~\ref{rem:d2verdict}) & $19+3$\\
\texttt{of\_softcheck.py} & odd relational floor: computed $=2$ up to a logged box, reduction to the single lemma (ODD-2), P4 transfer (Prop.~\ref{prop:oddfloorred}, Rem.~\ref{rem:p4transfer}) & $20+6$\\
\texttt{pf\_softcheck.py} & period frontier: the forced-identity ledger (with the new elliptic geometric identity), the Baker/Schanuel/external-open tool-gap map, and the validated-detector PSLQ bounded-negative (Rem.~\ref{rem:pfledger}, Rem.~\ref{rem:pffrontier}) & $13+13\,\mathrm{PSLQ}+5$\\
\midrule
 & \textbf{total (exact/computed $+$ guards $+$ PSLQ)} & $\mathbf{75+18+13}$\\
\bottomrule
\end{tabular}
\end{center}
Across the seventeen session harnesses (\texttt{tx}, \texttt{mx}, \texttt{rx},
\texttt{ux}, \texttt{ex}, \texttt{cx}, \texttt{lx}, \texttt{dx}, \texttt{cl},
\texttt{rd}, \texttt{ch}, \texttt{rn}, \texttt{pm}, \texttt{ra}, \texttt{qt},
\texttt{of}, \texttt{pf}) the running total is $\mathbf{424}$ exact checks $+$
$\mathbf{67}$ certified guards $+$ $\mathbf{13}$ PSLQ corroborations, all exit 0.
Environment pin: Python 3.12.10, sympy 1.14.0, mpmath 1.3.0 (the \texttt{py}
launcher); every exact decision in $\Q/\Q(\sqrt5)$ extended by radicals and $i$
(and by $\sqrt{13}$ where the Salem trace-down enters), with $\Q(\sqrt5)$ signs
by rational arithmetic and $K$ decided at the squared level; the $18$ guards are
interval-certified (mpmath \texttt{iv}, 60 dps), while the $13$ PSLQ
integer-relation searches (230 dps, coefficient-height $\le10^{12}$) are counted
as \tagC{} corroboration --- never a transcendence or algebraic-independence
proof; all four exit 0; each dossier passed two-lane adversarial verification,
with the companion papers read from the Echo-S-Research archive read-only and
their machinery re-encoded cold.

\smallskip\noindent
Every boxed identity of the source reports and every displayed identity of the
review thread appears above with a passing check ID; the six findings
(\ref{find:scope}, \ref{find:trunc}, \ref{find:stale}, \ref{find:hcount},
\ref{find:manifeststale}, \ref{find:ozsplit})
are the complete defect list, unchanged in v1.8: the five v1.8 verifiers
flagged only citation-hygiene in the derivation reports (companion-paper
theorem numbers), corrected on the way in, and found no new display defect. The
whitepaper's own harnesses (87$+$1 shipped; 89 independent) certify the upstream
structure this paper conditions on.

\end{document}
